% Definicja składni języka Kotlin dla pakietu listings (brak wsparcia wbudowanego)
\lstdefinelanguage{Kotlin}{
 keywords={package, import, class, object, interface, data, fun, val, var,
 private, public, internal, const, companion, return, if, else, when, for,
 while, do, break, continue, in, is, as, null, true, false, override,
 suspend, this, super, try, catch, finally, throw, by, out, sealed, enum},
 keywordstyle=\bfseries\color{EEUltramarine},
 ndkeywords={Int, Float, Long, Boolean, String, List, Map, Bitmap, Flow,
 Mutex, ArrayDeque, Dispatchers, Unit},
 ndkeywordstyle=\color{EEBlueDark},
 comment=[l]{//},
 morecomment=[s]{/*}{*/},
 commentstyle=\color{EEBlueDark},
 stringstyle=\color{EETangerine},
 morestring=[b]",
 sensitive=true
}
% Poprawna obsługa polskich znaków diakrytycznych wewnątrz listingów
\lstset{
 extendedchars=true,
 literate=%
 {ą}{{\k{a}}}1 {ć}{{\'c}}1 {ę}{{\k{e}}}1 {ł}{{\l{}}}1 {ń}{{\'n}}1
 {ó}{{\'o}}1 {ś}{{\'s}}1 {ź}{{\'z}}1 {ż}{{\.z}}1
 {Ą}{{\k{A}}}1 {Ć}{{\'C}}1 {Ę}{{\k{E}}}1 {Ł}{{\L{}}}1 {Ń}{{\'N}}1
 {Ó}{{\'O}}1 {Ś}{{\'S}}1 {Ź}{{\'Z}}1 {Ż}{{\.Z}}1
}

\chapter{Wstęp}\label{chap:wstep}
Przedmiotem pracy jest automatyczna identyfikacja gatunków grzybów na urządzeniach mobilnych pracujących autonomicznie, bez połączenia z serwerem. Grzybobranie – w Polsce i szerzej w Europie Środkowej – pozostaje powszechną praktyką rekreacyjną i kulinarną. Jednocześnie niesie realne ryzyko zdrowotne, gdy owocniki zostaną źle rozpoznane. Większość ciężkich zatruć nie wynika ze świadomego spożycia gatunków trujących, lecz z powierzchownej oceny cech morfologicznych przez osoby bez przygotowania mykologicznego. Ta rozbieżność między łatwą dostępnością grzybów w lasach a ograniczoną wiedzą przeciętnego zbieracza uzasadnia podjęcie tego tematu.

Za przeważającą część śmiertelnych zatruć grzybami na świecie, szacowaną na ponad 90\%, odpowiadają gatunki zawierające amatoksyny, przede wszystkim muchomor sromotnikowy (\textit{Amanita phalloides}) \cite{garcia2015amanita, mengs2012legalon}. Toksyny te cechują się wysoką hepatotoksycznością, a dawka śmiertelna $\alpha$-amanityny dla człowieka wynosi około 0,1~mg/kg masy ciała, co oznacza, że pojedynczy owocnik może stanowić zagrożenie dla życia osoby dorosłej \cite{garcia2015amanita}. Dodatkowym czynnikiem pogarszającym rokowanie jest wielogodzinny okres utajenia, w którym objawy kliniczne nie występują, przez co pomoc medyczna bywa udzielana ze znacznym opóźnieniem. Według danych europejskich i północnoamerykańskich ośrodków toksykologicznych liczbę zgonów spowodowanych zatruciami grzybami szacuje się na kilkadziesiąt do około dwustu przypadków rocznie \cite{mengs2012legalon}, przy czym rzeczywista skala zjawiska pozostaje trudna do precyzyjnego ustalenia z powodu niepełnej sprawozdawczości. W warunkach polskich zatrucia grzybami wykazują wyraźną sezonowość i nasilają się w okresie jesiennym, a ich bezpośrednią przyczyną jest najczęściej pomylenie gatunku jadalnego z jego trującym odpowiednikiem.

Z tych powodów sens ma narzędzie wspierające wstępną identyfikację grzybów bezpośrednio w terenie, na typowym smartfonie. System ma wspierać użytkownika przy ocenie owocników ułożonych na płaskiej powierzchni. Zwłaszcza w kwestii bezpieczeństwa mykologicznego i wyraźnego sygnalizowania gatunków potencjalnie trujących.

\section{Motywacja podjęcia tematu}
Motywacja do opracowania systemu wynika z trzech przesłanek: skali problemu zatruć, trudności rozróżniania gatunków bliźniaczych oraz ograniczeń dotychczasowych rozwiązań mobilnych. Kolejne podrozdziały omawiają je po kolei i uzasadniają wybór zastosowanych technologii.

\subsection{Problem gatunków bliźniaczych}
Kluczową trudnością w automatyzacji rozpoznawania grzybów jest zjawisko gatunków bliźniaczych (ang. \textit{look-alikes}), czyli par lub grup gatunków o zbliżonych cechach makroskopowych, spośród których część jest jadalna, a część silnie toksyczna. Wizualne podobieństwo owocników sprawia, że rozróżnienie oparte wyłącznie na kształcie kapelusza, barwie czy pokroju bywa zawodne nawet dla doświadczonych zbieraczy, a subtelne cechy diagnostyczne, takie jak budowa blaszek, obecność pierścienia czy pochwy u podstawy trzonu, są często pomijane.

Do najbardziej reprezentatywnych przykładów należy pieprznik jadalny (\textit{Cantharellus cibarius}), popularnie zwany kurką, mylony z lisówką pomarańczową (\textit{Hygrophoropsis aurantiaca}) oraz z toksycznym gatunkiem \textit{Omphalotus olearius}, którego spożycie wywołuje ostre dolegliwości żołądkowo-jelitowe. Błąd ten bywa śmiertelny. Młode owocniki pieczarki polnej (\textit{Agaricus campestris}) do złudzenia przypominają wczesne stadia muchomora sromotnikowego (\textit{Amanita phalloides}) oraz jadowitego (\textit{Amanita virosa}). We wczesnym stadium rozwoju oba muchomory przypominają jadalne gatunki o białych kapeluszach. Zdolność do niezawodnego rozróżniania takich par jest zatem nadrzędnym kryterium projektowym systemu, w którym błąd klasyfikacji może mieć konsekwencje zagrażające życiu.

Skala problemu wzrasta dodatkowo wraz z liczbą rozpoznawanych taksonów. Zadanie klasyfikacji drobnoziarnistej (ang. \textit{fine-grained classification}) na zbiorze obejmującym $147$ gatunków grzybów wymaga od modelu różnicowania klas na podstawie subtelnych, lokalnych cech morfologicznych – układu blaszek, faktury powierzchni kapelusza czy odcienia miazgi po przełamaniu – przy jednoczesnym zachowaniu odporności na wewnątrzklasową zmienność wynikającą z wieku owocnika, wilgotności podłoża oraz oświetlenia scenerii leśnej. W przeciwieństwie do zadań klasyfikacji gruboziarnistej, w których klasy różnią się istotnie w przestrzeni cech wysokopoziomowych, granice decyzyjne pomiędzy gatunkami bliźniaczymi bywają wąskie i nieliniowe, a błąd predykcji koncentruje się właśnie w ich sąsiedztwie. Zwiększenie liczby klas z rzędu kilkudziesięciu do $147$ pogłębia to zjawisko, ponieważ przestrzeń cech musi jednocześnie separować znacznie więcej par potencjalnie mylonych taksonów, co wymaga architektury o odpowiednio dużej pojemności reprezentacyjnej. Zagadnienie to rozwinięto w rozdziale \ref{chap:przeglad_rozwiazan}.

\subsection{Uzasadnienie wyboru detekcji wieloobiektowej}
Naturalnym podejściem do rozpoznawania gatunku grzyba na podstawie obrazu wydaje się klasyfikacja całego obrazu (ang. \textit{image classification}), przypisująca pojedynczą etykietę klasy do całej sceny. Podejście to zakłada jednak obecność dokładnie jednego dominującego obiektu w kadrze i traci użyteczność w typowym scenariuszu domowym, w którym zebrane owocniki są rozkładane obok siebie na stole lub w koszu. W takich warunkach kilka egzemplarzy, potencjalnie należących do różnych gatunków, występuje jednocześnie w jednym kadrze, a ich fragmenty wzajemnie się przysłaniają i nakładają na siebie, co uniemożliwia poprawne przypisanie jednej etykiety do całego obrazu.

Rozwiązaniem tego problemu jest detekcja wieloobiektowa (ang. \textit{multi-object detection}), która jednocześnie lokalizuje każdy owocnik za pomocą ramki otaczającej oraz przypisuje mu odrębną klasę gatunkową. Rozwiązanie takie pozwala na indywidualną ocenę każdego egzemplarza w scenie i – co ważne w kontekście bezpieczeństwa – na selektywne oznaczenie obiektów sklasyfikowanych jako trujące. Ponieważ narzędzie ma działać na strumieniu obrazu z kamery w czasie rzeczywistym, wybrano jednoprzebiegowy detektor z rodziny YOLO (ang. \textit{You Only Look Once}), której pierwotna koncepcja realizuje detekcję jako pojedyncze przejście sieci neuronowej i zapewnia wysoką przepustowość przetwarzania \cite{redmon2016yolo}. Jako model bazowy przyjęto architekturę YOLO11 Medium (\texttt{YOLO11m}), stanowiącą kompromis pomiędzy dokładnością a kosztem obliczeniowym, adekwatny do ograniczeń urządzeń mobilnych przy zadaniu rozróżniania wielu gatunków \cite{ultralytics2024yolo11}.

\subsection{Konieczność działania w trybie offline}
Całkowita niezależność od sieci to wymaganie twarde, nie preferencja. Grzybobranie odbywa się z reguły w głębi lasu, gdzie zasięg GSM i LTE bywa ograniczony albo zerowy. W takich warunkach rozwiązania oparte na architekturze klient–serwer, przesyłające obraz do zdalnego interfejsu chmurowego (ang. \textit{cloud-based API}), stają się bezużyteczne dokładnie w tym miejscu i momencie, gdy są najbardziej potrzebne.

Dlatego przyjęto, że cały proces inferencji musi być realizowany lokalnie, na urządzeniu użytkownika (ang. \textit{on-device inference}). Wymóg ten narzuca dodatkowe ograniczenia dotyczące rozmiaru modelu oraz zapotrzebowania na pamięć i moc obliczeniową. Podporządkowano im dobór architektury sieci oraz proces jej optymalizacji. Aby umożliwić efektywne uruchomienie modelu na sprzęcie mobilnym, przeprowadzono jego konwersję do formatu TensorFlow Lite z zastosowaniem kwantyzacji wag \cite{tflite}; szczegóły podano w rozdz.~\ref{chap:implementacja}.

\section{Cel i zakres pracy}
Celem pracy jest zaprojektowanie, implementacja i ewaluacja autonomicznego systemu detekcji oraz klasyfikacji gatunków grzybów na platformie Android. W czasie rzeczywistym i offline. Realizacja obejmuje model uczenia maszynowego oraz aplikację mobilną, która ma wspomagać identyfikację owocników w terenie, a nie zastępować ocenę eksperta.

Osiągnięcie celu głównego wymagało zrealizowania następujących zadań szczegółowych:
\begin{enumerate}[label=\textbf{Z\arabic*.}]
 \item \textbf{Opracowanie i wytrenowanie modelu detekcji} wieloobiektowej YOLO11m obejmującego $147$ gatunków grzybów, zwłaszcza par gatunków bliźniaczych, przy wykorzystaniu uczenia transferowego na bazie wag wstępnie wytrenowanych na zbiorze COCO \cite{lin2014coco}.
 \item \textbf{Zaprojektowanie półautomatycznego potoku anotacji} zbioru danych z wykorzystaniem modelu wizyjno-językowego GroundingDINO \cite{liu2023grounding}, ograniczającego pracochłonność ręcznego etykietowania ramek otaczających.
 \item \textbf{Przeprowadzenie optymalizacji i konwersji modelu} do postaci umożliwiającej lokalną inferencję na urządzeniu mobilnym, obejmującej dwuetapową ścieżkę PyTorch $\to$ TensorFlow SavedModel $\to$ TensorFlow Lite oraz dobór precyzji numerycznej FP16.
 \item \textbf{Zintegrowanie modelu z natywną aplikacją} systemu Android, napisaną w języku Kotlin z wykorzystaniem biblioteki CameraX oraz deklaratywnego interfejsu Jetpack Compose, wraz z implementacją asymetrycznej, programowej logiki bezpieczeństwa.
 \item \textbf{Przeprowadzenie ilościowej oceny skuteczności} rozwiązania na izolowanym zbiorze testowym, z naciskiem na wykrywalność gatunków trujących oraz macierz ryzyka bezpieczeństwa mykologicznego.
\end{enumerate}

Zakres pracy obejmuje stworzenie kompletnego potoku przetwarzania danych (ang. \textit{machine learning pipeline}), automatyzację procesu anotacji zbioru danych z wykorzystaniem modelu GroundingDINO \cite{liu2023grounding}, trening głębokiej sieci neuronowej YOLO11 Medium (\texttt{YOLO11m}), jej konwersję do formatów ONNX oraz TensorFlow Lite, a także integrację zoptymalizowanego modelu z aplikacją mobilną wykorzystującą interfejs CameraX. Poza zakresem pracy pozostają zagadnienia mykologiczne wykraczające poza identyfikację wizualną, takie jak analiza chemiczna czy mikroskopowa cech diagnostycznych owocników.

\section{Hipotezy badawcze}
Realizację celu głównego oparto na dwóch hipotezach badawczych, weryfikowanych ilościowo w rozdziale \ref{chap:badania}:
\begin{description}[topsep=2pt,itemsep=2pt]
 \item[$H_1$] Jednoprzebiegowa detekcja wieloobiektowa YOLO11m w precyzji FP16 umożliwia autonomiczną identyfikację $147$ gatunków grzybów w czasie rzeczywistym na urządzeniu mobilnym, bez dostępu do sieci.
 \item[$H_2$] Programowa, asymetryczna logika progów decyzyjnych ($0{,}18$ dla klas trujących, $0{,}60$ dla klas jadalnych) ogranicza wskaźnik krytycznego ryzyka (CFPR) na izolowanym zbiorze testowym względem pojedynczego progu globalnego.
\end{description}
Hipoteza $H_1$ dotyczy warstwy algorytmiczno-sprzętowej: weryfikuje się ją pomiarem $mAP$ wag PyTorch na izolowanym teście, porównaniem $mAP$ na próbie $2\,500$ kadrów między \texttt{.pt} a TFLite FP16 (Interpreter i dekoder zgodny z \texttt{MushroomDetector.kt}, nie LiteRT \texttt{val}; podrozdz.~\ref{sec:tflite_parity}) oraz pomiarem \texttt{detect()} na S23 Ultra (rozdz. \ref{chap:badania}). Hipoteza $H_2$ dotyczy pary progów $0{,}18$/$0{,}60$ względem jednego progu globalnego $0{,}30$ na \emph{tym samym} zbiorze testowym. Oba CFPR policzono skryptem \texttt{test\_h2\_baseline.py}: $1{,}50\%$ wobec $3{,}11\%$. Bufor konsensusu klatek ($4$ z $5$) nie wchodzi w treść $H_2$; opisano go w rozdz.~\ref{chap:implementacja} jako mechanizm stabilizacji czasowej nakładki.

\section{Układ pracy}
Praca składa się z sześciu rozdziałów, zgodnych z kolejnością prac nad projektem. Rozdział bieżący przedstawia motywację, cel i zakres. W rozdziale drugim omówiono stan wiedzy: aplikacje mobilne do identyfikacji grzybów, metody detekcji i klasyfikacji obiektów oraz optymalizację modeli pod urządzenia o ograniczonych zasobach. Trzeci rozdział dotyczy projektu systemu. Wymagań funkcjonalnych i niefunkcjonalnych oraz doboru architektury.

Czwarty opisuje implementację: przygotowanie danych, anotację, uczenie, konwersję modelu i aplikację mobilną. Piąty zawiera badania skuteczności i wydajności, w tym analizę gatunków bliźniaczych oraz jakościową ocenę detekcji w warunkach terenowych na danych zewnętrznych. Rozdział szósty podsumowuje wyniki i wskazuje kierunki dalszych prac.


\chapter{Przegląd istniejących rozwiązań}\label{chap:przeglad_rozwiazan}
Autonomiczny system identyfikacji owocników na urządzeniu wbudowanym wymaga uprzedniego przeglądu trzech obszarów: dostępnych aplikacji, metod detekcji i klasyfikacji w obrazie oraz technik dostosowania modeli do ograniczonego sprzętu. Takie tło pozwala ocenić istniejące podejścia, wskazać ich ograniczenia i uzasadnić decyzje projektowe omówione dalej.

\section{Analiza istniejących aplikacji mobilnych}
Dwie aplikacje wyznaczają górną granicę możliwości komercyjnych rozwiązań. iNaturalist \cite{inaturalist} gromadzi miliony obserwacji przyrodniczych i klasyfikuje je modelami trenowanymi na tej samej skali; Picture Mushroom \cite{picturemushroom} zawęża zakres do grzybów i osiąga imponującą trafność dla zdjęć studyjnych. Obie buduje się jednak na tym samym założeniu: smartfon to tylko aparat. Obraz wędruje na serwer, tam model inferuje i wraca wynik.

Ten wzorzec klient--serwer sprawdza się w parku z LTE. W lesie LTE zwykle nie ma. Polskie kompleksy leśne (Puszcza Białowieska, Bory Tucholskie, lasy Sudetów) to typowe tereny grzybobrania i typowe martwe strefy GSM. W głębi drzewostanu sygnał jest słaby albo go nie ma wcale. Gdy grzybiarz wyciąga telefon nad znaleziskiem, chmura mu nie odpowie.

Trzy konsekwencje tego wzorca są bezpośrednio szkodliwe dla scenariusza terenowego:
\begin{enumerate}
 \item Aplikacja traci funkcjonalność dokładnie tam, gdzie jest najbardziej potrzebna: w lesie bez zasięgu GSM/LTE.
 \item Dwukierunkowa transmisja zdjęcia wysokiej rozdzielczości wyklucza analizę strumienia wideo w czasie rzeczywistym.
 \item Żadna z aplikacji nie obsługuje jednoczesnej detekcji wielu owocników w kadrze -- jedno zdjęcie, jeden wynik.
\end{enumerate}

Porównanie właściwości aplikacji komercyjnych i systemu opracowanego w tej pracy zestawiono w tab.~\ref{tab:aplikacje_porownanie}. Porównanie prowadzono według pięciu kryteriów istotnych dla scenariusza terenowego: modelu wykonania inferencji, dostępności bez łączności sieciowej, zdolności do analizy strumienia wideo, obsługi wielu obiektów w jednym kadrze oraz obecności dedykowanej warstwy sygnalizacji ryzyka toksykologicznego.

\begin{table}[htbp]
 \centering
 \caption{Porównanie funkcjonalne popularnych aplikacji identyfikacyjnych z systemem opracowanym w pracy. Źródło: opracowanie własne na podstawie \cite{inaturalist, picturemushroom}.}
 \label{tab:aplikacje_porownanie}
 \begin{tabular}{|p{5.4cm}|c|c|c|}
 \hline
 \textbf{Kryterium} & \textbf{iNaturalist} & \textbf{Picture Mushroom} & \textbf{System własny} \\
 \hline \hline
 Model wykonania inferencji & klient--serwer & klient--serwer & na urządzeniu \\
 \hline
 Praca bez dostępu do sieci & nie & nie & tak \\
 \hline
 Analiza strumienia wideo & nie & nie & tak ($4$\,FPS) \\
 \hline
 Wiele obiektów w kadrze & nie & nie & tak \\
 \hline
 Warstwa ostrzegania toksykologicznego & brak & informacyjna & asymetryczna \\
 \hline
 \end{tabular}
\end{table}

Chmura zawodzi dokładnie w środowisku, dla którego narzędzie jest projektowane. Dlatego w tej pracy przyjęto odwrotny paradygmat: przetwarzanie krawędziowe (ang. \textit{edge AI}), gdzie cały model działa lokalnie na urządzeniu. To oznacza konieczność zmieszczenia wystarczająco dobrego detektora w budżecie pamięci i mocy obliczeniowej smartfona. Techniki tego dostosowania opisano w kolejnych podrozdziałach.

\section{Przegląd metod detekcji i klasyfikacji obiektów}
Głębokie sieci neuronowe zastąpiły ręcznie projektowane deskryptory obrazu – ich zdolność do hierarchicznej ekstrakcji cech pozwala odróżnić blaszki muchomora od blaszek pieczarki bez jawnego kodowania zasad mykologii. Poniżej opisano różnicę między klasyfikacją a detekcją, porównano rodziny detektorów i uzasadniono dobór architektury zastosowanej w pracy.

\subsection{Splotowe sieci neuronowe i transformatory}
Splotowe sieci neuronowe (CNN) działają na operacji splotu obrazu z wyuczalnymi filtrami \cite{lecun1998gradient}. Pierwsze warstwy wychwytują krawędzie i tekstury, głębsze składają je w coraz bardziej abstrakcyjne cechy, aż do poziomu rozróżnialnego gatunkowo układu blaszek czy faktury kapelusza. Współdzielenie wag nadaje niezmienniczość translacyjną, a pooling redukuje wymiarowość. Klasyczną ilustrację tej hierarchii pokazuje rys.~\ref{fig:lenet5_web} z architekturą LeNet-5 \cite{lecun1998gradient}. Połączenia rezydualne (ResNet) usunęły barierę zanikającego gradientu i umożliwiły efektywne uczenie sieci o dziesiątkach warstw \cite{he2016resnet}.

\begin{figure}[htbp]
 \centering
 \includegraphics[width=0.78\textwidth]{lenet5_official}
 \caption{Architektura klasycznej splotowej sieci neuronowej LeNet-5. Źródło: \cite{lecun1998gradient}.}
 \label{fig:lenet5_web}
\end{figure}

Transformatory wizyjne (ViT) \cite{dosovitskiy2021vit} zaadaptowały mechanizm samouwagi (ang. \textit{self-attention}) z NLP \cite{vaswani2017attention} do analizy obrazu. Samouwaga modeluje globalne zależności między odległymi fragmentami obrazu, czego sieć splotowa z ograniczonym polem receptywnym nie zrobi tanio. W detekcji obiektów mechanizmy uwagi działają w rodzinie DETR, na której opiera się model GroundingDINO \cite{liu2023grounding} – kluczowy dla auto-anotacji w tej pracy.

\subsection{Detekcja otwartego słownika: GroundingDINO}
YOLO zna tylko klasy, które dostał w treningu. Do budowy zbioru to fatalna cecha: trzeba najpierw narysować setki tysięcy ramek, zanim w ogóle ruszy uczenie gatunków. Ręczne BBOX-y dla pełnej bazy byłyby nierealne w skali pracy inżynierskiej.

GroundingDINO omija ten impas \cite{liu2023grounding}. To detektor otwartego słownika: zamiast indeksu klasy dostaje tekst. Enkoder Swin patrzy na obraz, BERT czyta frazę, a \textit{cross-modality attention} liczy zgodność regionu z tekstem (rys.~\ref{fig:grounding_dino_web}). W potoku anotacji zapytanie brzmiało po prostu \texttt{"mushroom"}. Model wskazywał owocnik bez wiedzy, czy to borowik, czy muchomor – gatunek brał się z folderu iNaturalist, nie z głowicy. Zero-shotowa lokalizacja oddzieliła „gdzie jest grzyb” od „jaki to gatunek”. Bez tego zbioru $147$ klas nie byłoby.

\begin{figure}[htbp]
 \centering
 \includegraphics[width=0.88\textwidth]{grounding_dino_arch_official}
 \caption{Architektura detektora otwartego słownika Grounding DINO. Źródło: \cite{liu2023grounding}.}
 \label{fig:grounding_dino_web}
\end{figure}

\subsection{Klasyfikacja obrazu a detekcja wieloobiektowa}
Klasyfikator widzi cały kadr i wydaje jeden werdykt. Formalnie odwzorowanie takie przypisuje obrazowi jedną etykietę spośród $K$ klas:
\begin{equation}
f_{\mathrm{cls}}:\; \mathbb{R}^{H\times W\times C} \rightarrow \{1, 2, \dots, K\},
\end{equation}
gdzie
\begin{description}[topsep=2pt,itemsep=-3pt]
 \item[$H, W, C$] wysokość, szerokość oraz liczba kanałów obrazu wejściowego,
 \item[$K$] liczba rozpoznawanych klas.
\end{description}
Softmax zamienia wynikowy wektor logitów w rozkład prawdopodobieństwa, a odpowiedzią jest klasa o najwyższym wyniku. Jeden kadr, jedna odpowiedź. Nie ma mowy o położeniu. Nie ma mowy o kilku gatunkach jednocześnie – a dokładnie taki jest las: kosz, w którym leżą kurka, podgrzybek i niejadalna gołąbka naraz.

Detekcja wieloobiektowa rozrywa to ograniczenie. Wynikiem nie jest etykieta, lecz zbiór krotek:
\begin{equation}
f_{\mathrm{det}}:\; \mathbb{R}^{H\times W\times C} \rightarrow \big\{(\mathbf{b}_i, c_i, s_i)\big\}_{i=1}^{N},
\end{equation}
gdzie
\begin{description}[topsep=2pt,itemsep=-3pt]
 \item[$\mathbf{b}_i$] ramka otaczająca (BBOX) $i$-tego obiektu, opisana wektorem współrzędnych $(x_i, y_i, w_i, h_i)$,
 \item[$c_i$] etykieta klasy przypisana obiektowi,
 \item[$s_i$] wartość pewności detekcji (ang. \textit{confidence score}),
 \item[$N$] liczba wykrytych obiektów.
\end{description}
Miarą jakości dopasowania ramki predykowanej $B_p$ do ramki referencyjnej $B_{gt}$ jest współczynnik przecięcia nad sumą (IoU), zdefiniowany jako:
\begin{equation}
\mathrm{IoU}(B_p, B_{gt}) = \frac{\lvert B_p \cap B_{gt} \rvert}{\lvert B_p \cup B_{gt} \rvert}.
\end{equation}
Jeden owocnik, jedna ramka i etykieta, co ilustruje rys.~\ref{fig:class_vs_det_web}. Klasyfikator wydaje jeden werdykt na cały kadr; detektor opisuje każdy obiekt z osobna. W kępie opieńki (\textit{Armillaria mellea}), gdzie kilkanaście trzonów styka się w ściółce, tylko detekcja pozwala odróżnić jadalny owocnik od sąsiedniego, który nim nie jest.

\begin{figure}[htbp]
 \centering
 \includegraphics[width=0.9\textwidth]{classification_vs_detection_official}
 \caption{Porównanie klasyfikacji całego obrazu a detekcji wieloobiektowej. Źródło: opracowanie własne na podstawie \cite{redmon2016yolo, lin2014coco}.}
 \label{fig:class_vs_det_web}
\end{figure}

Trening od zera na $147$ gatunkach przekroczyłby budżet czasu i pamięci VRAM. Dlatego wagi startują z COCO, a dostrajanie obejmuje głowicę i wyższe warstwy szkieletu \cite{alhadhrami2018transfer, lin2014coco}. Na rys.~\ref{fig:transfer_learning_web} pokazano tę logikę: niskopoziomowe filtry krawędzi zostają, głowica uczy się mykologii. W pracy zastosowano \texttt{yolo11m.pt} (rozdz.~\ref{chap:implementacja}).

\begin{figure}[htbp]
 \centering
 \includegraphics[width=0.85\textwidth]{transfer_learning_official}
 \caption{Schemat uczenia transferowego z zamrożonym szkieletem. Źródło: \cite{alhadhrami2018transfer}.}
 \label{fig:transfer_learning_web}
\end{figure}

\subsection{Detektory dwuetapowe i jednoetapowe}
Dwa przejścia sieci lub jedno – tyle się sprowadza podział rodzin detekcyjnych. Faster R-CNN \cite{ren2015fasterrcnn} generuje najpierw propozycje regionów, potem dopiero klasyfikuje i prostuje ramki. Precyzja wysoka; koszt wysoki. Na serwerze z GPU to akceptowalny kompromis. Na Snapdragonie 8 Gen 2 w lesie rozwiązanie dwuetapowe odpada: setki megabajtów, dwa przebiegi na klatkę i thermal throttling po kilku minutach.

Detektory jednoetapowe – SSD \cite{liu2016ssd} i cała rodzina YOLO \cite{redmon2016yolo} – regresują ramki i klasy bezpośrednio na siatce obrazu w jednym przebiegu. Jeden tensor wyjściowy. Przewidywalne opóźnienie. Model mieszczący się w kilkudziesięciu megabajtach flash.

W tab.~\ref{tab:detektory} zestawiono obie rodziny. Na urządzeniu z procesorem Snapdragon jednoetapowy detektor lepiej spełnia budżet opóźnienia: jeden przebieg sieci, przewidywalny czas inferencji i mniejszy artefakt binarny.

\begin{table}[h!]
 \centering
 \caption{Porównanie detektorów dwuetapowych i jednoetapowych. Źródło: opracowanie własne na podstawie \cite{ren2015fasterrcnn, redmon2016yolo, liu2016ssd}.}
 \label{tab:detektory}
 \begin{tabular}{|l|c|c|}
 \hline
 \textbf{Kryterium} & \textbf{Dwuetapowe} & \textbf{Jednoetapowe} \\
 \hline \hline
 Reprezentatywny model & Faster R-CNN & YOLO, SSD \\
 \hline
 Liczba przebiegów sieci & dwa & jeden \\
 \hline
 Dokładność detekcji & bardzo wysoka & wysoka \\
 \hline
 Szybkość inferencji & niska & bardzo wysoka \\
 \hline
 Przydatność na urządzeniach mobilnych & ograniczona & wysoka \\
 \hline
 \end{tabular}
\end{table}

\subsection{Ewolucja architektury YOLO}
Redmon i wsp. w 2016 roku postawili jedną tezę: detekcja to regresja na siatce: jeden przebieg wystarczy \cite{redmon2016yolo}. Teza słuszna. Przez osiem lat kolejne generacje zmniejszały lukę wobec detektorów dwuetapowych bez rezygnacji z tej zasady: YOLOv3 dołożył predykcję wieloskalową \cite{redmon2018yolov3}, YOLOv8 usunął kotwice i ujednolicił API Ultralytics \cite{ultralytics2023yolov8}, YOLO11 przeprojektował szkielet i szyję \cite{ultralytics2024yolo11}.

Historia wersji jest tu mniej istotna niż jeden konkretny wniosek inżynierski: YOLO11m z $\approx20{,}1$~mln parametrów mieści się w budżecie TFLite FP16 ($38{,}58$~MB) i jest wystarczająco głęboki, by odróżniać 147 gatunków bliźniaczych – łącznie z parami takimi jak \textit{Amanita phalloides}/\textit{Amanita citrina} czy \textit{Galerina marginata}/\textit{Kuehneromyces mutabilis}. Nano jest za płytki: ablacja to dokumentuje (rozdz.~\ref{chap:badania}). Large przekracza budżet RAM i opóźnienia na Snapdragonie.

Współczesne architektury YOLO dzielą się funkcjonalnie na trzy zasadnicze części:
\begin{enumerate}
 \item Szkielet (ang. \textit{backbone}) jest siecią splotową odpowiedzialną za ekstrakcję hierarchicznych map cech z obrazu wejściowego.
 \item Szyja (ang. \textit{neck}) agreguje cechy z różnych poziomów rozdzielczości, najczęściej z wykorzystaniem struktury piramidy cech typu PANet.
 \item Głowica (ang. \textit{head}) generuje końcowe predykcje ramek oraz klas, w nowszych wersjach w wariancie bezkotwicowym (ang. \textit{anchor-free}).
\end{enumerate}

Oficjalny schemat Ultralytics (rys.~\ref{fig:yolo11_web}) pokazuje ten podział wprost: wejście $640\times640\times3$ przechodzi przez szkielet z blokami \texttt{Conv}/\texttt{C3k2}, następnie przez szyję z \texttt{SPPF}, \texttt{C2PSA}, \texttt{Upsample} i \texttt{Concat}, a trzy gałęzie \texttt{Detect} na skalach $80\times80$, $40\times40$ i $20\times20$ zasilają zadania Object Detection \cite{ultralytics2024yolo11}.

\begin{figure}[htbp]
 \centering
 \includegraphics[width=0.88\textwidth]{yolo11_arch_official}
 \caption{Struktura detektora YOLO11 z podziałem na Backbone, Neck i Decoupled Head. Źródło: \cite{ultralytics2024yolo11}.}
 \label{fig:yolo11_web}
\end{figure}

Trzyskalowa głowica nie jest ozdobą diagramu. To odpowiedź na geometrię owocników. Mapa $80\times80$ łapie drobne kapelusze i młode owocniki na ściółce; $40\times40$ obejmuje typowy rozmiar w kadrze telefonu; $20\times20$ obejmuje duże skupiska, np.\ kępy opieńki miodowej (\textit{Armillaria mellea}), gdzie wiele trzonów styka się i nakłada. Bez fuzji skal model albo sklejałby sąsiednie owocniki w jedną ramkę, albo gubił drobne okazy między liśćmi.

Szkielet oparto na blokach \texttt{C3k2} (ang. \textit{Cross Stage Partial with 2 kernels}) – wariancie CSP z dwoma ścieżkami cech i lekkimi splotami \cite{ultralytics2024yolo11}. Jedna ścieżka liczy, druga omija część transformacji; wynik jest konkatenowany. Mniej redundantnego gradientu, mniej parametrów, a tekstury blaszek i łusek nadal przechodzą wyżej. Moduł \texttt{SPPF} na dnie szkieletu zbiera kontekst przestrzenny sekwencyjnym max-poolingiem zamiast kilku równoległych jąder SPP. Taniej obliczeniowo przy podobnym zasięgu recepcji. Blok \texttt{C2PSA} w szyi dokłada uwagę przestrzenną przed rozesłaniem cech do trzech detektorów.

Głowica YOLO11 jest bezkotwicowa i rozłączna: regresja ramki i klasyfikacja gatunku idą osobnymi podsieciami na tej samej mapie cech. Wspólna głowica w starszych YOLO pchała gradient lokalizacji przeciwko gradientowi klasy – tu konflikt jest mniejszy \cite{ultralytics2024yolo11}. Brak kotwic oznacza też brak ręcznego doboru \textit{anchor boxes} pod wąskie trzony i płaskie kapelusze; sieć regresuje ramkę względem środka komórki siatki.

Trzy człony budują łączną stratę: geometria ramki, precyzja brzegów i klasyfikacja gatunkowa. Geometrię obsługuje Complete IoU (CIoU) – klasyczne IoU to za mało, bo nie karze za przesunięty środek ani złe proporcje \cite{ultralytics2024yolo11}:
\begin{equation}
\mathcal{L}_{\mathrm{CIoU}} = 1 - \mathrm{IoU} + \frac{\rho^2(b, b^{gt})}{c^2} + \alpha v,
\end{equation}
gdzie
\begin{description}[topsep=2pt,itemsep=-3pt]
 \item[$\rho^2(b, b^{gt})$] kwadrat odległości euklidesowej pomiędzy środkami ramki predykowanej $b$ i referencyjnej $b^{gt}$,
 \item[$c$] przekątna najmniejszego prostokąta otaczającego obie ramki,
 \item[$v$] człon penalizujący niezgodność współczynnika proporcji (ang. \textit{aspect ratio}),
\end{description}
\begin{equation}
v = \frac{4}{\pi^2} \left(\arctan\frac{w^{gt}}{h^{gt}} - \arctan\frac{w}{h} \right)^2,
\end{equation}
gdzie $\alpha$ jest współczynnikiem wagowym równoważącym wpływ członu $v$ względem członu przecięcia, wyznaczanym adaptacyjnie jako:
\begin{equation}
\alpha = \frac{v}{(1-\mathrm{IoU}) + v}.
\end{equation}
Trzy sygnały w jednej liczbie. Człon $v$ z arcustangensami rozróżnia smukły trzon maślaka od płaskiego kapelusza koźlarza. Dokładnie ta geometria decyduje o parach bliźniaczych. Gdy ramki się nie nakładają, gradient klasycznego IoU zanika; CIoU podtrzymuje sygnał przez $\rho^2/c^2$ i skraca pierwsze, powolne epoki treningu \cite{ultralytics2024yolo11}. Geometrię $\rho$ i $c$ ilustruje rys.~\ref{fig:ciou_geometria}.

\begin{figure}[htbp]
 \centering
 \begin{tikzpicture}[scale=1.0]
 % najmniejszy prostokat otaczajacy obie ramki
 \draw[dashed, gray] (0,0) rectangle (6.2,3.8);
 \draw[dashed, gray, thick] (0,0) -- (6.2,3.8);
 \node[gray, font=\small] at (4.55,1.35) {$c$};

 % ramka referencyjna
 \draw[EEUltramarine, very thick] (0,0) rectangle (3.6,2.6);
 \fill[EEUltramarine] (1.8,1.3) circle (2pt);
 \node[EEUltramarine, font=\small, anchor=south west] at (0.1,0.1) {$b^{gt}$};

 % ramka predykowana
 \draw[EETangerine, very thick] (2.6,1.2) rectangle (6.2,3.8);
 \fill[EETangerine] (4.4,2.5) circle (2pt);
 \node[EETangerine, font=\small, anchor=north east] at (6.1,3.7) {$b$};

 % odleglosc srodkow
 \draw[-latex, semithick] (1.8,1.3) -- (4.4,2.5);
 \node[font=\small, anchor=south east] at (3.35,1.78) {$\rho$};

 % wymiary
 \draw[|-|, gray] (0,-0.35) -- node[below, font=\scriptsize, gray] {$w^{gt}$} (3.6,-0.35);
 \draw[|-|, gray] (-0.35,0) -- node[left, font=\scriptsize, gray] {$h^{gt}$} (-0.35,2.6);
 \draw[|-|, gray] (2.6,4.15) -- node[above, font=\scriptsize, gray] {$w$} (6.2,4.15);
 \draw[|-|, gray] (6.55,1.2) -- node[right, font=\scriptsize, gray] {$h$} (6.55,3.8);
 \end{tikzpicture}
 \caption{Interpretacja geometryczna składników miary CIoU: odległość środków ramek $\rho$, przekątna $c$ najmniejszego prostokąta otaczającego obie ramki oraz wymiary wykorzystywane w członie proporcji $v$. Źródło: opracowanie własne.}
 \label{fig:ciou_geometria}
\end{figure}

Precyzję krawędzi ramki podnosi DFL (\textit{Distribution Focal Loss}): zamiast regresować jedną liczbę, sieć modeluje każdą współrzędną jako rozkład dyskretny na sąsiednich węzłach \cite{ultralytics2024yolo11}:
\begin{equation}
\mathcal{L}_{\mathrm{DFL}}(\mathcal{S}_i, \mathcal{S}_{i+1}) = -\big(y_{i+1}-y\big)\log(\mathcal{S}_i) - \big(y - y_i\big)\log(\mathcal{S}_{i+1}),
\end{equation}
$y$ to rzeczywista współrzędna krawędzi, $y_i$/$y_{i+1}$ – sąsiednie węzły dyskretyzacji, $\mathcal{S}_i$/$\mathcal{S}_{i+1}$ – prognozowane prawdopodobieństwa przynależności. Interpolacja liniowa odtwarza wartość ciągłą z zachowaniem różniczkowalności. CIoU i DFL tworzą człon lokalizacyjny; Binary Cross-Entropy obsługuje klasyfikację gatunkową. Trzy składniki, jedna funkcja celu.

Inferencja: YOLO11m produkuje tysiące nakładających się kandydatów. NMS iteracyjnie wybiera ramkę o najwyższej pewności, odrzuca wszystkie pokrywające się z nią powyżej progu IoU i powtarza. Jeden owocnik. Jedna ramka wynikowa.

\section{Metody optymalizacji modeli dla urządzeń mobilnych}
Model wytrenowany na RTX~4070 musi działać na Snapdragonie bez serwera, bez chłodzenia i bez degradacji po godzinie w terenie. To nie jest proste skalowanie. Poniżej opisano ograniczenia sprzętowe telefonu, logikę kwantyzacji oraz drogę formatu modelu od PyTorch do pliku \texttt{.tflite} w pamięci flash.

\subsection{Ograniczenia sprzętowe urządzeń mobilnych}
Snapdragon nie jest RTX-em. Urządzenie mobilne dysponuje kilkoma gigabajtami RAM, ograniczoną przepustowością magistrali pamięci i nie posiada aktywnego chłodzenia – SoC po kilku minutach intensywnych obliczeń wchodzi w thermal throttling. Model musi zatem być mały i rzadko aktywowany. Stąd decyzja o FP16 i ograniczeniu inferencji do 4~FPS.

\subsection{Kwantyzacja potreningowa}
Podstawową techniką dostosowania wytrenowanego modelu do platformy mobilnej jest kwantyzacja potreningowa (ang. \textit{post-training quantization}, PTQ), polegająca na zastąpieniu 32-bitowej reprezentacji zmiennoprzecinkowej wag i aktywacji reprezentacją o mniejszej liczbie bitów \cite{jacob2018quantization}. Odwzorowanie wartości rzeczywistej na wartość skwantowaną opisuje przekształcenie afiniczne:
\begin{equation}
r = S\,(q - Z),
\end{equation}
gdzie
\begin{description}[topsep=2pt,itemsep=-3pt]
 \item[$r$] pierwotna wartość zmiennoprzecinkowa,
 \item[$q$] odpowiadająca jej wartość całkowitoliczbowa,
 \item[$S$] dodatni współczynnik skali (ang. \textit{scale}),
 \item[$Z$] przesunięcie punktu zerowego (ang. \textit{zero-point}).
\end{description}
Współczynnik skali dla reprezentacji $n$-bitowej wyznacza się na podstawie zakresu wartości pierwotnych zgodnie z zależnością:
\begin{equation}
S = \frac{r_{\max} - r_{\min}}{2^{n} - 1}.
\end{equation}

Kwantyzacja całkowitoliczbowa, opisana powyższym przekształceniem afinicznym, różni się fundamentalnie od redukcji precyzji w obrębie samej reprezentacji zmiennoprzecinkowej. Liczba zmiennoprzecinkowa w formacie IEEE~754 jest reprezentowana jako:
\begin{equation}
x = (-1)^{s} \cdot 2^{\,e-\mathrm{bias}} \cdot \left(1 + \frac{m}{2^{p}}\right),
\end{equation}
gdzie $s$ oznacza bit znaku, $e$ pole wykładnika, $m$ pole mantysy o precyzji $p$~bitów, a $\mathrm{bias}$ jest przesunięciem wykładnika. FP32: $p=23$ ($\approx1{,}19\times10^{-7}$ względnie). FP16: $p=10$ ($\approx9{,}77\times10^{-4}$). W obu przypadkach błąd rośnie z wielkością liczby. INT8 ma stały krok $S$ na cały zakres warstwy. Małe aktywacje giną w kwancie.

Nie jest to wyłącznie kwestia akademicka. Gradienty jasności na blaszkach (\textit{lamellae}), pierścieniu i łuskach kapelusza to właśnie te „małe” sygnały. One oddzielają jadalną czubajkę od trującej czubajeczki i młodego muchomora od pieczarki. INT8 spłaszcza je do stałego kwantu $S$ warstwy. FP16 zachowuje skalowanie wykładnicze, więc te detale nie giną w siatce całkowitoliczbowej.

Dlatego zrezygnowano z pełnej kwantyzacji INT8. Decyzja jest projektowa, nie empiryczna: w pracy nie raportowano $mAP$ ani opóźnienia wariantu INT8. Szacunkowy rozmiar artefaktu INT8 ($\approx20$\,MB) i zysk na NPU wynikają z szerokości reprezentacji, nie z pomiaru na telefonie. FP16 dla YOLO11m daje $38{,}58$\,MB – połowa FP32 – przy zachowaniu drobnych cech teksturalnych potrzebnych do par bliźniaczych. Stąd decyzja produkcyjna: tylko FP16.

Różnicę w podziale przestrzeni bitowej pomiędzy trzema rozważanymi reprezentacjami zilustrowano na rys.~\ref{fig:bity_formatow}. Redukcja FP32~$\rightarrow$~FP16 zachowuje strukturę zmiennoprzecinkową i skraca zarówno pole wykładnika, jak i mantysy, natomiast przejście do INT8 całkowicie eliminuje pole wykładnika, zastępując skalowanie wykładnicze pojedynczym, stałym dla całej warstwy współczynnikiem $S$.

\begin{figure}[htbp]
 \centering
 \begin{tikzpicture}[
 pole/.style={draw, minimum height=7mm, align=center, font=\scriptsize, inner sep=1pt},
 znak/.style={pole, fill=black!10},
 wykl/.style={pole, fill=EEBlue!25},
 mant/.style={pole, fill=EETangerine!30},
 calk/.style={pole, fill=EEUltramarine!20}
 ]
 \def\jb{0.34}
 % FP32
 \node[font=\small, anchor=east] at (-0.25,0) {FP32};
 \node[znak, minimum width=\jb cm] (a1) at (\jb/2, 0) {$s$};
 \node[wykl, minimum width=8*\jb cm] (a2) at (\jb+8*\jb/2, 0) {wykładnik: $8$ b};
 \node[mant, minimum width=23*\jb cm] (a3) at (9*\jb+23*\jb/2, 0) {mantysa: $23$ b};
 \node[font=\scriptsize, anchor=west] at (32*\jb+0.2, 0) {$32$ b};

 % FP16
 \node[font=\small, anchor=east] at (-0.25,-1.1) {FP16};
 \node[znak, minimum width=\jb cm] at (\jb/2, -1.1) {$s$};
 \node[wykl, minimum width=5*\jb cm] at (\jb+5*\jb/2, -1.1) {wykł.: $5$ b};
 \node[mant, minimum width=10*\jb cm] at (6*\jb+10*\jb/2, -1.1) {mantysa: $10$ b};
 \node[font=\scriptsize, anchor=west] at (16*\jb+0.2, -1.1) {$16$ b};

 % INT8
 \node[font=\small, anchor=east] at (-0.25,-2.2) {INT8};
 \node[calk, minimum width=8*\jb cm] at (8*\jb/2, -2.2) {liczba całkowita: $8$ b};
 \node[font=\scriptsize, anchor=west] at (8*\jb+0.2, -2.2) {$8$ b $+$ wspólne $S$, $Z$};
 \end{tikzpicture}
 \caption{Podział przestrzeni bitowej w formatach FP32, FP16 oraz INT8. Źródło: opracowanie własne na podstawie \cite{jacob2018quantization}.}
 \label{fig:bity_formatow}
\end{figure}

Ilościowe zestawienie właściwości numerycznych trzech rozważanych reprezentacji wag przedstawiono w tab.~\ref{tab:kwantyzacja}. Zestawienie obejmuje szerokość pola mantysy, wynikającą z niej względną precyzję reprezentacji, charakter propagacji błędu kwantyzacji oraz szacunkowy rozmiar artefaktu modelu YOLO11m w danym formacie. Wpływ kwantyzacji na jakość decyzji klasyfikacyjnych ocenia się w kategoriach macierzy pomyłek, której elementy ($TP$, $FP$, $FN$, $TN$) zdefiniowano na rys.~\ref{fig:conf_matrix_web}; miary te wykorzystano później przy ewaluacji bezpieczeństwa mykologicznego w rozdziale \ref{chap:badania}.

\begin{figure}[htbp]
 \centering
 \includegraphics[width=0.75\textwidth]{confusion_matrix_official}
 \caption{Definicje macierzy pomyłek (TP, FP, FN, TN). Źródło: \cite{fawcett2006roc}.}
 \label{fig:conf_matrix_web}
\end{figure}

\begin{table}[htbp]
 \centering
 \caption{Porównanie formatów reprezentacji numerycznej wag rozważanych na etapie konwersji modelu. Źródło: opracowanie własne na podstawie \cite{jacob2018quantization, tflite}.}
 \label{tab:kwantyzacja}
 \begin{tabular}{|p{4.6cm}|c|c|c|}
 \hline
 \textbf{Właściwość} & \textbf{FP32} & \textbf{FP16} & \textbf{INT8} \\
 \hline \hline
 Szerokość słowa & $32$\,b & $16$\,b & $8$\,b \\
 \hline
 Pole mantysy $p$ & $23$\,b & $10$\,b & brak \\
 \hline
 Precyzja względna & $1{,}19\cdot10^{-7}$ & $9{,}77\cdot10^{-4}$ & krok stały $S$ \\
 \hline
 Skalowanie błędu & proporcjonalne & proporcjonalne & absolutne \\
 \hline
 Rozmiar TFLite (YOLO11m) & ok.\ $77$\,MB & $38{,}58$\,MB & ok.\ $20$\,MB \\
 \hline
 Wymaga kalibracji & nie & nie & tak \\
 \hline
 Zastosowanie w pracy & referencja & \textbf{produkcja} & odrzucony \\
 \hline
 \end{tabular}
\end{table}

\subsection{Formaty pośrednie i docelowe modelu}
PyTorch nie trafia na telefon wprost. Graf obliczeniowy przechodzi przez dwa etapy konwersji. ONNX działa jak uniwersalny pośrednik, oddziela środowisko treningowe od docelowego i pozwala przełączać się między frameworkami \cite{onnx}. Z ONNX albo bezpośrednio z SavedModel konwerter \texttt{tf.lite.TFLiteConverter} kompiluje binarny artefakt TFLite, który Interpreter ładuje w czasie wykonania na urządzeniu \cite{tflite}. Tu stosuje się kwantyzację FP16. Schemat ścieżki SavedModel do TFLite FP16 pokazuje rys.~\ref{fig:tflite_flow_web}; wynikowy artefakt \texttt{.tflite} trafia bezpośrednio do aktywów aplikacji Android (rozdz.~\ref{chap:implementacja}).

\begin{figure}[htbp]
 \centering
 \includegraphics[width=0.9\textwidth]{tflite_conversion_flow_official}
 \caption{Schemat konwersji SavedModel do TFLite FP16. Źródło: \cite{tflite}.}
 \label{fig:tflite_flow_web}
\end{figure}


\chapter{Projekt systemu rozpoznawania}\label{chap:projekt_systemu}
\section{Wymagania funkcjonalne i niefunkcjonalne}
Punkt wyjścia projektu to jeden scenariusz: grzybiarz w Puszczy Białowieskiej, brak LTE, pełny kosz i telefon w dłoni. Wymagania wynikają z tego scenariusza wprost. Funkcjonalne opisują, co aplikacja robi dla użytkownika; niefunkcjonalne określają, ile zasobów można na to zużyć.

\subsection{Wymagania funkcjonalne}
Interakcja użytkownika z systemem zamyka się w kilku przypadkach użycia. Aplikacja musi spełniać pięć wymagań funkcjonalnych:
\begin{enumerate}
 \item Uruchamiać podgląd obrazu z kamery w czasie rzeczywistym jako główny widok roboczy.
 \item Automatycznie wykrywać wiele owocników jednocześnie w kadrze bez ręcznego wskazywania obiektów przez użytkownika.
 \item Nakładać ramki otaczające (BBOX) z etykietami gatunkowymi w wersji polskiej i łacińskiej.
 \item Generować wyraźne, krytyczne ostrzeżenia wizualne przy wykryciu gatunku trującego, wyróżnione kontrastową barwą względem detekcji jadalnych.
 \item Aktualizować wyniki detekcji na bieżąco wraz ze zmianą obserwowanej sceny.
\end{enumerate}
Reguła nadrzędna: priorytet czułości dla klas toksycznych. Ostrzeżenie o gatunku trującym wyprzedza komfort użytkowania. Lepiej oznaczyć jadalną czubajkę kania jako podejrzaną niż przeoczyć muchomora zielonego. Zakres interakcji przedstawia rys.~\ref{fig:przypadki_uzycia}.

\begin{figure}[htbp]
 \centering
 \begin{tikzpicture}[
 pu/.style={draw, ellipse, align=center, font=\scriptsize, minimum width=3.2cm,
 minimum height=1.0cm, inner sep=1pt, fill=EEBlue!10},
 rel/.style={-latex, semithick},
 ext/.style={dashed, -latex, semithick, gray}
 ]
 % aktor
 \draw[semithick] (0,0) circle (0.22);
 \draw[semithick] (0,-0.22) -- (0,-0.9);
 \draw[semithick] (-0.4,-0.45) -- (0.4,-0.45);
 \draw[semithick] (0,-0.9) -- (-0.3,-1.5);
 \draw[semithick] (0,-0.9) -- (0.3,-1.5);
 \node[font=\scriptsize, align=center] at (0,-1.9) {Grzybiarz\\(użytkownik)};

 % granica systemu
 \draw[gray, rounded corners=3pt] (2.0,-4.3) rectangle (10.6,2.1);
 \node[gray, font=\scriptsize, anchor=north west] at (2.15,2.0) {Aplikacja mobilna};

 \node[pu] (skan) at (5.0,0.9) {Skanowanie owocników\\w podglądzie kamery};
 \node[pu] (ostrz) at (5.0,-0.4) {Odbiór ostrzeżenia\\o gatunku trującym};
 \node[pu] (zapis) at (5.0,-1.7) {Zapis znaleziska\\do historii};
 \node[pu] (hist) at (5.0,-3.0) {Przegląd historii\\znalezisk};

 \draw[rel] (0.45,-0.3) -- (skan);
 \draw[rel] (0.45,-0.6) -- (ostrz);
 \draw[rel] (0.45,-1.1) -- (zapis);
 \draw[rel] (0.45,-1.5) -- (hist);

 \node[pu, fill=EETangerine!18] (offl) at (8.6,0.9) {Praca offline};
 \node[pu, fill=EETangerine!18] (kadr) at (8.6,-1.7) {Kadrowanie\\i zapis miniatury};
 \draw[ext] (offl) -- node[right, font=\scriptsize, gray, xshift=1pt] {} (skan);
 \draw[ext] (kadr) -- (zapis);
 \node[gray, font=\scriptsize, anchor=west] at (6.7,1.55) {};
 \end{tikzpicture}
 \caption{Diagram przypadków użycia aplikacji mobilnej; relacje przerywane oznaczają zależności rozszerzające. Źródło: opracowanie własne.}
 \label{fig:przypadki_uzycia}
\end{figure}

\subsection{Wymagania niefunkcjonalne}
Telefon w lesie, bez zasięgu, bateria w połowie. Cztery ograniczenia twarde:
\begin{enumerate}
 \item Praca offline jest bezwzględna: zero odwołań sieciowych; aplikacja działa w martwej strefie GSM/LTE bez degradacji funkcji.
 \item Rozmiar artefaktu pozostaje poniżej $40$\,MB, tak aby model mieścił się w pamięci flash i nie groził wyrzutem procesu przez Android Low-Memory Killer.
 \item Minimalny poziom API to 24 (Android 7.0), co obejmuje zdecydowaną większość urządzeń w obiegu.
 \item Niski pobór energii: ograniczenie tempa inferencji do $4$\,FPS zmniejsza nagrzewanie i chroni baterię podczas kilkugodzinnej pracy w terenie.
\end{enumerate}

Cztery ograniczenia, trzy budżety inżynierskie.

\textbf{Budżet opóźnienia}: interwał throttlingu kamery wynosi $250$\,ms ($4$\,FPS). Do tego w projekcie doliczono rezerwę $50$\,ms na przygotowanie klatki CameraX (RGBA, letterbox $640\times640$), co daje projektowy limit $300$\,ms na \emph{cały} cykl. Zmierzono jednak wyłącznie \texttt{detect()} (rozdz.~\ref{chap:badania}): na S23 Ultra średnia to $156{,}5$\,ms na zimnym SoC i $245{,}9$\,ms po obciążeniu, a $p_{95}$ serii gorącej wynosi $278{,}3$\,ms. Czas pobrania i letterboxingu nie wchodzi w te liczby, więc z pomiaru nie wynika, że pełny cykl zawsze mieści się poniżej $300$\,ms.

\textbf{Budżet binarny modelu}: poniżej $40$\,MB. Artefakt TFLite FP16 YOLO11m zajmuje $38{,}58$\,MB. Margines $1{,}42$\,MB chroni przed limitami Android App Bundle i zmniejsza obciążenie transferu przy instalacji w terenie.

\textbf{Budżet RAM}: poniżej $250$\,MB procesu. Składowe: wagi rezydentne ${\approx}39$\,MB + bufor wejściowy $640{\times}640{\times}3{\times}4$\,B~${\approx}4{,}69$\,MB + bufor wyjściowy $151{\times}8400{\times}4$\,B~${\approx}4{,}84$\,MB + bitmapa podglądu CameraX/Compose + narzut ART. Łącznie margines względem domyślnego limitu sterty API~24+ eliminuje ryzyko wyrzutu procesu przez Low-Memory Killer podczas kilkugodzinnej sesji terenowej.

Zbiorcze zestawienie wyprowadzonych budżetów inżynierskich wraz z ich składowymi oraz sposobem weryfikacji przedstawiono w tab.~\ref{tab:budzety}.

\begin{table}[htbp]
 \centering
 \caption{Budżety inżynierskie systemu wyprowadzone z wymagań niefunkcjonalnych. Źródło: opracowanie własne.}
 \label{tab:budzety}
 \begin{tabular}{|p{2.9cm}|c|p{4.4cm}|p{3.1cm}|}
 \hline
 \textbf{Budżet} & \textbf{Limit} & \textbf{Składowe} & \textbf{Weryfikacja} \\
 \hline \hline
 Opóźnienie \texttt{detect()} & interwał $250$\,ms ($4$\,FPS) & zmierzono samą inferencję TFLite, nie pełny cykl CameraX & tab.~\ref{tab:latency_s23} \\
 \hline
 Rozmiar binarny modelu & $<40$\,MB & artefakt TFLite FP16 $38$--$39$\,MB $+$ margines pakietu & pomiar rozmiaru pliku \\
 \hline
 Pamięć operacyjna RAM & $<250$\,MB & wagi $\approx39$\,MB, bufor wejściowy $4{,}69$\,MB, bufor wyjściowy $4{,}84$\,MB, podgląd i ART & profilowanie pamięci procesu \\
 \hline
 \end{tabular}
\end{table}

Trzy budżety, jedno źródło: YOLO11m w formacie TFLite FP16. Dlaczego akurat ten model – uzasadnia kolejny podrozdział.

Sprzęt użyty do treningu i testów mobilnych zestawia tab.~\ref{tab:sprzet}. Stos programowy przedstawiono w tab.~\ref{tab:stos_programowy}. Obliczenia treningowe wykonywano wyłącznie na stacji z CUDA; weryfikację inferencji krawędziowej przeprowadzono na fizycznym urządzeniu mobilnym.

\begin{table}[htbp]
 \centering
 \caption{Specyfikacja środowiska sprzętowego wykorzystanego w pracy. Źródło: opracowanie własne.}
 \label{tab:sprzet}
 \begin{tabular}{|p{3.4cm}|p{4.6cm}|p{4.6cm}|}
 \hline
 \textbf{Komponent} & \textbf{Stacja robocza (trening)} & \textbf{Platforma mobilna (test)} \\
 \hline \hline
 Jednostka centralna & AMD Ryzen~5~3600, $6$/$12$, $3{,}60$\,GHz & Qualcomm Snapdragon~8 Gen~2 \\
 \hline
 Pamięć operacyjna & $32$\,GB DDR4 $3200$\,MHz & $12$\,GB LPDDR5X \\
 \hline
 Akcelerator & NVIDIA GeForce RTX~4070, $12$\,GB VRAM & Adreno~740 (delegat GPU) \\
 \hline
 Pamięć masowa & $1$\,TB NVMe PCIe Gen4 & pamięć wewnętrzna UFS \\
 \hline
 System operacyjny & Windows~11 & Android~16 \\
 \hline
 \end{tabular}
\end{table}

\begin{table}[htbp]
 \centering
 \caption{Stos programowy zastosowany na poszczególnych etapach realizacji projektu. Źródło: opracowanie własne.}
 \label{tab:stos_programowy}
 \begin{tabular}{|p{3.3cm}|p{5.0cm}|p{4.3cm}|}
 \hline
 \textbf{Etap potoku} & \textbf{Technologia} & \textbf{Rola w systemie} \\
 \hline \hline
 Akwizycja danych & Python $3.11$, \texttt{requests.Session} & pobieranie obserwacji iNaturalist \\
 \hline
 Auto-anotacja & GroundingDINO, \texttt{autodistill} & lokalizacja zero-shot owocników \\
 \hline
 Trening & PyTorch $2.x$, CUDA~$12.x$, Ultralytics & uczenie transferowe YOLO11m \\
 \hline
 Konwersja & TensorFlow, \texttt{tf.lite.TFLiteConverter} & eksport SavedModel do TFLite FP16 \\
 \hline
 Inferencja mobilna & TensorFlow Lite Interpreter, GpuDelegate & lokalna detekcja na urządzeniu \\
 \hline
 Interfejs użytkownika & Kotlin, Jetpack Compose, CameraX & podgląd, nakładka AR, gesty \\
 \hline
 Persystencja & Room (SQLite), Kotlin Coroutines & historia znalezisk \\
 \hline
 \end{tabular}
\end{table}

Dobór bibliotek wynika z podziału na stację roboczą (akwizycja, anotacja, trening, konwersja) oraz telefon (inferencja TFLite, interfejs Compose, persystencja Room). Dzięki temu graf uczenia nie trafia na urządzenie mobilne, a aplikacja nie zależy od Pythona w czasie wykonania.

\section{Architektura systemu i dobór sieci neuronowej}
Budżety z poprzedniego podrozdziału wymuszają konkretne decyzje architektoniczne: jeden przebieg sieci, artefakt poniżej 40 MB, praca offline. Poniżej opisano wynikającą z tych decyzji strukturę systemu i model, który je spełnia.

\subsection{Trójwarstwowa architektura systemu}
Trzy moduły, trzy odpowiedzialności. CameraX dostarcza klatki i skaluje je do $640\times640$. Interpreter TFLite uruchamia model FP16 i zwraca surowy tensor. Kotlin odrzuca słabe kandydatury, rysuje ramki AR i wyzwala ostrzeżenie toksykologiczne.

Granica między przygotowaniem modelu a czasem wykonania przebiega na artefakcie \texttt{.tflite}. Karta graficzna pracuje tylko raz: trening i konwersja na stacji roboczej. Telefon wykonuje wyłącznie inferencję i logikę decyzyjną -- żadnego połączenia sieciowego, żadnego serwera. Na rys.~\ref{fig:warstwy} pokazano ten podział.

\begin{figure}[htbp]
 \centering
 \begin{tikzpicture}[
 w/.style={draw, rounded corners=3pt, align=center, font=\small,
 text width=12.4cm, inner sep=5pt},
 modul/.style={draw, rounded corners=2pt, align=center, font=\scriptsize,
 text width=3.6cm, minimum height=8mm, inner sep=2pt},
 strz/.style={-latex, thick}
 ]
 % warstwa 1
 \node[w, fill=black!5] (w1) at (0,0) {\textbf{Warstwa przygotowania danych i uczenia} \\[2pt] {\footnotesize stacja robocza: AMD Ryzen 5 3600, NVIDIA RTX 4070, Python}};
 \node[modul, fill=EEBlue!12] (m11) at (-4.2,-1.15) {Akwizycja\\i walidacja obrazów};
 \node[modul, fill=EEBlue!12] (m12) at (0,-1.15) {Auto-anotacja\\GroundingDINO};
 \node[modul, fill=EEBlue!12] (m13) at (4.2,-1.15) {Trening YOLO11m\\(PyTorch, AMP)};

 % warstwa 2
 \node[w, fill=black!5] (w2) at (0,-2.9) {\textbf{Warstwa konwersji modelu} \\[2pt] {\footnotesize PyTorch $\rightarrow$ TensorFlow SavedModel $\rightarrow$ TensorFlow Lite FP16}};
 \node[modul, fill=EETangerine!22] (m21) at (0,-4.05) {Artefakt \texttt{.tflite} FP16, $38{,}58$\,MB};

 % warstwa 3
 \node[w, fill=black!5] (w3) at (0,-5.8) {\textbf{Warstwa wykonawcza (urządzenie mobilne)} \\[2pt] {\footnotesize Android, Kotlin, Jetpack Compose, CameraX, TFLite Interpreter}};
 \node[modul, fill=EEUltramarine!14] (m31) at (-4.2,-6.95) {Akwizycja klatek\\CameraX};
 \node[modul, fill=EEUltramarine!14] (m32) at (0,-6.95) {Inferencja\\i dekodowanie tensora};
 \node[modul, fill=EEUltramarine!14] (m33) at (4.2,-6.95) {Prezentacja, ostrzeżenia\\i baza Room};

 \draw[strz] (m12) -- (m13);
 \draw[strz] (m11) -- (m12);
 \draw[strz] (m13) -- ++(0,-0.75) -- (4.2,-2.15) -- (0,-2.15) -- (w2);
 \draw[strz] (m21) -- ++(0,-0.75) -- (0,-5.05) -- (w3);
 \draw[strz] (m31) -- (m32);
 \draw[strz] (m32) -- (m33);
 \end{tikzpicture}
 \caption{Podział systemu na warstwę przygotowania danych i uczenia, warstwę konwersji modelu oraz warstwę wykonawczą urządzenia mobilnego. Źródło: opracowanie własne.}
 \label{fig:warstwy}
\end{figure}
Przepływ danych przez system to jednokierunkowy ośmioetapowy potok. Klasa \texttt{ImageAnalysis} CameraX dostarcza klatki ze strategią \texttt{STRATEGY\_KEEP\_ONLY\_LATEST}. Starsze klatki są odrzucane, zanim zdążą się kolejkować. \texttt{Mutex.tryLock()} ogranicza tempo do $4$\,FPS w trybie standardowym ($250$\,ms) lub do $1{,}5$\,s \textit{heartbeat} po ustabilizowaniu sceny. Interpreter TFLite uruchamia model FP16 i zwraca tensor $[1,\,151,\,8400]$: 4 współrzędne ramki i 147 prawdopodobieństw klas dla każdego z $8400$ punktów predykcyjnych na mapach cech. Fast-fail obcina listę do 50 najwyżej ocenionych kandydatów; NMS ($\mathrm{IoU}>0{,}45$) usuwa nadmiary. Należy zaznaczyć, że zaimplementowany w aplikacji mechanizm NMS (\texttt{MushroomDetector.kt}) jest z zasady niewrażliwy na klasę (ang. \textit{class-agnostic}). Oznacza to, że pętla porównuje wyłącznie IoU geometryczny bez sprawdzania indeksu klasy. Ramka gatunku B o niższym wyniku pewności zostanie bezwzględnie usunięta, jeśli jej IoU z nakładającą się ramką gatunku A przekroczy zadany próg. Jest to celowe uproszczenie ukierunkowane na wydajność mobilną.

Osiem etapów przetwarzania pojedynczej klatki ilustruje rys.~\ref{fig:potok_architektura}.

\begin{figure}[htbp]
 \centering
 \begin{tikzpicture}[
 blok/.style={draw, rounded corners=2pt, align=center, text width=9.2cm,
 inner sep=3pt, font=\small, minimum height=7.5mm},
 akw/.style={blok, fill=EEBlue!12},
 inf/.style={blok, fill=EEUltramarine!12},
 bez/.style={blok, fill=EETangerine!18},
 strz/.style={-latex, thick}
 ]
 \node[akw] (n1) at (0,0) {1. CameraX \texttt{ImageAnalysis} – \texttt{STRATEGY\_KEEP\_ONLY\_LATEST}};
 \node[akw] (n2) at (0,-1.25) {2. Ograniczanie częstotliwości \texttt{Mutex.tryLock()}, $250$\,ms ($4$\,FPS) lub $1{,}5$\,s};
 \node[inf] (n3) at (0,-2.50) {3. Interpreter TensorFlow Lite – model YOLO11m FP16};
 \node[inf] (n4) at (0,-3.75) {4. Surowy tensor wyjściowy $[1,\,151,\,8400]$};
 \node[inf] (n5) at (0,-5.00) {5. Szybka eliminacja (\textit{fast-fail}), top-$50$ kandydatów};
 \node[inf] (n6) at (0,-6.25) {6. Tłumienie niemaksymalne (NMS), $\mathrm{IoU}>0{,}45$};
 \node[bez] (n7) at (0,-7.50) {7. Bufor konsensusu klatek, $4$ z $5$ ostatnich klatek};
 \node[bez] (n8) at (0,-8.75) {8. Asymetryczny filtr bezpieczeństwa – $0{,}18$ oraz $0{,}60$};
 \node[blok, fill=black!6] (n9) at (0,-10.00) {Nakładka AR (Jetpack Compose) oraz baza Room (\texttt{MushroomFind})};

 \draw[strz] (n1) -- (n2);
 \draw[strz] (n2) -- (n3);
 \draw[strz] (n3) -- (n4);
 \draw[strz] (n4) -- (n5);
 \draw[strz] (n5) -- (n6);
 \draw[strz] (n6) -- (n7);
 \draw[strz] (n7) -- (n8);
 \draw[strz] (n8) -- (n9);
 \end{tikzpicture}
 \caption{Ośmioetapowy potok przetwarzania pojedynczej klatki obrazu w zaprojektowanym systemie. Źródło: opracowanie własne.}
 \label{fig:potok_architektura}
\end{figure}

Przefiltrowane detekcje trafiają do bufora konsensusu klatek (\textit{multi-frame voting buffer}): gatunek uważa się za potwierdzony wyłącznie wtedy, gdy dominuje w co najmniej $4$ z $5$ ostatnich klatek. Bufor wygasza jednoklatkowe fluktuacje predykcji i stanowi mechanizm stabilizacji czasowej nakładki, nie składnik hipotezy $H_2$. Potwierdzone detekcje przechodzą przez asymetryczny filtr progów: $0{,}18$ dla gatunków trujących, $0{,}60$ dla jadalnych. Wynik trafia równolegle do nakładki AR Jetpack Compose, rysującej ramki i etykiety, oraz – przy ręcznym zapisie – do bazy Room, gdzie encja \texttt{MushroomFind} przechowuje gatunek, kategorię toksyczności, wartość pewności, znacznik czasu i ścieżkę wycinka obrazu. Polityka progów dotyczy warstwy wyświetlania: ostrzeżenie toksyczne od $0{,}18$, asercja jadalności od $0{,}60$; nie implikuje wyższego $R$ detektora na klasach trujących.

\subsection{Uzasadnienie wyboru modelu YOLO11 Medium}
Wybór modelu to czterowymiarowy kompromis: dokładność, rozmiar binarny, opóźnienie inferencji i zużycie energii, wszystko przy spełnieniu budżetów z poprzedniego podrozdziału. Przeanalizowano cztery alternatywy:
\begin{enumerate}
 \item YOLOv8n, reprezentujący wcześniejszą generację rodziny YOLO, o gorszym stosunku dokładności do liczby parametrów niż jego następca \cite{ultralytics2023yolov8}.
 \item Lżejszy wariant YOLO11 Nano (\texttt{YOLO11n}), atrakcyjny ze względu na minimalny rozmiar artefaktu TFLite, lecz o niedostatecznej pojemności reprezentacyjnej w zadaniu rozróżniania wielu gatunków bliźniaczych \cite{ultralytics2024yolo11}.
 \item Detektor MobileNet SSD, o niższej precyzji detekcji, niewystarczającej do rozróżniania gatunków bliźniaczych \cite{liu2016ssd}.
 \item Rodzina EfficientDet, której czas inferencji jest zbyt długi względem założeń pracy w reżimie czasu rzeczywistego.
\end{enumerate}
Wybrano YOLO11 Medium (\texttt{YOLO11m}) \cite{ultralytics2024yolo11}. $\approx20{,}1$~mln parametrów daje pojemność wystarczającą do modelowania subtelnych cech morfologicznych $147$ gatunków, a konwersja do TFLite FP16 ($38{,}58$~MB) mieści artefakt w pamięci flash telefonu. Na S23 Ultra zmierzono czas \texttt{detect()} (rozdz.~\ref{chap:badania}); pełnego cyklu CameraX nie profilowano. Wariant Nano, mimo mniejszego pliku, nie radził sobie z gatunkami bliźniaczymi – potwierdza to ablacja w rozdz.~\ref{chap:badania}. Wejście modelu to $640\times640$ pikseli; wyjście po dekodowaniu ma postać etykiet YOLO: \textit{class\_id cx cy w h}. Ilościową ocenę skuteczności przedstawia rozdz.~\ref{chap:badania}.


\chapter{Implementacja rozwiązania}\label{chap:implementacja}
\section{Przygotowanie zbioru danych i półautomatyczna anotacja}
Baza obejmuje $147$ gatunków i nie ma gotowych etykiet lokalizacyjnych. Ręczne rysowanie ramek dla całego zbioru jest nierealne: to kilka miesięcy pracy dla jednej osoby. GroundingDINO rozwiązał ten problem bez ręcznego etykietowania geometrii na pełnej skali: zero-shotowa lokalizacja frazą \texttt{"mushroom"} wygenerowała etykiety automatycznie, a nazwa gatunku brała się z folderu iNaturalist, nie z głowicy modelu. Jakość lokalizacji potwierdzono ręcznym audytem geometrii $1\,065$ ramek oraz spisem pominięć na $250$ kadrach izolowanego testu (podrozdz.~\ref{sec:grounding_dino}). Potok: masowa akwizycja z filtrem \texttt{quality\_grade=research}, auto-anotacja, oczyszczenie duplikatów, podział korpusu uczącego na \texttt{train}/\texttt{val} oraz niezależny zbiór testowy izolowany po identyfikatorze obserwacji \texttt{obs\_id}. Liczebności obrazów i ramek zestawia tab.~\ref{tab:podzial_zbioru}; te same wartości obowiązują w rozdz.~\ref{chap:badania}. Schemat przedstawia rys.~\ref{fig:potok_danych}.

\begin{figure}[htbp]
 \centering
 \begin{tikzpicture}[
 etap/.style={draw, rounded corners=2pt, align=center, font=\scriptsize,
 text width=2.95cm, minimum height=9mm, inner sep=2.5pt},
 src/.style={etap, fill=black!7},
 acq/.style={etap, fill=EEBlue!14},
 ann/.style={etap, fill=EEUltramarine!14},
 cln/.style={etap, fill=EETangerine!20},
 strz/.style={-latex, semithick}
 ]
 \node[src] (api) at (0,0) {iNaturalist REST API\\\texttt{quality\_grade}\\\texttt{=research}};
 \node[acq] (dl) at (3.9,0) {Skrypt \texttt{gather\_data}\\$105$ taksonów};
 \node[acq] (val) at (7.8,0) {Cel $\geq 1000$ zdjęć\\skip / retry HTTP};
 \node[acq] (purge) at (11.7,0) {Rozszerzenie\\do $147$ gatunków};

 \node[ann] (dino) at (11.7,-2.1) {GroundingDINO\\zero-shot};
 \node[ann] (reg) at (7.8,-2.1) {Rejestr klas\\\texttt{class\_registry}\\ID $0$--$146$};
 \node[cln] (nms) at (3.9,-2.1) {Czyszczenie\\duplikatów ramek};
 \node[cln] (split) at (0,-2.1) {Izolacja testu\\po \texttt{obs\_id}};

 \node[src] (yaml) at (0,-3.9) {\texttt{data.yaml}\\\texttt{nc: 147}};

 \draw[strz] (api) -- (dl);
 \draw[strz] (dl) -- (val);
 \draw[strz] (val) -- (purge);
 \draw[strz] (purge) -- (dino);
 \draw[strz] (dino) -- (reg);
 \draw[strz] (reg) -- (nms);
 \draw[strz] (nms) -- (split);
 \draw[strz] (split) -- (yaml);
 \end{tikzpicture}
 \caption{Potok przygotowania zbioru: akwizycja z iNaturalist, dociągnięcie klas do progu $1000$ zdjęć, rozszerzenie do $147$ gatunków, auto-anotacja oraz izolacja zbioru testowego po \texttt{obs\_id}. Źródło: opracowanie własne.}
 \label{fig:potok_danych}
\end{figure}

Liczebność podziałów podano w tab.~\ref{tab:podzial_zbioru}. Obrazy i ramki zliczono plik po pliku na fizycznym zbiorze (\texttt{images/} oraz \texttt{labels/} każdego podziału). Korpus uczący (\texttt{train}/\texttt{val}) pochodzi z masowej akwizycji i auto-anotacji. Skrypt weryfikacji przecięcia zliczył w korpusie uczącym $129\,607$ unikalnych \texttt{obs\_id}. Zbiór testowy zbudowano osobno: z iNaturalist pobrano nowe obserwacje, których \texttt{obs\_id} nie występuje w \texttt{train} ani \texttt{val} (skrypt \texttt{harvest\_clean\_test.py}), a ramki wygenerowano tym samym protokołem GroundingDINO. Weryfikacja przecięcia dała $0$ wspólnych \texttt{obs\_id} na $32\,727$ obrazach testowych. Cel operacyjny to do $250$ zdjęć na klasę; dla $23$ taksonów iNaturalist nie dostarczył tylu niewidzianych obserwacji research-grade, więc te klasy weszły do testu w mniejszej liczbie (od $8$ zdjęć dla \textit{Hygrophorus marzuolus} do $217$ dla \textit{Scleroderma verrucosum}). W dalszej części pracy obowiązują wartości z tej tabeli.

\begin{table}[htbp]
 \centering
 \caption{Liczba obrazów i ramek otaczających w podziałach \texttt{train}/\texttt{val}/\texttt{test} ($147$ klas). Test jest niezależnym zbiorem, izolowanym po \texttt{obs\_id}, a nie wycinkiem $10\%$ korpusu uczącego. Źródło: opracowanie własne.}
 \label{tab:podzial_zbioru}

Skrypt \texttt{split\_dataset.py} dzieli zbiór uczący i walidacyjny losowo na poziomie poszczególnych obrazów, nie grupując ich po unikalnym identyfikatorze obserwacji (\texttt{obs\_id}). Stanowi to pewne ograniczenie metodologiczne, wiążące się z ryzykiem wycieku cech (zdjęcia tego samego owocnika pod różnym kątem w \texttt{train} i \texttt{val}), co może prowadzić do niewielkiego zawyżenia metryk walidacyjnych doboru hiperparametrów. Należy jednak podkreślić, że ostateczny zbiór testowy został rygorystycznie odizolowany na poziomie \texttt{obs\_id}, co gwarantuje wiarygodność ewaluacji końcowej modelu na całkowicie niewidzianych obiektach.

 \begin{tabular}{|l|c|r|r|}
 \hline
 \textbf{Podział} & \textbf{Rola} & \textbf{Obrazy} & \textbf{Ramki} \\
 \hline \hline
 \texttt{train} & uczenie & $287\,423$ & $602\,251$ \\
 \hline
 \texttt{val} & dobór hiperparametrów & $31\,223$ & $64\,648$ \\
 \hline
 \texttt{test} & ewaluacja (izolacja \texttt{obs\_id}) & $32\,727$ & $73\,272$ \\
 \hline
 łącznie & --- & $351\,373$ & $740\,171$ \\
 \hline
 \end{tabular}
\end{table}

\paragraph{Uwaga o podziale train/val.}
Podział \texttt{train}/\texttt{val} dla klas rozszerzonych ($105$--$146$) wykonano losowym podziałem plików (\texttt{split\_dataset.py}), bez grupowania po \texttt{obs\_id}. Różne zdjęcia tej samej obserwacji mogą zatem trafić do obu zbiorów, co może zawyżać metryki walidacyjne. Izolacja po \texttt{obs\_id} obowiązuje wyłącznie dla zbioru testowego (weryfikacja: $0$ wspólnych \texttt{obs\_id}). Stanowi to ograniczenie metodologiczne obecnego podziału.

\subsection{Akwizycja danych z serwisu iNaturalist}
Zbiór uczący pochodzi z iNaturalist \cite{inaturalist} – platformy citizen science z publicznym interfejsem REST API (\url{https://api.inaturalist.org/v1/observations}). Kluczowy filtr: \texttt{quality\_grade=research}. Obserwacja na tym poziomie przeszła weryfikację przez co najmniej dwóch niezależnych identyfikatorów ze społeczności. Skrypt \texttt{gather\_data.py} odpytuje API stronicowanymi żądaniami dla każdego taksonu z listy $105$ gatunków, podmienia w adresie miniaturę \texttt{square} na wariant \texttt{large} i zapisuje plik lokalnie pod nazwą \texttt{\{obs\_id\}\_\{photo\_id\}.jpg}. Format anotacji lokalizacyjnych powstaje dopiero w kolejnym etapie (GroundingDINO); na tym kroku zbierane są same obrazy. Przykładowe ramki COCO przedstawia rys.~\ref{fig:coco_sample_web} \cite{lin2014coco}.

\begin{figure}[htbp]
 \centering
 \includegraphics[width=0.9\textwidth]{coco_dataset_sample_official}
 \caption{Przykładowe ramki BBOX ze zbioru COCO. Źródło: \cite{lin2014coco}.}
 \label{fig:coco_sample_web}
\end{figure}

Pobieranie jest sekwencyjne w ramach gatunku: jedna sesja HTTP (\texttt{requests.Session}) obsługuje wszystkie taksony, a między zapisami plików obowiązuje pauza $0{,}3$\,s, żeby nie przeciążać API. Cel operacyjny to $\geq 1000$ niepustych plików na klasę (do $20$ stron $\times$ $50$ obserwacji). Jeśli katalog gatunku już spełnia próg, skrypt pomija takson (\texttt{Skip}) i przechodzi dalej. Dzięki temu przerwany proces można wznawiać bez dublowania pracy.

Błędy sieciowe i kody $429$/$500$/$503$ nie kończą całego przebiegu. Dla każdej strony API przewidziano do trzech ponowień z pauzą $5$\,s. Uszkodzony pojedynczy link do zdjęcia jest pomijany w bloku \texttt{except} – pętla idzie dalej. Po przejściu całej listy skrypt uruchamia kolejny cykl audytu, aż wszystkie $105$ klas osiągną próg; dopiero później zbiór rozszerzono o dodatkowe taksony do finalnych $147$ klas.

Rdzeń mechanizmu akwizycji (sesja HTTP, stronicowanie, wznowienie oraz cel $1000$ zdjęć) przedstawiono na listingu \ref{lst:scraper}.

\begin{lstlisting}[language=Python,caption={Akwizycja obrazów z iNaturalist (\texttt{gather\_data.py}): sesja HTTP, wznowienie pobierania i cel $1000$ zdjęć na gatunek},label={lst:scraper}]
BASE_DATA_DIR = "E:/Kiril_Horobets/mushroom_data"
TARGET_PER_CLASS = 1000
PAGES, PER_PAGE, MAX_RETRIES = 20, 50, 3

def download_species(session, species):
 folder = f"{BASE_DATA_DIR}/{species.replace(' ', '_')}"
 os.makedirs(folder, exist_ok=True)
 existing = count_valid_files(folder) # pliki > 0 B
 if existing >= TARGET_PER_CLASS:
 print(f"Skip: {species} ({existing} photos)")
 return True

 for page in range(1, PAGES + 1):
 url = (
 "https://api.inaturalist.org/v1/observations"
 f"?taxon_name={species.replace(' ', '+')}"
 f"&has[]=photos&quality_grade=research"
 f"&per_page={PER_PAGE}&page={page}"
)
 results = []
 for retry in range(MAX_RETRIES):
 response = session.get(url, headers={"User-Agent": "Mozilla/5.0"},
 timeout=15)
 if response.status_code == 200:
 results = response.json().get("results", [])
 break
 if response.status_code in (429, 500, 503):
 time.sleep(5)
 if not results:
 break

 for obs in results:
 for photo in obs.get("photos", []):
 img_url = photo["url"].replace("square", "large")
 path = os.path.join(folder, f"{obs['id']}_{photo['id']}.jpg")
 if os.path.exists(path):
 continue
 try:
 data = session.get(img_url, timeout=10).content
 with open(path, "wb") as f:
 f.write(data)
 time.sleep(0.3)
 except Exception:
 continue
 return count_valid_files(folder) >= TARGET_PER_CLASS
\end{lstlisting}

Funkcja \texttt{download\_species} najpierw liczy niepuste pliki w katalogu gatunku. Przy $\geq 1000$ zwraca \texttt{True} bez kontaktu z API – to warunek wznowienia po przerwie. W przeciwnym razie stronicuje obserwacje (\texttt{quality\_grade=research}, \texttt{has[]=photos}). Adres miniatury \texttt{square} zamieniany jest na \texttt{large}. Istniejący plik o tej samej nazwie nie jest pobierany ponownie.

Pętla \texttt{MAX\_RETRIES} chroni przed chwilowymi błędami serwera ($429$, $500$, $503$). Zapis obrazu i pauza $0{,}3$\,s ograniczają ryzyko limitu zapytań. Główna pętla \texttt{main()} owija to w cykle audytu: dopóki którakolwiek z $105$ klas nie osiągnie progu, skrypt odczekuje $60$\,s i skanuje listę od nowa. Sesja \texttt{requests.Session} utrzymuje połączenia TCP między taksonami.

Zbiorcze zestawienie parametrów skryptu akwizycji przedstawiono w tab.~\ref{tab:scraper}.

\begin{table}[htbp]
 \centering
 \caption{Parametry skryptu \texttt{gather\_data.py}. Źródło: opracowanie własne.}
 \label{tab:scraper}
 \begin{tabular}{|p{5.3cm}|p{3.1cm}|p{4.2cm}|}
 \hline
 \textbf{Parametr} & \textbf{Wartość} & \textbf{Uzasadnienie} \\
 \hline \hline
 Interfejs HTTP & \texttt{requests.Session} & ponowne użycie połączeń TCP \\
 \hline
 Liczba taksonów (lista bazowa) & $105$ & gatunki + look-alikes \\
 \hline
 Stopień weryfikacji obserwacji & \texttt{research} & wyłącznie rekordy zweryfikowane \\
 \hline
 Rekordów na stronę / liczba stron & $50$ / $20$ & do $1000$ obserwacji na takson \\
 \hline
 Cel liczebności klasy & $\geq 1000$ plików & próg \texttt{Skip} / sukces cyklu \\
 \hline
 Ponowienia HTTP & $3\times$ ($5$\,s) & kody $429$, $500$, $503$ \\
 \hline
 Pauza między zdjęciami & $0{,}3$\,s & ograniczenie obciążenia API \\
 \hline
 Katalog wyjściowy & \texttt{E:/Kiril\_Horobets/}\newline\texttt{mushroom\_data} & lokalny dysk stacji uczelnianej \\
 \hline
 Liczba klas po rozszerzeniu & $147$ & finalny rejestr produkcyjny \\
 \hline
 \end{tabular}
\end{table}

\subsection{Otwarto-słownikowa auto-anotacja GroundingDINO i rejestr klas}\label{sec:grounding_dino}
GroundingDINO dostaje tekst, nie indeks klasy. Ontologia \texttt{CaptionOntology(\{"mushroom": "mushroom"\})} mówi modelowi: znajdź obiekt semantycznie zgodny z pojęciem \texttt{mushroom} \cite{liu2023grounding}. Gatunek (\textit{Boletus edulis}, \textit{Amanita muscaria}, cokolwiek) pochodzi z nazwy katalogu iNaturalist, nie z głowicy klasyfikacyjnej. Zero-shot rozdziela dwa zadania: lokalizację owocnika od etykiety taksonomicznej. Predykcje trafiają do formatu YOLO: pięć wartości znormalizowanych – \texttt{class\_id cx cy w h}.

Mapowanie nazw gatunków na identyfikatory numeryczne obsługuje rejestr \texttt{class\_registry.txt}. Funkcja \texttt{get\_or\_create\_id()} czyta istniejące mapowania i przy nowym taksonie przydziela kolejny inkrementowany identyfikator. Reguła \textit{append-only}. Żadnej alfabetycznej renumeracji, która unieważniłaby wcześniejsze etykiety. Rejestr końcowy: 147 gatunków, identyfikatory od 0 do 146, jedno źródło prawdy dla warstwy anotacji, pliku YAML i aplikacji mobilnej.

Ograniczenia VRAM wymuszają przetwarzanie partiami po 300 obrazów (\texttt{BATCH\_SIZE = 300}). Każda partia izoluje się w katalogach przejściowych \texttt{staging\_in} i \texttt{staging\_out} – niepełny przebieg inferencji nie zanieczyszcza zbioru wynikowego, a po przerwie demon wznawia pracę bezpiecznie. Po każdej partii \texttt{gc.collect()} zwalnia obiekty Pythona; dopiero potem \texttt{torch.cuda.empty\_cache()} oddaje bloki CUDA. Kolejność istotna. Potok ma dwustopniowy mechanizm awaryjny: gdy inferencja całej partii zawiedzie, demon przełącza się na tryb pojedynczego elementu (\textit{single-item fallback}). Uszkodzony obraz, który wywołuje wyjątek w trybie izolowanym, trwale usuwa \texttt{os.remove} – nigdy więcej nie blokuje procesu.

Po zakończeniu etykietowania skrypt regeneruje plik \texttt{data.yaml}: synchronizuje ścieżkę zbioru, liczbę klas \texttt{nc: 147} i tablicę nazw z aktualnym stanem rejestru. Konfiguracja przekazywana do Ultralytics zawsze odpowiada fizycznej strukturze katalogów \texttt{images/} i \texttt{labels/}.

Zasadniczą pętlę demona auto-anotacji, obejmującą inicjalizację ontologii, przetwarzanie partiami, mechanizm awaryjny oraz zwalnianie pamięci karty graficznej, przedstawiono na listingu \ref{lst:autolabel}.

\begin{lstlisting}[language=Python,caption={Auto-anotacja zero-shot modelem GroundingDINO z przetwarzaniem partiami i zwalnianiem pamięci VRAM (\texttt{scripts/auto\_label.py})},label={lst:autolabel}]
BATCH_SIZE = 300

base_model = GroundingDINO(ontology=CaptionOntology({"mushroom": "mushroom"}))

for folder_name in os.listdir(BASE_DIR):
 if folder_name in IGNORE_DIRS:
 continue
 src_path = os.path.join(BASE_DIR, folder_name)
 class_id = get_or_create_id(folder_name)
 files_to_process = collect_pending_images(src_path)

 for i in range(0, len(files_to_process), BATCH_SIZE):
 batch_files = files_to_process[i: i + BATCH_SIZE]
 clear_staging()

 for f in batch_files:
 shutil.copy2(os.path.join(src_path, f), os.path.join(STAGING_IN, f))

 try:
 with torch.inference_mode():
 base_model.label(input_folder=STAGING_IN, output_folder=STAGING_OUT)

 generated_labels = index_generated_labels(STAGING_OUT)
 for f in batch_files:
 save_single_label(f, src_path, class_id, generated_labels)

 except Exception as batch_error:
 # Batch failed: retry one file at a time
 print(f"[WARN] Batch failed ({batch_error}), falling back to single-item")
 for single_f in batch_files:
 clear_staging()
 try:
 shutil.copy2(os.path.join(src_path, single_f),
 os.path.join(STAGING_IN, single_f))
 with torch.inference_mode():
 base_model.label(input_folder=STAGING_IN,
 output_folder=STAGING_OUT)
 save_single_label(single_f, src_path, class_id,
 index_generated_labels(STAGING_OUT))
 except Exception as single_error:
 print(f"[ERR] Dropping broken file {folder_name}/{single_f}")
 os.remove(os.path.join(src_path, single_f))
 finally:
 clear_staging()
 gc.collect()
 if torch.cuda.is_available():
 torch.cuda.empty_cache()
\end{lstlisting}

Linia 3 inicjalizuje model nauczyciela raz. Poza pętlą po katalogach. Wielokrotne ładowanie wag do VRAM byłoby kosztowne. Ontologia \texttt{CaptionOntology} mapuje zapytanie \texttt{"mushroom"} na etykietę wynikową; model otwartego słownika lokalizuje owocniki wszystkich 147 gatunków jednym pojęciem językowym.

Linia 9 przypisuje identyfikator z rejestru \textit{append-only} na podstawie nazwy katalogu – nie z predykcji GroundingDINO. Model odpowiada wyłącznie za geometrię ramki. \texttt{collect\_pending\_images} (linia 10) pomija obrazy z istniejącymi etykietami: proces jest idempotentny i odporny na przerwania.

Pętla w linii 12 dzieli pliki na partie po 300. \texttt{clear\_staging()} (linia 14) czyści katalogi przejściowe przed każdą partią. Indeksowanie wyników w linii 23 obejmuje wyłącznie etykiety bieżącego przebiegu. \texttt{torch.inference\_mode()} (linia 20) wyłącza budowę grafu gradientów i redukuje zapotrzebowanie na pamięć.

Blok \texttt{except} (linie 27--41) realizuje fallback: wyjątek partii nie przerywa procesu. Pliki wracają w trybie izolowanym – uszkodzony obraz odpada na linii 41, $299$ pozostałych jest zachowanych. Obraz, który zawiedzie w trybie izolowanym, nigdy nie wraca do przetwarzania.

Blok \texttt{finally} (linie 42--46) działa bezwarunkowo. Kolejność wywołań ma znaczenie: \texttt{gc.collect()} zwalnia referencje Pythona do tensorów, dopiero potem \texttt{torch.cuda.empty\_cache()} oddaje bloki CUDA alokatorowi. Odwrotna kolejność nie zwalniałaby bloków wciąż referencjonowanych przez żywe obiekty, co prowadzi do fragmentacji VRAM i awarii demona przy długiej pracy.

Rozkład $147$ gatunków objętych rejestrem względem kategorii toksykologicznej, zgodny z mapowaniem zaimplementowanym w rejestrze aplikacji mobilnej oraz z wykazem zamieszczonym w załączniku, przedstawiono w tab.~\ref{tab:toksycznosc} oraz, w ujęciu graficznym, na rys.~\ref{fig:rozklad_kategorii}. Gatunki objęte agresywnym progiem ostrzegawczym $0{,}18$ stanowią łącznie $17{,}7\%$ rejestru, natomiast pozostałe $82{,}3\%$ podlega progom zachowawczym.

\begin{figure}[htbp]
 \centering
 \begin{tikzpicture}[scale=1.0]
 \def\sk{0.075} % cm na jeden gatunek
 \foreach \y/\n/\etyk/\kol in {
 0/74/{Jadalny}/{EEBlue!45},
 -0.9/47/{Niejadalny}/{black!25},
 -1.8/19/{Trujący}/{EETangerine!70},
 -2.7/7/{Śmiertelnie trujący}/{red!55}}{
 \node[font=\scriptsize, anchor=east] at (-0.15,\y) {\etyk};
 \fill[\kol, draw=black!50] (0,\y-0.28) rectangle (\n*\sk,\y+0.28);
 \node[font=\scriptsize, anchor=west] at (\n*\sk+0.12,\y) {\n};
 }
 % os
 \draw[semithick] (0,-3.15) -- (80*\sk,-3.15);
 \foreach \x in {0,20,40,60,80}{
 \draw[semithick] (\x*\sk,-3.15) -- (\x*\sk,-3.28)
 node[below, font=\scriptsize] {\x};
 }
 \node[font=\scriptsize, anchor=north] at (40*\sk,-3.72) {liczba gatunków};
 \end{tikzpicture}
 \caption{Liczebność kategorii toksykologicznych w rejestrze $147$ gatunków. Źródło: opracowanie własne.}
 \label{fig:rozklad_kategorii}
\end{figure}

\begin{table}[htbp]
 \centering
 \caption{Rozkład $147$ gatunków objętych systemem względem kategorii toksykologicznej. Źródło: opracowanie własne.}
 \label{tab:toksycznosc}
 \begin{tabular}{|l|c|c|p{5.6cm}|}
 \hline
 \textbf{Kategoria} & \textbf{Liczba} & \textbf{Udział} & \textbf{Traktowanie w polityce} \\
 \hline \hline
 Jadalny & $74$ & $50{,}3\%$ & asercja od $0{,}60$ (konsensus $4$/$5$ tylko na nakładce) \\
 \hline
 Niejadalny & $47$ & $32{,}0\%$ & status neutralny \\
 \hline
 Trujący & $19$ & $12{,}9\%$ & ostrzeżenie od $0{,}18$ \\
 \hline
 Śmiertelnie trujący & $7$ & $4{,}8\%$ & ostrzeżenie od $0{,}18$ \\
 \hline \hline
 \textbf{Razem} & $\mathbf{147}$ & $\mathbf{100{,}0\%}$ & --- \\
 \hline
 \end{tabular}
\end{table}

Do kategorii najwyższego ryzyka zaliczono siedem taksonów: \textit{Amanita phalloides}, \textit{Amanita virosa}, \textit{Amanita pantherina}, \textit{Galerina marginata}, \textit{Cortinarius rubellus}, \textit{Entoloma sinuatum} oraz \textit{Gyromitra esculenta}. Gatunki te stanowią jedynie $4{,}8\%$ rejestru, lecz odpowiadają za przeważającą część ciężkich zatruć rejestrowanych w Europie Środkowej \cite{garcia2015amanita}, co uzasadnia przypisanie im najniższego progu sygnalizacji.

Automatyczna anotacja zero-shot nie zwalnia z obowiązku kontroli jakości etykiet. Geometria ramki jest jedynym sygnałem lokalizacyjnym przekazywanym dalej do YOLO11m, więc systematyczny błąd nauczyciela przeniósłby się na cały trening. Po zakończeniu przebiegu GroundingDINO przeprowadzono więc ręczny audyt lokalizacji na próbie auto-anotowanych obrazów, etykietowanych tym samym protokołem co korpus. Skontrolowano $N=1\,065$ obrazów. Audyt ocenia jakość nauczyciela, a nie zastępuje ewaluacji detektora na bieżącym harvestcie $32\,727$ zdjęć. Próba nie była równomierna względem rejestru $147$ klas: celowo zwiększono udział taksonów o najwyższym ryzyku toksykologicznym, tak aby błąd nauczyciela na klasach krytycznych dla bezpieczeństwa nie umknął w próbie zdominowanej przez gatunki jadalne. W materiale znalazło się $175$ obrazów klas śmiertelnie trujących oraz $285$ obrazów klas trujących; łącznie $460$ kadrów ($43{,}2\%$ próby) wobec $17{,}7\%$ udziału tych kategorii w rejestrze. Pozostałe $605$ obrazów pochodziło z klas jadalnych i niejadalnych.

Każdą ramkę oceniano binarnie jako poprawną albo niepoprawną względem widocznego pokroju owocnika (habitus: kapelusz oraz -- gdy widoczny -- trzon). Ramkę uznawano za niepoprawną, gdy obrys wciągał mech, ściółkę liściową, sąsiednie drewno albo pomijał istotną część owocnika.

Audyt $1\,065$ kadrów dotyczył \emph{istniejących} ramek: oceniano geometrię tego, co GroundingDINO już narysował. Nie liczono tam owocników bez ramki. $mAP$, $P$ i $R$ YOLO11m na teście są względem etykiet auto-anotacji, nie względem pełnego ręcznego spisu. To ograniczenie ground truth. Uzupełnieniem jest stratyfikowana próba $250$ kadrów izolowanego testu (\texttt{seed}=42; $100$ toksycznych, $80$ jadalnych, $70$ niejadalnych; \texttt{sample\_fn\_audit.py}, \texttt{fn\_audit\_gui.py}): na każdym zdjęciu liczono owocniki widoczne okiem, ramki leżące na grzybie, pominięcia (FN) i ramki puste (FP). Wyniki są w \texttt{fn\_audit\_sheet.csv} (tab.~\ref{tab:audyt_fn}).

Wynik inspekcji geometrii: $1\,011$ z $1\,065$ ramek ($94{,}93\%$) oceniono jako poprawne. W $54$ przypadkach ($5{,}07\%$) lokalizacja była wadliwa. Tych $54$ obrazów nie poprawiano ręcznym rysowaniem ramek -- etykiety całego zbioru, w tym nadmiarowe obrysy, poddano następnie automatycznemu czyszczeniu \texttt{clean\_duplicate\_boxes\_v2} (kolejny podrozdział). Audyt nie zastępuje pełnej ręcznej anotacji zbioru. Zestawienie podsumowano w tab.~\ref{tab:audyt_grounding}.

\begin{table}[htbp]
 \centering
 \caption{Wyniki ręcznego audytu lokalizacji GroundingDINO na próbie $1\,065$ auto-anotowanych obrazów. Źródło: opracowanie własne na podstawie \texttt{manual\_audit\_summary.json}.}
 \label{tab:audyt_grounding}
 \begin{tabular}{|l|c|}
 \hline
 \textbf{Wielkość} & \textbf{Wartość} \\
 \hline \hline
 Obrazy poddane audytowi & $1\,065$ \\
 \hline
 Ramki poprawne & $1\,011$ \quad ($94{,}93\%$) \\
 \hline
 Ramki niepoprawne & $54$ \quad ($5{,}07\%$) \\
 \hline
 Obrazy klas śmiertelnie trujących & $175$ \\
 \hline
 Obrazy klas trujących & $285$ \\
 \hline
 Obrazy klas jadalnych i niejadalnych & $605$ \\
 \hline
 Pominięcia owocników (FN nauczyciela) & próba $250$ kadrów, tab.~\ref{tab:audyt_fn} \\
 \hline
 \end{tabular}
\end{table}

W próbie $250$ kadrów doliczono $441$ owocników. Nauczyciel pominął $34$ z nich ($7{,}71\%$; $22$ kadry, $8{,}80\%$ próby). $17$ ramek nie leżało na owocniku. $471$ ramek uznano za leżące na grzybie -- więcej niż $441$ widocznych, bo na części okazów stało kilka niemal identycznych obrysów (artefakt z następnego podrozdziału). W czterech kadrach suma ``ramka na grzybie'' + FN różniła się od liczby widocznych o więcej niż $2$; pozostawiono ocenę audytora (kępy kontra zdublowane ramki). Odsetek FN jest wyższy na klasach toksycznych ($19/200$, $9{,}50\%$) niż na jadalnych ($6/117$, $5{,}13\%$). $mAP$ YOLO11m na teście pozostaje pomiarem względem etykiet auto-anotacji: owocniki pominięte przez nauczyciela nie wchodzą do mianownika czułości detektora.

\begin{table}[htbp]
 \centering
 \caption{Ręczny spis owocników i ramek GroundingDINO na stratyfikowanej próbie izolowanego testu ($N=250$, \texttt{seed}=42). FN -- widoczny owocnik bez ramki; FP -- ramka bez owocnika. Procent FN liczony względem liczby widocznych owocników w kolumnie. Źródło: \texttt{fn\_audit\_sheet.csv}.}
 \label{tab:audyt_fn}
 \begin{tabular}{|l|c|c|c|c|}
 \hline
 \textbf{Wielkość} & \textbf{Jadalne} & \textbf{Niejadalne} & \textbf{Toksyczne} & \textbf{Razem} \\
 \hline \hline
 Obrazy & $80$ & $70$ & $100$ & $250$ \\
 \hline
 Widoczne owocniki & $117$ & $124$ & $200$ & $441$ \\
 \hline
 Ramki na owocniku & $135$ & $136$ & $200$ & $471$ \\
 \hline
 FN (owocnik bez ramki) & $6$ ($5{,}13\%$) & $9$ ($7{,}26\%$) & $19$ ($9{,}50\%$) & $34$ ($7{,}71\%$) \\
 \hline
 FP (ramka bez owocnika) & $7$ & $5$ & $5$ & $17$ \\
 \hline
 Kadry z $\ge 1$ FN & $5$ & $7$ & $10$ & $22$ \\
 \hline
 Kadry z $\ge 1$ FP & $7$ & $4$ & $5$ & $16$ \\
 \hline
 \end{tabular}
\end{table}

\paragraph{Kontaminacja wielogatunkowa.}
Obecny potok anotacji przypisuje wszystkim ramkom wykrytym na danym obrazie etykietę klasy wynikającą z~nazwy folderu iNaturalist. Jeżeli na fotografii gatunku~$A$ widoczny jest również owocnik gatunku~$B$, obie ramki otrzymają etykietę~$A$ (ang. \textit{label noise}). Zjawisko to jest częściowo ograniczone przez charakter danych iNaturalist (dominujący jeden obiekt na zdjęciu, \texttt{quality\_grade=research}), lecz jego ilościowa skala nie została zbadana. Stanowi to ograniczenie metodologiczne; potencjalnym rozwiązaniem byłoby dwuetapowe filtrowanie: wstępna detekcja ramek, a~następnie klasyfikacja każdej wyciętej ramki odrębnym modelem wizerunkowym.

\subsection{Dwukryterialne czyszczenie duplikatów ramek}
Surowe etykiety GroundingDINO zawierają charakterystyczny artefakt: kilka niemal identycznych ramek na tym samym owocniku. Klasyczny NMS z łagodnym progiem IoU zawodzi tu z konkretnego powodu – naturalne kępy opieńki (\textit{Armillaria mellea}) mają wysokie przestrzenne nakładanie sąsiednich osobników. Agresywne tłumienie usunęłoby prawdziwe detekcje w gęstych koloniach. Skrypt \texttt{clean\_duplicate\_boxes\_v2.py} rozwiązuje ten konflikt koniunkcją dwóch kryteriów. Dla pary ramek tej samej klasy w formacie YOLO ze środkami $(c_x, c_y)$ i wymiarami $(w, h)$ definiuje się odległość środków:
\begin{equation}
\mathrm{Center\_Distance}(B_i, B_j)
=
\sqrt{\bigl(c_{x,i}-c_{x,j}\bigr)^{2}+\bigl(c_{y,i}-c_{y,j}\bigr)^{2}},
\end{equation}
a także współczynnik przecięcia względem mniejszej ramki (ang. \textit{Intersection over Smallest}):
\begin{equation}
\mathrm{IoS}(B_i, B_j)
=
\frac{\lvert B_i \cap B_j \rvert}{\min\bigl(\lvert B_i \rvert,\,\lvert B_j \rvert\bigr)}.
\end{equation}
Ramkę uznawano za fantomowy duplikat wyłącznie wtedy, gdy jednocześnie spełnione były dwa warunki: wysokie pokrycie przestrzenne $\bigl(\mathrm{IoU} > 0{,}75$ albo $\mathrm{IoS} > 0{,}85\bigr)$ oraz bliskość geometrycznych środków $\bigl(\mathrm{Center\_Distance} < 0{,}06\bigr)$. Koniunkcja kryterium nakładania z kryterium bliskości środków pozwalała usuwać niemal identyczne kopie tej samej detekcji, jednocześnie zachowując sąsiadujące, lecz odrębne owocniki w naturalnych kępach. Geometryczną różnicę pomiędzy obiema konfiguracjami zilustrowano na rys.~\ref{fig:ios_geometria}.

\begin{figure}[htbp]
 \centering
 \begin{tikzpicture}[
 ramka/.style={draw, thick},
 srodek/.style={circle, fill=EETangerine, inner sep=1.3pt},
 opis/.style={font=\footnotesize, align=center}
 ]
 % Przypadek A: duplikat fantomowy (zagniezdzenie, wspolny srodek)
 \draw[ramka, EEUltramarine] (0,0) rectangle (3.0,2.4);
 \draw[ramka, EEBlue] (0.45,0.4) rectangle (2.6,2.05);
 \node[srodek] at (1.5,1.2) {};
 \node[srodek] at (1.525,1.225) {};
 \node[opis] at (1.5,-0.55) {(a) duplikat fantomowy\\$\mathrm{IoU}\approx0{,}49$, $\mathrm{IoS}\approx0{,}99$\\$\mathrm{Center\_Distance}\approx0{,}00$};
 \node[opis, EETangerine] at (1.5,-1.75) {\textbf{usunięty}};

 % Przypadek B: sasiadujace owocniki w kepie
 \draw[ramka, EEUltramarine] (6.0,0) rectangle (8.4,2.4);
 \draw[ramka, EEBlue] (7.6,0.15) rectangle (10.0,2.25);
 \node[srodek] at (7.2,1.2) {};
 \node[srodek] at (8.8,1.2) {};
 \node[opis] at (8.0,-0.55) {(b) sąsiadujące owocniki\\$\mathrm{IoU}\approx0{,}17$, $\mathrm{IoS}\approx0{,}30$\\$\mathrm{Center\_Distance}$ znaczna};
 \node[opis, EEUltramarine] at (8.0,-1.75) {\textbf{zachowane}};
 \end{tikzpicture}
 \caption{Rozróżnienie duplikatu fantomowego od sąsiadujących owocników przez koniunkcję kryteriów $\mathrm{IoS}$ i $\mathrm{Center\_Distance}$. Źródło: opracowanie własne.}
 \label{fig:ios_geometria}
\end{figure} Procedurę uruchomiono na wszystkich podziałach zbioru (\texttt{train}, \texttt{val}, \texttt{test}); wynik: $22\,327$ usuniętych fantomowych ramek w $19\,670$ plikach etykiet. Implementację obu miar i warunku koniunkcyjnego pokazuje listing \ref{lst:nms_cleaner}.

\begin{lstlisting}[language=Python,caption={Dwukryterialna eliminacja fantomowych duplikatów ramek (\texttt{scripts/utils/clean\_duplicate\_boxes\_v2.py})},label={lst:nms_cleaner}]
IOU_THRESHOLD = 0.75
IOS_THRESHOLD = 0.85
CENTER_DIST_THRESHOLD = 0.06

def calculate_box_metrics(box1, box2):
 # IoU, IoS, and center distance for two YOLO boxes [cls, cx, cy, w, h]
 b1_x1, b1_x2 = box1[1] - box1[3] / 2, box1[1] + box1[3] / 2
 b1_y1, b1_y2 = box1[2] - box1[4] / 2, box1[2] + box1[4] / 2
 b2_x1, b2_x2 = box2[1] - box2[3] / 2, box2[1] + box2[3] / 2
 b2_y1, b2_y2 = box2[2] - box2[4] / 2, box2[2] + box2[4] / 2

 center_dist = math.sqrt((box1[1] - box2[1]) ** 2 + (box1[2] - box2[2]) ** 2)

 inter_x1, inter_y1 = max(b1_x1, b2_x1), max(b1_y1, b2_y1)
 inter_x2, inter_y2 = min(b1_x2, b2_x2), min(b1_y2, b2_y2)
 inter_area = max(0, inter_x2 - inter_x1) * max(0, inter_y2 - inter_y1)

 b1_area, b2_area = box1[3] * box1[4], box2[3] * box2[4]
 min_area = min(b1_area, b2_area)
 union_area = b1_area + b2_area - inter_area + 1e-6

 iou = inter_area / union_area
 ios = inter_area / (min_area + 1e-6)
 return iou, ios, center_dist

def filter_duplicates(boxes):
 filtered_boxes = []
 removed_in_file = 0

 for box_a in boxes:
 is_duplicate = False
 for box_b in filtered_boxes:
 if box_a[0] == box_b[0]:
 iou, ios, center_dist = calculate_box_metrics(box_a, box_b)
 if center_dist < CENTER_DIST_THRESHOLD and (
 iou > IOU_THRESHOLD or ios > IOS_THRESHOLD
):
 is_duplicate = True
 removed_in_file += 1
 break
 if not is_duplicate:
 filtered_boxes.append(box_a)

 return filtered_boxes, removed_in_file
\end{lstlisting}

Linie 7--10 przeliczają reprezentację YOLO, opisującą ramkę środkiem i wymiarami $(c_x, c_y, w, h)$, na współrzędne narożnikowe wymagane do wyznaczenia pola przecięcia. Przeliczenie wykonywane jest w przestrzeni znormalizowanej $[0,1]$, dzięki czemu progi decyzyjne pozostają niezależne od rozdzielczości obrazu źródłowego – ta sama wartość $\mathrm{Center\_Distance} < 0{,}06$ oznacza to samo kryterium bliskości dla zdjęcia o boku $1024$ i $4096$ pikseli.

Linia 12 realizuje bezpośrednio wzór na odległość środków. Linie 14--16 wyznaczają pole prostokąta przecięcia, przy czym funkcja \texttt{max(0, \ldots)} zeruje wynik dla ramek rozłącznych, zapobiegając ujemnym wartościom pola. Kluczowa dla przyjętego rozwiązania jest linia 23, wyznaczająca współczynnik $\mathrm{IoS}$ jako iloraz pola przecięcia i pola \emph{mniejszej} z ramek. W przeciwieństwie do klasycznego $\mathrm{IoU}$, miara ta przyjmuje wartość bliską jedności również wtedy, gdy mała ramka jest niemal całkowicie zawarta w większej. A taka właśnie konfiguracja charakteryzuje zagnieżdżone duplikaty generowane przez model nauczyciela, dla których $\mathrm{IoU}$ pozostaje niskie z powodu istotnej różnicy pól. Stała $10^{-6}$ w mianownikach (linie 20 i 23) zabezpiecza przed dzieleniem przez zero dla ramek zdegenerowanych.

Właściwa decyzja zapada w warunku z linii 34--36. Koniunkcja kryterium bliskości środków z alternatywą dwóch miar pokrycia oznacza, że wysokie nakładanie samo w sobie nie wystarcza do usunięcia ramki – konieczne jest dodatkowo, aby oba owocniki miały praktycznie wspólny środek geometryczny. Dzięki temu sąsiadujące, silnie stykające się osobniki w kępie opieńki zostają zachowane, mimo że ich wzajemne $\mathrm{IoS}$ bywa wysokie. Porównanie ograniczone jest do ramek tej samej klasy (linia 32), a algorytm ma charakter zachłanny: każda nowa ramka zestawiana jest wyłącznie ze zbiorem już zaakceptowanym, co daje złożoność $O(n^2)$ względem liczby ramek w pliku, akceptowalną przy typowej liczbie kilku--kilkunastu anotacji na obraz.

Wyniki zastosowania procedury do poszczególnych podziałów zbioru zestawiono w tab.~\ref{tab:nms_cleanup}.

\begin{table}[htbp]
 \centering
 \caption{Rezultaty dwukryterialnego czyszczenia duplikatów ramek w podziałach zbioru. Źródło: opracowanie własne.}
 \label{tab:nms_cleanup}
 \begin{tabular}{|p{5.2cm}|c|p{4.8cm}|}
 \hline
 \textbf{Kryterium / wynik} & \textbf{Wartość} & \textbf{Rola} \\
 \hline \hline
 Próg $\mathrm{IoU}$ & $0{,}75$ & pokrycie względem sumy pól \\
 \hline
 Próg $\mathrm{IoS}$ & $0{,}85$ & pokrycie względem mniejszej ramki \\
 \hline
 Próg $\mathrm{Center\_Distance}$ & $0{,}06$ & bliskość środków geometrycznych \\
 \hline
 Warunek usunięcia & \multicolumn{2}{p{9.5cm}|}{$\bigl(\mathrm{IoU}>0{,}75 \;\lor\; \mathrm{IoS}>0{,}85\bigr) \;\land\; \mathrm{Center\_Distance}<0{,}06$} \\
 \hline \hline
 Zmodyfikowane pliki etykiet & $19\,670$ & podziały \texttt{train}, \texttt{val}, \texttt{test} \\
 \hline
 Usunięte ramki fantomowe & $22\,327$ & duplikaty tego samego owocnika \\
 \hline
 Średnio na zmodyfikowany plik & $1{,}13$ & średnia liczba usuniętych ramek \\
 \hline
 \end{tabular}
\end{table}

\section{Proces treningu i optymalizacja modelu}
Zbiór gotowy. Pora uruchomić trening. Wagi YOLO11m wychodzą ze środowiska PyTorch i muszą trafić do pliku \texttt{.tflite} nadającego się do lokalnej inferencji na Snapdragonie.

\subsection{Środowisko uczenia i parametry treningu YOLO11m}
Stacja robocza: AMD Ryzen~5~3600, $32$\,GB DDR4, NVIDIA RTX~4070 ($12$\,GB VRAM), CUDA. Stos programowy: Python $3.11+$, PyTorch $2.x$, CUDA~$12.x$, Ultralytics \cite{pytorch, ultralytics2024yolo11}. Skrypt \texttt{train.py} inicjalizuje model, wczytuje \texttt{data.yaml} i uruchamia trening na GPU bez ręcznej interwencji. Dwa etapy uczenia różnią się rozmiarem wsadu i liczbą procesów ładujących dane – powód: po rozszerzeniu głowicy do 147 klas $12$\,GB VRAM przestaje wystarczać dla \texttt{batch=16}.

Architektura bazowa: YOLO11 Medium, $20\,143\,369$ parametrów, złożoność $68{,}3$\,GFLOPs \cite{ultralytics2024yolo11}. Wariant Nano nie radzi sobie z 147 gatunkami bliźniaczymi; Large przekracza budżet RAM. Medium mieści się w obu ograniczeniach. Wagi startują z COCO (\texttt{yolo11m.pt}) – uczenie transferowe \cite{lin2014coco}. Niskopoziomowe filtry krawędzi i tekstur wyuczone na ogólnym zbiorze COCO stanowią punkt startowy; model dostosowuje wyższe warstwy i głowicę do mykologii.

Rozdzielczość wejściowa: $640\times640$ pikseli w obu etapach. Biblioteka Ultralytics domyślnie aktywuje mieszaną precyzję (AMP/FP16), co na rdzeniach Tensor Core układu RTX~4070 skraca czas epoki bez destabilizacji gradientów. Domyślny harmonogram uczenia dba o stopniowe wygaszanie współczynnika w końcowej fazie zbieżności.

Trening podzielono na dwa etapy odpowiadające przyrostowej budowie rejestru klas.

Etap podstawowy: wcześniejszy, mniejszy zestaw gatunków, \texttt{batch=16}, wczesne zatrzymanie z \texttt{patience=40} na mAP walidacyjnym. Brak poprawy w ciągu 40 epok kończy etap i zachowuje najlepsze wagi.

Etap dostrajania: pełny zestaw 147 klas, wagi z etapu podstawowego. Rozszerzona głowica klasyfikacyjna kosztuje VRAM – \texttt{batch=16} nie mieści się, więc ucięto do \texttt{batch=8} i \texttt{workers=2}. Adaptacja koncentruje się na wyższych warstwach i głowicy detekcyjnej. Czas trwania: 15 epok (\texttt{epochs=15}), bez wczesnego zatrzymania. Pełny horyzont gwarantuje powtarzalną adaptację głowicy do rozszerzonego zestawu.

Konfigurację wywołania treningu dostrajającego, wraz z rozstrzyganiem urządzenia obliczeniowego, przedstawiono na listingu \ref{lst:train}.

\begin{lstlisting}[language=Python,caption={Konfiguracja dostrajania modelu YOLO11m na $147$ klasach (\texttt{scripts/train.py})},label={lst:train}]
DEFAULT_MODEL = "weights/yolo11m.pt"

def resolve_device(requested: str) -> str:
 # Map CLI device string to Ultralytics device id
 cuda_available = torch.cuda.is_available()
 if requested == "auto":
 if cuda_available:
 print(f"[device] GPU: {torch.cuda.get_device_name(0)} -> '0'")
 return "0"
 print("[device] No CUDA, using CPU")
 return "cpu"
 if requested == "cpu":
 return "cpu"
 return "0" if requested == "cuda" else requested

def train(args) -> str:
 device = resolve_device(args.device)
 project = str(Path(args.project).resolve())

 model = YOLO(args.model)

 results = model.train(
 data=args.data, # data.yaml, nc: 147
 epochs=args.epochs, # 15 for the fine-tuning stage
 imgsz=args.imgsz, # 640 x 640
 batch=args.batch, # 8 (VRAM-bound after head expansion)
 workers=args.workers, # 2 dataloader processes
 device=device,
 project=project,
 name=args.name,
 plots=True,
 show=False,
 verbose=True,
)

 save_dir = getattr(results, "save_dir", os.path.join(project, args.name))
 return os.path.join(str(save_dir), "weights", "best.pt")
\end{lstlisting}

\texttt{resolve\_device} (linie 3--14) mapuje argument CLI na identyfikator Ultralytics. Tryb \texttt{auto} bez CUDA degraduje do CPU zamiast rzucać wyjątkiem – ten sam skrypt działa na stacji treningowej i na maszynie bez akceleratora do testów poprawności potoku.

Linia 20 ładuje wagi COCO (\texttt{yolo11m.pt}) \cite{lin2014coco}. \texttt{model.train} (linie 22--38) uruchamia uczenie transferowe. Kluczowe hiperparametry:

\texttt{epochs=15} (linia 24). Stała długość etapu dostrajania bez wczesnego zatrzymania; pełny horyzont gwarantuje powtarzalność adaptacji głowicy. \texttt{batch=8} (linia 26) – rozszerzona głowica 147 klas wyczerpuje $12$\,GB VRAM przy \texttt{batch=16}. \texttt{workers=2} (linia 27) – przy NVMe wąskim gardłem jest GPU, nie dysk.

Proces uczenia wykorzystuje domyślne, zoptymalizowane pod maską mechanizmy biblioteki Ultralytics. Automatycznie uruchamiana jest mieszana precyzja (AMP): przebieg w przód i wsteczny liczony jest w FP16 na układach Tensor Core, natomiast akumulacja gradientów następuje w FP32, co zapobiega utracie dokładności numerycznej. Wbudowany harmonogram dba również o wygaszanie współczynnika uczenia, redukując oscylacje w końcowej fazie treningu.

Zestawienie hiperparametrów obu etapów uczenia przedstawiono w tab.~\ref{tab:parametry_treningu}.

\begin{table}[htbp]
 \centering
 \caption{Hiperparametry dwuetapowego procesu uczenia modelu YOLO11m. Źródło: opracowanie własne.}
 \label{tab:parametry_treningu}
 \begin{tabular}{|p{4.5cm}|c|c|}
 \hline
 \textbf{Parametr} & \textbf{Etap podstawowy} & \textbf{Dostrajanie ($147$ klas)} \\
 \hline \hline
 Wagi inicjalizujące & \texttt{yolo11m.pt} (COCO) & \texttt{best.pt} etapu podstawowego \\
 \hline
 Rozdzielczość wejściowa & $640\times640$ & $640\times640$ \\
 \hline
 Rozmiar wsadu (\texttt{batch}) & $16$ & $8$ \\
 \hline
 Procesy wczytujące (\texttt{workers}) & $8$ & $2$ \\
 \hline
 Zamrożone warstwy (\texttt{freeze}) & – & $10$ \\
 \hline
 Liczba epok & do $200$ & $15$ (stała) \\
 \hline
 Wczesne zatrzymanie & \texttt{patience=40} & wyłączone \\
 \hline
 Precyzja obliczeń & AMP FP16 & AMP FP16 \\
 \hline
 Harmonogram uczenia & kosinusoidalny & kosinusoidalny \\
 \hline
 $\texttt{lr0}$ / $\texttt{lrf}$ & $0{,}01$ / $0{,}01$ & $0{,}01$ / $0{,}01$ \\
 \hline
 \end{tabular}
\end{table}

\subsection{Konwersja do formatu TensorFlow Lite}
Najlepszy punkt kontrolny wag (\texttt{best.pt}) musi trafić na telefon jako plik \texttt{.tflite}. Tu pojawia się przeszkoda platformowa: bezpośredni eksport \texttt{format='tflite'} w Ultralytics $8.4+$ korzysta z LiteRT, które kompiluje tylko na Linux~x86\_64 i macOS. Na stacji Windows wywołanie kończy się asercją – bez artefaktu. Rozwiązanie: dwuetapowy potok w skrypcie \texttt{scripts/export\_tflite.py}. Etap pierwszy: Ultralytics eksportuje graf do TensorFlow SavedModel. Obsługiwany na wszystkich platformach. Etap drugi: natywny \texttt{tf.lite.TFLiteConverter} kompiluje SavedModel do binarnego \texttt{.tflite} z żądaną precyzją. Żadnego przenoszenia wag na Linuksa; proces powtarzalny na stacji Windows. Oba etapy konwersji pokazuje listing \ref{lst:export}.

\begin{lstlisting}[language=Python,caption={Dwuetapowy eksport PyTorch $\rightarrow$ SavedModel $\rightarrow$ TFLite FP16 (\texttt{scripts/export\_tflite.py})},label={lst:export}]
def _run_export(weights: Path, quant: str, imgsz: int, data: str) -> Path:
 # Two-stage export: SavedModel then TFLiteConverter. Returns.tflite path.
 import tensorflow as tf

 model = YOLO(str(weights))

 export_kwargs = {"format": "saved_model", "imgsz": imgsz}
 if quant == "int8":
 export_kwargs["int8"] = True
 export_kwargs["data"] = data

 saved_model_path = Path(str(model.export(**export_kwargs)))

 converter = tf.lite.TFLiteConverter.from_saved_model(str(saved_model_path))

 if quant == "int8":
 converter.optimizations = [tf.lite.Optimize.DEFAULT]
 elif quant == "fp16":
 converter.optimizations = [tf.lite.Optimize.DEFAULT]
 converter.target_spec.supported_types = [tf.float16]
 # fp32: converter defaults

 tflite_model = converter.convert()

 tflite_path = saved_model_path.parent / f"{saved_model_path.name}_{quant}.tflite"
 tflite_path.write_bytes(tflite_model)
 return tflite_path
\end{lstlisting}

Linia 12: \texttt{model.export} z \texttt{format='saved\_model'} omija LiteRT i generuje graf TensorFlow SavedModel – obsługiwany wszędzie. \texttt{imgsz} utrwala kształt tensora wejściowego $640\times640$ statycznie: Interpreter mobilny nie może realokować buforów przy każdej klatce.

Linia 14 tworzy konwerter ze ścieżki SavedModel. Rozgałęzienie w liniach 16--21 konfiguruje precyzję. FP16 (linie 19--20) wymaga dwóch właściwości jednocześnie: \texttt{Optimize.DEFAULT} uruchamia kwantyzację wag, a \texttt{target\_spec.supported\_types = [tf.float16]} wymusza połowiczną precyzję zamiast domyślnej INT8. INT8 (linie 17--18) potrzebuje dodatkowo zbioru kalibracyjnego w \texttt{data} do wyznaczenia zakresów $[r_{\min}, r_{\max}]$. Ten wariant odrzucono z powodów opisanych wcześniej.

\texttt{converter.convert()} (linia 23) kompiluje graf do FlatBuffer, który linia 26 zapisuje na dysk. Artefakt: \texttt{best\_saved\_model\_fp16.tflite}, $38{,}58$\,MB – mieści się w budżecie $40$\,MB (tab.~\ref{tab:budzety}).

Strukturę tensora wyjściowego skompilowanego modelu przedstawiono na rys.~\ref{fig:tensor_geometria}, a odwzorowanie kanałów na wielkości fizyczne zestawiono w tab.~\ref{tab:tensor}.

\begin{figure}[htbp]
 \centering
 \begin{tikzpicture}[
 opis/.style={font=\footnotesize, align=center},
 strz/.style={-latex, semithick}
 ]
 % Kanaly (os pionowa)
 \fill[EETangerine!35] (0,0) rectangle (10.5,0.55);
 \draw (0,0) rectangle (10.5,0.55);
 \node[opis] at (5.25,0.275) {kanały $0$--$3$: współrzędne ramki $(\mathit{cx},\,\mathit{cy},\,w,\,h)$};

 \fill[EEUltramarine!22] (0,-2.35) rectangle (10.5,-0.05);
 \draw (0,-2.35) rectangle (10.5,-0.05);
 \node[opis] at (5.25,-1.2) {kanały $4$--$150$: $147$ prawdopodobieństw gatunkowych};

 % Os pozioma - punkty predykcyjne
 \draw[strz] (0,-2.75) -- (10.5,-2.75);
 \node[opis] at (5.25,-3.2) {$8400$ punktów predykcyjnych na mapach cech};

 \draw[dashed, black!45] (5.0,-2.65) -- (5.0,0.65);
 \draw[dashed, black!45] (7.5,-2.65) -- (7.5,0.65);
 \node[opis] at (2.5,0.95) {siatka $80\times80$\\($6400$)};
 \node[opis] at (6.25,0.95) {$40\times40$\\($1600$)};
 \node[opis] at (9.0,0.95) {$20\times20$\\($400$)};

 % Wymiar pionowy
 \draw[latex-latex, semithick] (-0.45,-2.35) -- (-0.45,0.55);
 \node[opis, rotate=90] at (-0.85,-0.9) {$151$ kanałów};
 \end{tikzpicture}
 \caption{Geometria tensora wyjściowego $[1,\,151,\,8400]$ z podziałem punktów predykcyjnych na trzy skale map cech (architektura \textit{anchor-free}). Źródło: opracowanie własne.}
 \label{fig:tensor_geometria}
\end{figure}

\begin{table}[htbp]
 \centering
 \caption{Geometria tensora wyjściowego modelu TFLite FP16 o kształcie $[1,\,151,\,8400]$. Źródło: opracowanie własne.}
 \label{tab:tensor}
 \begin{tabular}{|c|c|p{7.3cm}|}
 \hline
 \textbf{Wymiar} & \textbf{Rozmiar} & \textbf{Interpretacja} \\
 \hline \hline
 $0$ & $1$ & rozmiar wsadu (pojedyncza klatka) \\
 \hline
 $1$ & $151$ & $4$ współrzędne ramki $(\mathit{cx},\mathit{cy},w,h)$ oraz $147$ prawdopodobieństw gatunkowych \\
 \hline
 $2$ & $8400$ & punkty predykcyjne na mapach cech: $80^2 + 40^2 + 20^2 = 6400+1600+400$ \\
 \hline \hline
 \multicolumn{3}{|l|}{Indeksowanie w dekoderze: $\mathit{idx} = \mathit{kanal} \cdot 8400 + \mathit{punkt}$} \\
 \hline
 \multicolumn{3}{|l|}{Rozmiar bufora wyjściowego: $151 \cdot 8400 \cdot 4$\,B $\approx 4{,}84$\,MB (Float32)} \\
 \hline
 \end{tabular}
\end{table}

Pełnej kwantyzacji INT8 nie wdrożono. Artefakt produkcyjny to \texttt{best\_saved\_model\_fp16.tflite}, $38{,}58$\,MB. INT8 schodziłby do około $20$\,MB i potencjalnie szybciej na NPU, lecz stały kwant $S$ spłaszcza gradienty jasności na blaszkach, pierścieniu i łuskach kapelusza -- cechach, które odróżniają jadalną czubajkę od trującej czubajeczki oraz borowika od goryczaka. W pracy nie mierzono $mAP$ wariantu INT8; rezygnacja wynika z własności reprezentacji (rozdz.~\ref{chap:przeglad_rozwiazan}), nie z nieprzeprowadzonej ablacji. FP16 daje niemal dwukrotną redukcję wobec FP32 przy zachowaniu struktury zmiennoprzecinkowej. Decyzja produkcyjna: tylko FP16.

Tensor wyjściowy modelu ma kształt $[1,\;151,\;8400]$: $8400$ punktów predykcyjnych na trzech skalowych mapach cech ($80{\times}80$, $40{\times}40$, $20{\times}20$), $151$ kanałów. Cztery współrzędne ramki i 147 prawdopodobieństw gatunkowych. Geometria zgodna z dekoderem po stronie Androida; jedyna transformacja to mapowanie indeksów klas z rejestru \texttt{class\_registry.txt}. Po poprawnej kompilacji FP16 plik \texttt{mushrooms\_yolo11m.tflite} trafia do aktywów aplikacji mobilnej.

\section{Implementacja aplikacji mobilnej Android}
Aplikacja mobilna to natywny Kotlin na Androida. Trzy warstwy: deklaratywny interfejs w Jetpack Compose, akwizycja klatek przez CameraX i warstwa ML z interpreterem TFLite, logiką śledzenia i filtrem bezpieczeństwa. Granice między warstwami są ostre – każdą można testować niezależnie. Razem realizują jeden cel: lokalna inferencja na żywo, bez serwera, w lesie bez LTE.

\subsection{Architektura interfejsu w Jetpack Compose i CameraX}
Ekran główny to Compose: \texttt{MushroomCameraScreen} trzyma uprawnienia, podgląd i nakładkę; \texttt{DetectionOverlay} rysuje ramki osobno od kamery \cite{jetpackcompose}. CameraX dostarcza przypadki użycia z rys.~\ref{fig:camerax_web}: \texttt{Preview} zasila \texttt{PreviewView} przez \texttt{setSurfaceProvider(...)}, a \texttt{ImageAnalysis} podłącza analizator przez \texttt{setAnalyzer(...)} \cite{camerax}. W tej aplikacji liczy się właśnie \texttt{ImageAnalysis} – stamtąd idą klatki do TFLite. \texttt{ImageCapture} i \texttt{VideoCapture} zostawiono na uboczu; ciągła detekcja nie potrzebuje zapisu wideo.

\begin{figure}[htbp]
 \centering
 \includegraphics[width=0.9\textwidth]{camerax_arch_official}
 \caption{Architektura biblioteki CameraX z modułem ImageAnalysis. Źródło: \cite{camerax}.}
 \label{fig:camerax_web}
\end{figure}

\texttt{ImageAnalysis.Builder} dostał \texttt{STRATEGY\_KEEP\_ONLY\_LATEST}: w buforze zostaje tylko najnowsza klatka, reszta leci do kosza. Podgląd: \texttt{FILL\_CENTER}. Analizator pracuje na własnym wątku; przed inferencją mapa bitowa jest obracana do pionu i konwertowana do RGBA.

Telefon nie wytrzyma 30 inferencji na sekundę. W \texttt{MushroomViewModel} \texttt{Mutex.tryLock()} przepuszcza co najwyżej jedną inferencję naraz. Interwał: $250$\,ms ($4$\,FPS). Na S23 Ultra ciągła analiza bez throttlingu szybko wchodziła w thermal throttling. Nakładka zaczynała się zacinać. Gdy \texttt{isSceneLocked()} zgłasza stabilną scenę, interwał rośnie do $1{,}5$\,s (\textit{heartbeat}): nie ma sensu liczyć w kółko tego samego owocnika. Listing \ref{lst:throttling} pokazuje obie bramki.

\begin{lstlisting}[language=Kotlin,caption={Ograniczanie częstotliwości inferencji i tryb ekonomiczny (\texttt{MushroomViewModel.kt})},label={lst:throttling}]
private val inferenceMutex = Mutex()
private var lastInferenceAt = 0L

fun onFrame(upright: Bitmap) {
 val activeDetector = detector
 if (activeDetector == null) {
 upright.recycle()
 return
 }

 val now = SystemClock.elapsedRealtime()
 val interval = if (tracker.isSceneLocked()) LOCK_HEARTBEAT_MS else INFERENCE_INTERVAL_MS
 if (now - lastInferenceAt < interval) {
 upright.recycle()
 return
 }
 // Skip if a previous inference is still running (keep only the latest frame).
 if (!inferenceMutex.tryLock()) {
 upright.recycle()
 return
 }
 lastInferenceAt = now

 viewModelScope.launch(Dispatchers.Default) {
 try {
 val lb = FrameCropper.letterbox(upright, aspect, activeDetector.inputResolution)
 val detections = activeDetector.detect(lb.bitmap)
.map { it.copy(boundingBox = FrameCropper.toVisibleNormalized(it.boundingBox, lb)) }
 // Asymmetric floor: toxic @ 0.18, other classes @ 0.20.
.filter { it.confidence >= MushroomRegistry.displayThresholdFor(it.species.classId) }
 lb.bitmap.recycle()

 val tracked = tracker.update(detections, now)
 currentDetections = tracked
 _uiState.value = MushroomUiState.Scanning(tracked, tracker.isSceneLocked())
 } catch (t: Throwable) {
 Log.e(TAG, "Inference failed for frame", t)
 upright.recycle()
 } finally {
 inferenceMutex.unlock()
 }
 }
}

companion object {
 /** Inference cadence: 4 FPS (one frame every 250 ms). */
 private const val INFERENCE_INTERVAL_MS = 250L

 /** Slow heartbeat used while the scene is locked, to conserve battery. */
 private const val LOCK_HEARTBEAT_MS = 1_500L
}
\end{lstlisting}

Linia 12 realizuje adaptacyjny dobór interwału. Predykat \texttt{tracker.isSceneLocked()} zwraca wartość prawdziwą, gdy wszystkie potwierdzone tory śledzenia zostały zamrożone, co odpowiada scenie statycznej – użytkownik utrzymuje kamerę nieruchomo nad rozłożonymi owocnikami. W takiej sytuacji interwał wydłuża się z $250$\,ms do $1500$\,ms, redukując liczbę wywołań interpretera sześciokrotnie bez zauważalnej zmiany zawartości nakładki. Wykorzystano zegar \texttt{SystemClock.elapsedRealtime()} (linia 11), monotoniczny i odporny na zmiany czasu systemowego, w przeciwieństwie do \texttt{System.currentTimeMillis()}.

Linie 13--16 stanowią pierwszą, tanią bramkę: klatka nadchodząca przed upływem interwału jest natychmiast zwalniana wywołaniem \texttt{recycle()}. Jawne zwalnianie map bitowych jest tu konieczne, ponieważ CameraX dostarcza klatki z częstotliwością rzędu $30$\,FPS, a przy $4$\,FPS inferencji około $87\%$ z nich zostaje odrzuconych. Pozostawienie ich odśmiecaczowi prowadziłoby do gwałtownego przyrostu zajętości sterty i widocznych zacięć interfejsu.

Linie 18--21 realizują drugą bramkę, opartą na wzajemnym wykluczaniu. Metoda \texttt{tryLock()} obiektu \texttt{Mutex} z biblioteki \texttt{kotlinx.coroutines} jest nieblokująca: w przypadku zajętej blokady zwraca \texttt{false} zamiast wstrzymywać wątek wywołujący. Ma to zasadnicze znaczenie, ponieważ metoda \texttt{onFrame} wykonywana jest na wątku analizatora CameraX – użycie blokującego \texttt{lock()} spowodowałoby narastanie kolejki zaległych klatek i wzrost opóźnienia prezentacji. Przyjęta konstrukcja gwarantuje, że w danej chwili wykonywana jest co najwyżej jedna inferencja, a klatki nadmiarowe są porzucane, co odpowiada semantyce strategii \texttt{STRATEGY\_KEEP\_ONLY\_LATEST} zastosowanej w warstwie akwizycji.

Zwolnienie blokady umieszczono w bloku \texttt{finally} (linie 38--40), dzięki czemu wyjątek zgłoszony podczas inferencji nie prowadzi do trwałego zakleszczenia potoku. Filtr w linii 30 stosuje asymetryczny próg dekodera, omówiony szczegółowo w kolejnym podrozdziale.

Wyniki trafiają na płótno Jetpack Compose (\texttt{Canvas}) w komponencie \texttt{DetectionOverlay}. Znormalizowane współrzędne ramek (przestrzeń wejściowa modelu $640{\times}640$) mapują się na piksele podglądu w geometrii \texttt{FILL\_CENTER}: nakładka pokrywa się z obrazem kamery niezależnie od proporcji ekranu. Kolory obrysów ramek rejestr \texttt{MushroomRegistry} wiąże z klasą toksyczności \texttt{Toxicity} – zielony, szary, pomarańczowy i ciemnoczerwony rozróżniają gatunki bez zaglądania do etykiety tekstowej. Gesty dotykowe obsługuje \texttt{modifier.pointerInput} z \texttt{detectTapGestures}: po dotknięciu hit-test przestrzenny wskazuje aktywny prostokąt \texttt{Rect} i uruchamia ręczny zapis wybranej detekcji do historii znalezisk.

\subsection{Dekodowanie tensorów i dwustopniowy filtr bezpieczeństwa}
Interpreter TFLite zwraca tensor $[1,\;151,\;8400]$ w transponowanym układzie Ultralytics \cite{ultralytics2024yolo11}: $8400$ punktów predykcyjnych na mapach cech, $151$ kanałów. Cztery współrzędne ramki i $147$ prawdopodobieństw gatunkowych. Głowica YOLO11 jest \textit{anchor-free}: te $8400$ lokacji to komórki siatek $80{\times}80$, $40{\times}40$ i $20{\times}20$, a nie klasyczne \textit{anchor boxes}. Dekoder przegląda wszystkie punkty, wyznacza klasę o najwyższym wyniku i odrzuca kandydatów poniżej progu – globalnego i klasowo-zależnego – zanim powstanie jakakolwiek alokacja na stercie.

Ograniczenie obciążenia GC wymusza fast-fail: do NMS trafia co najwyżej 50 najwyżej ocenionych kandydatów (\texttt{MAX\_NMS\_INPUT=50}). Zachłanny NMS usuwa nakładające się ramki powyżej progu IoU \cite{redmon2016yolo}. Stały górny pułap wejść eliminuje patologiczny koszt kwadratowy przy gęstych scenach i stabilizuje czas przetwarzania klatki na Snapdragonie. Dekoder i NMS implementuje listing \ref{lst:decoder}.

\begin{lstlisting}[language=Kotlin,caption={Dekodowanie tensora \texttt{[1, 151, 8400]} z szybką eliminacją i tłumieniem niemaksymalnym (\texttt{MushroomDetector.kt})},label={lst:decoder}]
private fun decode(data: FloatArray): List<RawDetection> {
 val detections = ArrayList<RawDetection>()
 val globalMin = MushroomRegistry.minSafetyThreshold

 for (anchor in 0 until numAnchors) {
 var bestClass = -1
 var bestScore = 0f

 for (c in 0 until numClasses) {
 val score = data[(BOX_COORDS + c) * numAnchors + anchor]
 if (score > bestScore) {
 bestScore = score
 bestClass = c
 }
 }

 // Cheap global gate first, then the per-class safety threshold.
 if (bestClass < 0 || bestScore < globalMin) continue
 if (bestScore < MushroomRegistry.thresholdFor(bestClass)) continue

 val cx = data[0 * numAnchors + anchor]
 val cy = data[1 * numAnchors + anchor]
 val w = data[2 * numAnchors + anchor]
 val h = data[3 * numAnchors + anchor]
 val half = 0.5f

 detections.add(
 RawDetection(
 x1 = cx - w * half, y1 = cy - h * half,
 x2 = cx + w * half, y2 = cy + h * half,
 score = bestScore, classIndex = bestClass
)
)
 }
 return detections
}

private fun nonMaxSuppression(input: List<RawDetection>): List<RawDetection> {
 if (input.isEmpty()) return emptyList()

 val sorted = input.sortedByDescending { it.score }
.take(MAX_NMS_INPUT)
.toMutableList()
 val kept = ArrayList<RawDetection>()

 while (sorted.isNotEmpty()) {
 val best = sorted.removeAt(0)
 kept.add(best)
 if (kept.size >= MAX_DETECTIONS) break

 val iterator = sorted.iterator()
 while (iterator.hasNext()) {
 if (iou(best, iterator.next()) > IOU_THRESHOLD) iterator.remove()
 }
 }
 return kept
}

private const val IOU_THRESHOLD = 0.45f
private const val MAX_NMS_INPUT = 50
private const val MAX_DETECTIONS = 20
\end{lstlisting}

Pętla zewnętrzna (linia 5) przebiega kolejno wszystkie $8400$ punktów predykcyjnych na mapach cech (w kodzie indeks \texttt{anchor}), natomiast pętla wewnętrzna (linia 9) wyszukuje wśród $147$ kanałów gatunkowych wynik maksymalny. Adresowanie w linii 10 odzwierciedla transponowany układ tensora: element o indeksie kanału $c$ i punktu $a$ znajduje się pod pozycją $(4+c)\cdot 8400 + a$, ponieważ dane zapisane są kanałami, a nie punktami przestrzennymi. Bezpośrednia praca na spłaszczonej tablicy \texttt{FloatArray} zamiast na strukturach wielowymiarowych eliminuje pośrednictwo obiektów opakowujących w ścieżce wykonywanej $8400 \cdot 147 \approx 1{,}23$ mln razy na klatkę.

Kluczowa dla wydajności jest kolejność bramek w liniach 18--19. Najpierw wykonywane jest porównanie z progiem globalnym \texttt{globalMin} (równym $0{,}18$), będące pojedynczym porównaniem zmiennoprzecinkowym; dopiero kandydat, który je przejdzie, podlega kosztowniejszemu odpytaniu rejestru o próg właściwy dla swojej klasy. Odwrócenie tej kolejności oznaczałoby wykonanie odczytu z mapy dla każdego z $8400$ punktów, podczas gdy w praktyce warunek z linii 18 odrzuca zdecydowaną większość z nich. Istotne jest również, że alokacja obiektu \texttt{RawDetection} (linia 27) następuje dopiero po przejściu obu bramek. Dzięki temu klatka pozbawiona owocników nie generuje żadnej alokacji na stercie.

Metoda \texttt{nonMaxSuppression} (linie 37--56) realizuje szybką eliminację oraz właściwe tłumienie. Wywołanie \texttt{take(MAX\_NMS\_INPUT)} w linii 41 ogranicza wejście algorytmu do $50$ najwyżej ocenionych kandydatów. Ponieważ złożoność zachłannego NMS jest kwadratowa względem liczby wejść, ograniczenie to zamienia potencjalnie nieograniczony koszt na stały górny pułap rzędu $50^2/2 = 1250$ porównań, co gwarantuje przewidywalny czas przetwarzania niezależnie od zawartości sceny – także w przypadku patologicznym, gdy tekstura tła generuje setki słabych detekcji. Pętla główna (linie 45--55) w każdej iteracji przenosi kandydata o najwyższym wyniku do zbioru zaakceptowanego, a następnie usuwa z listy roboczej wszystkie ramki wykazujące z nim $\mathrm{IoU} > 0{,}45$. Zastosowanie iteratora z metodą \texttt{remove()} (linie 51--53) pozwala na usuwanie elementów w trakcie przechodzenia kolekcji bez tworzenia list pośrednich. Ogranicznik \texttt{MAX\_DETECTIONS} (linia 48) przerywa pętlę po $20$ zachowanych detekcjach, co odpowiada praktycznemu maksimum liczby owocników czytelnie prezentowanych jednocześnie na ekranie urządzenia mobilnego.

Zastosowane tłumienie niemaksymalne jest klasowo-niezależne (ang. \textit{class-agnostic}): porównanie $\mathrm{IoU}$ w linii 52 zestawia wyłącznie geometrię ramek, bez uwzględniania indeksu klasy gatunkowej. Oznacza to, że ramka gatunku $B$ o~niższym wyniku zostanie usunięta, jeżeli jej $\mathrm{IoU}$ z~ramką gatunku $A$ przekroczy próg $0{,}45$. W~typowym scenariuszu terenowym, gdzie owocniki różnych gatunków leżą obok siebie na stole lub w~koszu i~ich projekcje na obraz nie nachodzą na siebie znacząco, sytuacja ta nie występuje. Może się jednak zdarzyć w~przypadku gęstych kęp, w~których owocniki różnych gatunków częściowo przysłaniają się nawzajem. Wariant klasowo-zależny (ang. \textit{per-class NMS}), w~którym tłumienie dotyczy wyłącznie ramek tej samej klasy, wyeliminowałby to ograniczenie kosztem wyższej złożoności obliczeniowej.

Detekcje przeżywające etap NMS podlegają następnie asymetrycznej polityce bezpieczeństwa mykologicznego, scentralizowanej w metodzie \texttt{MushroomRegistry.displayStatus} oraz w buforze konsensusu klatek klasy \texttt{DetectionTracker}. Filtr ten stanowi jedyne źródło prawdy dla etykiet i kolorów prezentowanych zarówno na nakładce na żywo, jak i w historii zapisów. Zanim kandydat trafi do warstwy decyzyjnej, podlega on jeszcze absolutnemu progowi odrzucenia na poziomie dekodera: wynik poniżej $0{,}18$ dla klas \texttt{POISONOUS} i \texttt{DEADLY} lub poniżej $0{,}20$ dla pozostałych klas jest odrzucany jeszcze przed alokacją struktur wynikowych, co eliminuje fałszywe alarmy tła, takie jak tekstury biurek czy klawiatury. Powyżej tego progu obowiązują trzy reguły prezentacji:
\begin{enumerate}
 \item Wczesne ostrzeżenie toksyczne ($\ge 0{,}18$ dla klas \texttt{POISONOUS} i \texttt{DEADLY}): natychmiastowe wyświetlenie czerwonej lub ciemnoczerwonej etykiety „Trujący” / „Śmiertelnie trujący”, bez wymogu potwierdzenia w buforze konsensusu klatek.
 \item Status jadalności ($\ge 0{,}60$ dla gatunku klasy \texttt{EDIBLE}): wyświetlenie zielonej etykiety „Jadalny” wyłącznie wtedy, gdy dodatkowo spełniony jest warunek \texttt{voteConfirmed}, to znaczy dany gatunek stanowił najwyżej ocenioną predykcję w co najmniej $4$ z ostatnich $5$ klatek śledzonych w danym torze przez \texttt{DetectionTracker}.
 \item Status neutralny (pasmo pośrednie): gdy detekcja nie spełnia żadnego z powyższych warunków, to znaczy dla pewności z przedziału orientacyjnie $0{,}20$--$0{,}55$, a także dla detekcji gatunku jadalnego o pewności $\ge 0{,}60$, która nie została jeszcze potwierdzona przez bufor konsensusu klatek, wyświetlana jest szara etykieta „Obiekt nierozpoznany / Zbliż kamerę”. W tym paśmie nigdy nie prezentowano zielonego statusu jadalności, aby zredukować ryzyko błędnej asercji jadalności.
\end{enumerate}
Rozdzielenie warunku pewności od warunku czasowego bufora konsensusu wynika z założenia projektowego, że jednoklatkowy skok pewności powyżej progu $0{,}60$ bywa chwilowym artefaktem szumu inferencji. Bufor ma ograniczać liczbę fałszywie potwierdzonych statusów jadalności na strumieniu kamery; nie przeprowadzono osobnego eksperymentu A/B z buforem i bez niego na sekwencjach wideo. Dwustopniowość rozwiązania polega zatem na rozdzieleniu filtracji geometryczno-statystycznej (dekodowanie tensora, szybka eliminacja, NMS) od filtracji semantycznej (asymetryczna polityka progów oraz, niezależnie, konsensus czasowy). Pierwszy etap ogranicza złożoność obliczeniową i liczbę fałszywych detekcji przestrzennych. Drugi priorytetyzuje czułość dla klas toksycznych: ostrzeżenia prezentowane są wcześniej i bez opóźnienia buforowego, natomiast asercja jadalności wymaga jednocześnie wyższej pewności. Bufor $4$ z $5$ klatek jest mechanizmem stabilizacji czasowej nakładki, a nie składnikiem hipotezy $H_2$. Rozkład progów decyzyjnych na osi pewności, wraz z zaznaczeniem pasma neutralnego oraz asymetrii pomiędzy ścieżką ostrzeżenia a ścieżką asercji jadalności, przedstawiono na rys.~\ref{fig:progi_asymetryczne}.

\begin{figure}[htbp]
 \centering
 \begin{tikzpicture}[
 opis/.style={font=\footnotesize, align=center},
 strz/.style={-latex, semithick}
 ]
 % Os pewnosci
 \draw[strz] (0,0) -- (12.4,0);
 \node[opis] at (12.9,0) {$\mathrm{conf}$};
 \foreach \x/\l in {0/0{,}00, 2.16/0{,}18, 2.4/{}, 6.6/0{,}55, 7.2/0{,}60, 12/1{,}00}{
 \draw (\x,-0.1) -- (\x,0.1);
 }
 \node[opis] at (0,-0.4) {$0{,}00$};
 \node[opis] at (2.16,-0.4) {$0{,}18$};
 \node[opis] at (2.75,-0.85) {$0{,}20$};
 \node[opis] at (6.6,-0.4) {$0{,}55$};
 \node[opis] at (7.35,-0.85) {$0{,}60$};
 \node[opis] at (12,-0.4) {$1{,}00$};

 % Sciezka toksyczna (gora)
 \fill[black!12] (0,0.35) rectangle (2.16,1.05);
 \fill[EETangerine!45] (2.16,0.35) rectangle (12,1.05);
 \draw (0,0.35) rectangle (12,1.05);
 \node[opis] at (1.08,0.7) {odrzucone};
 \node[opis] at (7.1,0.7) {\textbf{ostrzeżenie toksyczne} – natychmiast, bez konsensusu};
 \node[opis, anchor=east] at (-0.2,0.7) {klasy\\trujące};

 % Sciezka jadalna (dol)
 \fill[black!12] (0,-2.2) rectangle (2.4,-1.5);
 \fill[black!25] (2.4,-2.2) rectangle (7.2,-1.5);
 \fill[EEBlue!35] (7.2,-2.2) rectangle (12,-1.5);
 \draw (0,-2.2) rectangle (12,-1.5);
 \node[opis] at (1.2,-1.85) {odrzucone};
 \node[opis] at (4.8,-1.85) {status neutralny (szary)};
 \node[opis] at (9.6,-1.85) {„Jadalny”\\wymaga konsensusu $4$/$5$};
 \node[opis, anchor=east] at (-0.2,-1.85) {klasy\\jadalne};
 \end{tikzpicture}
 \caption{Asymetryczna polityka progów decyzyjnych: ścieżka ostrzeżenia toksycznego (próg $0{,}18$) wobec ścieżki asercji jadalności (próg $0{,}60$ i konsensus czasowy). Źródło: opracowanie własne.}
 \label{fig:progi_asymetryczne}
\end{figure}

Scentralizowaną implementację tej polityki przedstawiono na listingu \ref{lst:safety}.

\begin{lstlisting}[language=Kotlin,caption={Asymetryczny filtr bezpieczeństwa mykologicznego (\texttt{MushroomRegistry.kt})},label={lst:safety}]
// Asymmetric thresholds: toxic warns early, edible needs a high bar + vote buffer.
const val UNSAFE_THRESHOLD = 0.18f
const val SAFE_THRESHOLD = 0.60f

// Mid band: detection exists but app asserts neither SAFE nor UNSAFE.
const val NEUTRAL_MIN_CONFIDENCE = 0.20f
const val NEUTRAL_MAX_CONFIDENCE = 0.55f

// SAFE badge: species must win 4 of the last 5 frames.
const val VOTE_REQUIRED = 4
const val VOTE_WINDOW = 5

// Decode / pre-tracker gate: toxic from 0.18, everything else from 0.20.
fun thresholdFor(classId: Int): Float =
 if (speciesMap[classId]?.toxicity?.isDangerous == true) UNSAFE_THRESHOLD
 else NEUTRAL_MIN_CONFIDENCE

// Badge policy:
// - toxic/deadly @ >= 0.18 -> immediate warning colour
// - edible @ >= 0.60 AND voteConfirmed -> green "Jadalny"
// - otherwise -> neutral gray
fun displayStatus(
 toxicity: Toxicity,
 confidence: Float,
 voteConfirmed: Boolean = true
): DisplayStatus =
 when {
 toxicity.isDangerous && confidence >= UNSAFE_THRESHOLD ->
 DisplayStatus(toxicity.labelPl, toxicity.color)

 toxicity == Toxicity.EDIBLE &&
 confidence >= SAFE_THRESHOLD &&
 voteConfirmed ->
 DisplayStatus(toxicity.labelPl, toxicity.color)

 else ->
 DisplayStatus(NEUTRAL_LABEL, NEUTRAL_COLOR)
 }
\end{lstlisting}

Linie 2--3 definiują dwa niesymetryczne progi decyzyjne hipotezy $H_2$. Ponad trzykrotna różnica pomiędzy progiem ostrzeżenia ($0{,}18$) a progiem asercji jadalności ($0{,}60$) odzwierciedla nierówny koszt obu typów błędu: fałszywe ostrzeżenie skutkuje pozostawieniem jadalnego owocnika w lesie, natomiast fałszywe potwierdzenie jadalności może prowadzić do spożycia gatunku toksycznego. Warunek \texttt{voteConfirmed} w dalszej części listingu należy do bufora konsensusu i nie jest składnikiem $H_2$.

Funkcja \texttt{thresholdFor} (linie 14--16) stanowi bramkę stosowaną jeszcze na poziomie dekodera, przed przekazaniem detekcji do trackera. Zwraca ona $0{,}18$ dla gatunków oznaczonych predykatem \texttt{isDangerous}, obejmującym kategorie \texttt{POISONOUS} oraz \texttt{DEADLY}, natomiast $0{,}20$ dla pozostałych. Dolny próg $0{,}20$ dla klas nietoksycznych pełni funkcję filtra szumu tła. Eliminuje przypadkowe aktywacje na teksturach niebędących owocnikami, zanim wygenerują one obiekt na stercie.

Właściwa polityka prezentacji zawarta jest w wyrażeniu \texttt{when} (linie 27--38), którego kolejność gałęzi ma znaczenie semantyczne. Gałąź pierwsza (linie 28--29) obsługuje gatunki niebezpieczne i jest ewaluowana jako pierwsza, wobec czego ostrzeżenie ma bezwarunkowe pierwszeństwo przed jakimkolwiek innym statusem. Nie występuje w niej odwołanie do zmiennej \texttt{voteConfirmed} – ostrzeżenie prezentowane jest natychmiast, już przy pierwszej klatce przekraczającej próg, bez oczekiwania na potwierdzenie czasowe.

Gałąź druga (linie 31--34) formułuje warunek koniunkcyjny dla statusu jadalności: wymagana jest jednocześnie przynależność do kategorii \texttt{EDIBLE}, pewność nie mniejsza niż $0{,}60$ oraz pozytywny wynik bufora konsensusu. Niespełnienie któregokolwiek z trzech warunków przenosi detekcję do gałęzi domyślnej (linie 36--37), zwracającej neutralną, szarą etykietę. Konstrukcja ta realizuje projekt bezpieczny w razie awarii (ang. \textit{fail-safe}): każdy stan nieprzewidziany jawnie, w tym detekcja o pewności z pasma $0{,}20$--$0{,}55$ oraz gatunek jadalny oczekujący na potwierdzenie konsensusu, skutkuje statusem neutralnym, nigdy zaś zielonym. Parametr \texttt{voteConfirmed} posiada domyślną wartość \texttt{true} dla wariantów wywołań, które chcą pominąć głosowanie czasowe, jednakże w głównym przepływie aplikacji, a także w ekranie historii, weryfikowana jest jego jawna wartość.

Zasadę działania okna przesuwnego, wraz z przykładowym przebiegiem prowadzącym do potwierdzenia i do jego cofnięcia, zilustrowano na rys.~\ref{fig:bufor_konsensusu}.

\begin{figure}[htbp]
 \centering
 \begin{tikzpicture}[
 klatka/.style={draw, minimum width=1.15cm, minimum height=0.72cm,
 font=\footnotesize, inner sep=1pt},
 zgodna/.style={klatka, fill=EEBlue!28},
 odstajaca/.style={klatka, fill=EETangerine!45},
 opis/.style={font=\footnotesize, align=left}
 ]
 % Naglowek osi klatek
 \foreach \i/\x in {1/0, 2/1.25, 3/2.5, 4/3.75, 5/5.0}{
 \node[font=\scriptsize] at (\x,0.75) {$k_{\i}$};
 }

 % Przebieg A: 4 z 5 -> potwierdzone
 \node[zgodna] at (0,0) {A};
 \node[zgodna] at (1.25,0) {A};
 \node[odstajaca] at (2.5,0) {B};
 \node[zgodna] at (3.75,0) {A};
 \node[zgodna] at (5.0,0) {A};
 \draw[EEUltramarine, thick, rounded corners=3pt] (-0.65,-0.5) rectangle (5.65,0.5);
 \node[opis, anchor=west, EEUltramarine] at (6.1,0) {$4$ z $5$ dla klasy A\\$\Rightarrow$ \texttt{voteConfirmed} $=$ \texttt{true}};

 % Przebieg B: 3 z 5 -> niepotwierdzone
 \node[zgodna] at (0,-1.55) {A};
 \node[odstajaca] at (1.25,-1.55) {B};
 \node[zgodna] at (2.5,-1.55) {A};
 \node[odstajaca] at (3.75,-1.55) {C};
 \node[zgodna] at (5.0,-1.55) {A};
 \draw[EETangerine, thick, rounded corners=3pt] (-0.65,-2.05) rectangle (5.65,-1.05);
 \node[opis, anchor=west, EETangerine] at (6.1,-1.55) {$3$ z $5$ dla klasy A\\$\Rightarrow$ status neutralny};
 \end{tikzpicture}
 \caption{Bufor konsensusu klatek: okno pięcioelementowe dopuszcza jedną klatkę odstającą, lecz blokuje status jadalności przy dwóch lub więcej niezgodnościach. Źródło: opracowanie własne.}
 \label{fig:bufor_konsensusu}
\end{figure}

Klasę \texttt{DetectionTracker} realizującą ten bufor prezentuje listing \ref{lst:voting}.

\begin{lstlisting}[language=Kotlin,caption={Bufor konsensusu klatek – potwierdzenie gatunku w $4$ z ostatnich $5$ klatek (\texttt{DetectionTracker.kt})},label={lst:voting}]
private class Track(val id: Int, seed: DetectionResult, nowMs: Long) {
 var species: MushroomSpecies = seed.species
 var confidence: Float = seed.confidence
 var box: NormalizedBox = seed.boundingBox
 var frozen: Boolean = false

 // Number of consecutive frames this track kept the same species class.
 var confirmCount: Int = 1

 // Rolling top-class ids for SAFE vote: need VOTE_REQUIRED hits in VOTE_WINDOW.
 private val recentClassIds = ArrayDeque<Int>(MushroomRegistry.VOTE_WINDOW)

 init {
 recentClassIds.addLast(seed.species.classId)
 }

 fun absorb(detection: DetectionResult, nowMs: Long) {
 lastSeenMs = nowMs
 confidence = detection.confidence

 confirmCount = if (detection.species.classId == species.classId) {
 confirmCount + 1
 } else {
 1
 }
 species = detection.species

 recentClassIds.addLast(detection.species.classId)
 if (recentClassIds.size > MushroomRegistry.VOTE_WINDOW) {
 recentClassIds.removeFirst()
 }

 if (nowMs - firstSeenMs >= FREEZE_MS) {
 frozen = true
 return
 }
 box = smooth(box, detection.boundingBox)
 }

 // True when current species won >= VOTE_REQUIRED of the recent window.
 fun voteConfirmed(): Boolean =
 recentClassIds.count { it == species.classId } >= MushroomRegistry.VOTE_REQUIRED

 fun toTracked(): TrackedDetection =
 TrackedDetection(id, species, confidence, box, frozen, voteConfirmed())
}
\end{lstlisting}

Strukturą nośną bufora jest kolejka dwustronna \texttt{ArrayDeque} (linia 11), zainicjalizowana pojemnością równą rozmiarowi okna. Wybór tej struktury wynika z wzorca dostępu: w każdej klatce wykonywane jest jedno dopisanie na końcu oraz co najwyżej jedno usunięcie z początku, a obie operacje realizowane są w czasie stałym. Zastosowanie zwykłej listy \texttt{ArrayList} wymuszałoby przesuwanie wszystkich elementów przy usuwaniu z czoła.

Metoda \texttt{absorb} (linie 17--38) wywoływana jest po dopasowaniu toru do detekcji z bieżącej klatki metodą IoU. Linie 28--31 realizują właściwy mechanizm okna przesuwnego: nowa etykieta klasowa dopisywana jest na końcu kolejki, a po przekroczeniu rozmiaru okna najstarszy element jest usuwany. Kolejka utrzymuje zatem stale co najwyżej pięć ostatnich rozstrzygnięć klasowych danego toru.

Metoda \texttt{voteConfirmed} (linie 41--42) zlicza w oknie wystąpienia identyfikatora klasy aktualnie przypisanej do toru i porównuje wynik z progiem $4$. Warto zwrócić uwagę na różnicę względem licznika \texttt{confirmCount} (linie 21--25), zliczającego wyłącznie serie \emph{następujących po sobie} klatek o tej samej klasie i resetowanego przy każdej zmianie etykiety. Kryterium głosowania jest łagodniejsze (dopuszcza pojedynczą klatkę odstającą wewnątrz okna) co odpowiada rzeczywistemu charakterowi zakłóceń: chwilowe migotanie predykcji pomiędzy dwoma zbliżonymi gatunkami nie powinno całkowicie unieważniać stabilnego rozpoznania, natomiast systematyczna niepewność, przejawiająca się dwiema lub więcej klatkami odstającymi w oknie pięcioelementowym, blokuje wyświetlenie zielonego statusu.

Metoda \texttt{toTracked} (linie 44--45) przekazuje wynik głosowania do warstwy prezentacji jako pole \texttt{voteConfirmed} struktury \texttt{TrackedDetection}, gdzie jest ono odczytywane przez omówioną wcześniej funkcję \texttt{displayStatus}. Rozdzielenie odpowiedzialności jest tu ścisłe: tracker odpowiada wyłącznie za stabilność czasową predykcji, natomiast rejestr – za interpretację bezpieczeństwa. Pełny cykl życia pojedynczego toru, od utworzenia po wygaśnięcie, przedstawiono na rys.~\ref{fig:cykl_toru}.

\begin{figure}[htbp]
 \centering
 \begin{tikzpicture}[
 stan/.style={draw, rounded corners=4pt, align=center, font=\scriptsize,
 text width=2.5cm, minimum height=9mm, inner sep=2.5pt},
 nowy/.style={stan, fill=black!7},
 akt/.style={stan, fill=EEBlue!14},
 pot/.style={stan, fill=EETangerine!25},
 kon/.style={stan, fill=black!12},
 p/.style={-latex, semithick, font=\scriptsize}
 ]
 \node[nowy] (n) at (0,0) {\textsc{Nowy}\\okno: $1$ klatka};
 \node[akt] (a) at (4.6,0) {\textsc{Śledzony}\\dopasowanie IoU};
 \node[pot] (p) at (9.2,0) {\textsc{Potwierdzony}\\$4$/$5$ zgodnych};
 \node[kon] (k) at (4.6,-3.4) {\textsc{Wygasły}\\brak dopasowania};

 \draw[p] (-1.7,0) -- node[above] {detekcja} (n);
 \draw[p] (n) -- node[above, align=center] {kolejna\\klatka} (a);
 \draw[p] (a) to[bend left=22] node[above, align=center] {konsensus\\osiągnięty} (p);
 \draw[p] (p) to[bend left=22] node[below, align=center] {zmiana klasy\\w oknie} (a);
 \draw[p] (a) -- node[right, align=left, xshift=2pt] {przekroczony\\limit czasu} (k);
 \draw[p] (p.south) |- node[pos=0.75, above] {utrata obiektu} (k.east);
 \end{tikzpicture}
 \caption{Cykl życia pojedynczego toru w module \texttt{DetectionTracker}: przejście do stanu potwierdzonego wymaga konsensusu $4$ z $5$ klatek i jest odwracalne. Źródło: opracowanie własne.}
 \label{fig:cykl_toru}
\end{figure}

\subsection{Lokalna baza danych Room i zapis historii znalezisk}
Historia znalezisk nie może ginąć po zamknięciu aplikacji. Persystencję realizuje Room nad SQLite \cite{androidroom}. Przepływ danych jest prosty: aplikacja pobiera DAO z \texttt{MushroomDatabase}, DAO odczytuje lub zapisuje encje, warstwa interfejsu czyta i ustawia pola. Warstwa danych, DAO i interfejs nie mają dodatkowych pośredników. Schemat ilustruje rys.~\ref{fig:room_db_web}.

\begin{figure}[htbp]
 \centering
 \includegraphics[width=0.85\textwidth]{room_db_arch_official}
 \caption{Struktura lokalnej bazy danych Room. Źródło: \cite{androidroom}.}
 \label{fig:room_db_web}
\end{figure}

Encja \texttt{MushroomFind} mapuje tabelę \texttt{mushroom\_finds}: klucz \texttt{id}, nazwa łacińska \texttt{speciesName}, etykieta toksyczności \texttt{toxicityLabel}, \texttt{confidence}, \texttt{timestamp} oraz \texttt{imagePath} do JPEG w \texttt{filesDir}. DAO to \texttt{MushroomFindDao}; repozytorium \texttt{FindsRepository} woła go z \texttt{Dispatchers.IO}, żeby dysk nie blokował Compose. Lista historii to \texttt{Flow}. Ekran odświeża się sam po każdym zapisie.

Interfejs DAO i warstwę repozytorium pokazano na listingu \ref{lst:room}.

\begin{lstlisting}[language=Kotlin,caption={Asynchroniczna warstwa persystencji Room z reaktywnym strumieniem \texttt{Flow} (\texttt{MushroomFindDao.kt}, \texttt{FindsRepository.kt})},label={lst:room}]
@Dao
interface MushroomFindDao {

 /** Insert a new find and return its generated row id. */
 @Insert
 suspend fun insert(find: MushroomFind): Long

 /** Observe all finds, newest first. */
 @Query("SELECT * FROM mushroom_finds ORDER BY timestamp DESC")
 fun observeAll(): Flow<List<MushroomFind>>

 /** Delete a find by id (does not remove the backing image file). */
 @Query("DELETE FROM mushroom_finds WHERE id =:id")
 suspend fun deleteById(id: Long)
}

class FindsRepository(context: Context) {

 private val appContext = context.applicationContext
 private val dao = MushroomDatabase.get(appContext).findDao()

 /** All saved finds, newest first, as a reactive stream. */
 val finds: Flow<List<MushroomFind>> = dao.observeAll()

 suspend fun save(
 speciesName: String,
 toxicityLabel: String,
 confidence: Float,
 voteConfirmed: Boolean,
 crop: Bitmap
): MushroomFind = withContext(Dispatchers.IO) {
 val dir = File(appContext.filesDir, FINDS_DIR).apply { mkdirs() }
 val file = File(dir, "find_${UUID.randomUUID()}.jpg")

 FileOutputStream(file).use { out ->
 crop.compress(Bitmap.CompressFormat.JPEG, JPEG_QUALITY, out)
 }

 val find = MushroomFind(
 speciesName = speciesName,
 toxicityLabel = toxicityLabel,
 confidence = confidence,
 voteConfirmed = voteConfirmed,
 timestamp = System.currentTimeMillis(),
 imagePath = file.absolutePath
)
 val id = dao.insert(find)
 find.copy(id = id)
 }
}
\end{lstlisting}

Adnotacje \texttt{@Dao}, \texttt{@Insert} oraz \texttt{@Query} (linie 1--14) przetwarzane są przez procesor adnotacji Room na etapie kompilacji, który generuje implementację interfejsu wraz z kodem wiążącym parametry zapytań. Zapytania SQL są przy tym weryfikowane składniowo i semantycznie względem schematu bazy w czasie budowy projektu, wobec czego literówka w nazwie kolumny skutkuje błędem kompilacji, a nie wyjątkiem czasu wykonania.

Istotne jest rozróżnienie sygnatur w liniach 6 i 10. Metody modyfikujące \texttt{insert} oraz \texttt{deleteById} zadeklarowano ze słowem kluczowym \texttt{suspend}, co wymusza ich wywołanie z wnętrza korutyny i uniemożliwia przypadkowe wykonanie operacji dyskowej na wątku głównym – próba takiego wywołania zostanie odrzucona przez kompilator. Metoda \texttt{observeAll} nie jest natomiast funkcją zawieszalną, lecz zwraca strumień \texttt{Flow}; Room monitoruje tabelę \texttt{mushroom\_finds} i emituje nową listę przy każdej modyfikacji jej zawartości. Warstwa interfejsu użytkownika nie musi zatem jawnie odświeżać widoku po zapisie. Ekran historii aktualizuje się jako konsekwencja propagacji strumienia.

Metoda \texttt{save} (linie 25--47) łączy dwie operacje o odmiennym charakterze: zapis pliku binarnego oraz wstawienie rekordu do bazy. Całość opakowano w \texttt{withContext(Dispatchers.IO)} (linia 30), przenosząc wykonanie do puli wątków przeznaczonej dla operacji blokujących. Katalog docelowy tworzony jest wewnątrz \texttt{context.filesDir} (linia 31), czyli w prywatnej piaskownicy aplikacji – wycinki nie trafiają do publicznej galerii urządzenia i są automatycznie usuwane wraz z odinstalowaniem aplikacji. Nazwę pliku generuje identyfikator \texttt{UUID} (linia 32), co eliminuje ryzyko kolizji przy zapisach wykonanych w tej samej milisekundzie.

Konstrukcja \texttt{FileOutputStream(file).use \{ \ldots \}} (linie 34--36) wykorzystuje funkcję rozszerzającą \texttt{use}, stanowiącą odpowiednik konstrukcji \textit{try-with-resources}: strumień jest zamykany również w przypadku zgłoszenia wyjątku podczas kompresji, co zapobiega wyciekowi deskryptorów plików. Kolejność operacji jest przy tym zamierzona. Plik obrazu zapisywany jest przed wstawieniem rekordu (linia 45), dzięki czemu w bazie nie może powstać wiersz wskazujący na nieistniejący plik. Zwracany obiekt \texttt{find.copy(id = id)} (linia 46) uzupełnia encję o wygenerowany klucz główny, ponieważ instancje danych w języku Kotlin są niemodyfikowalne.

Ręczny zapis uruchamia \texttt{FrameCropper}: aplikacja przechwytuje pełnorozdzielczą, zorientowaną pionowo klatkę i wycina region znormalizowanej ramki otaczającej. Wycinek trafia do JPEG w prywatnym katalogu piaskownicy (\texttt{context.filesDir}) – poza publiczną galerią urządzenia. Dane użytkownika pozostają izolowane i znikają automatycznie przy odinstalowaniu.

Ekran historii (\texttt{HistoryScreen}) to deklaratywna lista \texttt{LazyColumn}. Biblioteka Coil (\texttt{coil-compose}, \texttt{AsyncImage}) ładuje miniatury leniwie. Pamięć nie rośnie liniowo wraz z liczbą zapisów \cite{coil}. Status wizualny każdego elementu wyznacza \texttt{MushroomRegistry.displayStatus} na podstawie zapisanej pewności. Ważnym aspektem bezpieczeństwa jest tu uwzględnienie utrwalonej w bazie danych flagi \texttt{voteConfirmed} z encji \texttt{MushroomFind}. Dzięki temu na ekranie historii zielona etykieta „Jadalny” (wymagająca w trybie wideo konsensusu 4 z 5 klatek) pojawia się wyłącznie wtedy, gdy ów rygorystyczny warunek został faktycznie spełniony w momencie dokonywania zapisu predykcji oraz utrwalonego wyniku głosowania (\texttt{voteConfirmed}). Bufor konsensusu klatek jest strukturą czasową (retrospektywnie niedostępną), lecz stan jego weryfikacji w chwili zapisu jest persystowany w bazie danych. Dzięki temu ekran historii odtwarza wiernie status z momentu znalezienia: zielone „Jadalny” widnieje wyłącznie dla rekordów, które w chwili kliknięcia posiadały pewność $\ge 0{,}60$ i domknięty konsensus $4$ z $5$ klatek. W przeciwnym razie wyświetlana jest szara etykieta neutralna. Stały czerwony baner na górze ekranu przypomina: „Ostrzeżenie: Aplikacja ma charakter pomocniczy. Nigdy nie spożywaj grzybów bez pewności eksperta”.

\chapter{Badanie i analiza skuteczności systemu}\label{chap:badania}

\section{Metodyka testowania i miary oceny}
Do końcowej ewaluacji użyto niezależnego zbioru testowego (\texttt{images/test}, \texttt{labels/test}), złożonego z obserwacji iNaturalist, których \texttt{obs\_id} nie występuje w \texttt{train} ani \texttt{val}. Zapobiega to wyciekowi informacji ze zbioru uczącego do etapu oceny: dwa zdjęcia tego samego owocnika nie mogą jednocześnie uczestniczyć w uczeniu i w pomiarze. Zgodnie z tab.~\ref{tab:podzial_zbioru} zbiór liczy $32\,727$ obrazów oraz $73\,272$ anotowanych instancji owocników ($147$ klas; przecięcie \texttt{obs\_id} z korpusem uczącym: $0$). Wszystkie raportowane w tym rozdziale wartości $mAP$, $P$ i $R$ uzyskano wyłącznie na tym izolowanym podzbiorze.

Jakościowe dopasowanie ramki predykowanej $B_p$ do ramki referencyjnej $B_{gt}$ mierzono współczynnikiem przecięcia nad sumą (IoU), zgodnie z definicją wprowadzoną wcześniej:
\begin{equation}
\mathrm{IoU}(B_p, B_{gt}) = \frac{\lvert B_p \cap B_{gt} \rvert}{\lvert B_p \cup B_{gt} \rvert}.
\end{equation}
Interpretację geometryczną licznika i mianownika powyższego ilorazu, wraz z wpływem wartości progu na kwalifikację detekcji, przedstawiono na rys.~\ref{fig:iou_definicja}.

\begin{figure}[htbp]
 \centering
 \begin{tikzpicture}[scale=0.92]
 % – (a) przeciecie
 \begin{scope}[shift={(0,0)}]
 \fill[EEBlue!30] (0,0) rectangle (2.4,1.8);
 \fill[EETangerine!35] (1.1,0.6) rectangle (3.5,2.4);
 \fill[EEUltramarine!55] (1.1,0.6) rectangle (2.4,1.8);
 \draw[EEUltramarine, very thick] (0,0) rectangle (2.4,1.8);
 \draw[EETangerine, very thick] (1.1,0.6) rectangle (3.5,2.4);
 \node[font=\scriptsize, anchor=north] at (1.75,-0.25) {(a) $\lvert B_p \cap B_{gt}\rvert$};
 \end{scope}
 %, (b) suma
 \begin{scope}[shift={(5.0,0)}]
 \fill[black!20] (0,0) rectangle (2.4,1.8);
 \fill[black!20] (1.1,0.6) rectangle (3.5,2.4);
 \draw[EEUltramarine, very thick] (0,0) rectangle (2.4,1.8);
 \draw[EETangerine, very thick] (1.1,0.6) rectangle (3.5,2.4);
 \node[font=\scriptsize, anchor=north] at (1.75,-0.25) {(b) $\lvert B_p \cup B_{gt}\rvert$};
 \end{scope}
 % legenda
 \draw[EEUltramarine, very thick] (9.4,1.9) -- (10.0,1.9);
 \node[font=\scriptsize, anchor=west] at (10.1,1.9) {$B_{gt}$ – ramka referencyjna};
 \draw[EETangerine, very thick] (9.4,1.35) -- (10.0,1.35);
 \node[font=\scriptsize, anchor=west] at (10.1,1.35) {$B_p$. Ramka predykowana};
 \node[font=\scriptsize, anchor=west, align=left] at (9.4,0.45)
 {$\mathrm{IoU}\ge0{,}50 \Rightarrow TP$\\[1pt] $\mathrm{IoU}<0{,}50 \Rightarrow FP$ oraz $FN$};
 \end{tikzpicture}
 \caption{Geometryczna interpretacja współczynnika IoU oraz reguła kwalifikacji detekcji przy progu $0{,}50$. Źródło: opracowanie własne.}
 \label{fig:iou_definicja}
\end{figure}

Na podstawie progu IoU oraz zgodności etykiet klasowych klasyfikowano detekcje jako prawdziwie dodatnie ($TP$), fałszywie dodatnie ($FP$) lub fałszywie ujemne ($FN$). Precyzję $P$ oraz czułość (ang. \textit{recall}) $R$ zdefiniowano odpowiednio jako:
\begin{equation}
P = \frac{TP}{TP + FP}, \qquad R = \frac{TP}{TP + FN}.
\end{equation}
Średnią precyzję dla pojedynczej klasy ($AP$) obliczano jako pole pod krzywą precyzja--czułość:
\begin{equation}
AP = \int_{0}^{1} p(r)\,dr,
\end{equation}
natomiast średnią precyzję dla wielu klas ($mAP$) jako średnią arytmetyczną wartości $AP$ po $K$ gatunkach:
\begin{equation}
mAP = \frac{1}{K}\sum_{i=1}^{K} AP_i.
\end{equation}

W praktyce raportowano dwa warianty tej miary, zgodne ze standardem protokołu COCO \cite{lin2014coco}:
\begin{enumerate}
 \item $mAP_{50}$ (mAP@$0{,}50$) -- ewaluacja przy łagodnym progu przestrzennym $\mathrm{IoU}=0{,}50$, akcentująca poprawność klasyfikacji przy umiarkowanych wymaganiach lokalizacyjnych.
 \item $mAP_{50\text{--}95}$ (mAP@$0{,}50{:}0{,}95$) -- średnia z dziesięciu progów IoU od $0{,}50$ do $0{,}95$ z krokiem $0{,}05$, mierząca ścisłość lokalizacji ramek.
\end{enumerate}
Pierwsza miara jest bardziej odporna na drobne przesunięcia geometrii, druga natomiast surowiej karze niedokładne dopasowanie przestrzenne.

W kontekście identyfikacji grzybów profil ryzyka ma charakter asymetryczny. Fałszywie dodatnia detekcja ($FP$) (np. błędne wskazanie gatunku jadalnego) skutkuje co najwyżej irytacją użytkownika. Fałszywie ujemny wynik ($FN$), polegający na przeoczeniu gatunku trującego lub jego zaklasyfikowaniu jako jadalnego, niesie natomiast poważne konsekwencje zdrowotne. Z tego względu maksymalizacja czułości ($R$) dla klas toksycznych jest nadrzędnym kryterium projektowym, nawet kosztem pewnego spadku precyzji. Asymetria ta uzasadnia również polityki progów pewności wdrożone w aplikacji mobilnej, opisane w rozdziale \ref{chap:implementacja}.

\section{Analiza wyników detekcji na zbiorze testowym}
Model wytrenowany na architekturze YOLO11m poddano ewaluacji na izolowanym zbiorze testowym (tab.~\ref{tab:podzial_zbioru}). Obserwacje tego podzbioru nie brały udziału w uczeniu ani w strojeniu hiperparametrów: nazwy plików zachowują konwencję \texttt{\{obs\_id\}\_\{photo\_id\}.jpg}, a zbiór odrzucono i odbudowano tak, by żaden \texttt{obs\_id} z \texttt{train}/\texttt{val} nie wrócił do testu. Pomiarów dokonano standardowym potokiem walidacyjnym biblioteki Ultralytics, zapewniającym porównywalność wyników z protokołem raportowania mAP w detekcji obiektów \cite{ultralytics2024yolo11, lin2014coco}.

Globalne metryki na całym zbiorze testowym ($147$ gatunków) wyniosły: $mAP_{50}=0{,}8275$ ($82{,}75\%$), $mAP_{50\text{--}95}=0{,}7823$ ($78{,}23\%$), precyzja $P=0{,}8090$ ($80{,}90\%$) oraz czułość $R=0{,}7532$ ($75{,}32\%$). Są to wyniki wag PyTorch (\texttt{best.pt}) w protokole Ultralytics \texttt{val} (\texttt{split=test}). Artefakt produkcyjny na telefonie to TFLite FP16; próbę porównania $mAP$ obu formatów opisano w podrozdz.~\ref{sec:tflite_parity}. Wartość $mAP_{50}$ potwierdza, że model lokalizuje i rozpoznaje owocniki na niewidzianych obserwacjach, w zróżnicowanych warunkach oświetlenia naturalnego oraz na tle leśnym. Poza wynikiem zagregowanym po wszystkich klasach obliczono również metryki odrębnie dla dwóch podzbiorów gatunkowych zdefiniowanych kategorią toksykologiczną: gatunków jadalnych oraz gatunków trujących i śmiertelnie trujących (średnia arytmetyczna $AP$ po klasach danego podzbioru, zgodnie z definicją $mAP$). Charakterystykę zbioru testowego wraz z globalnymi metrykami zagregowanymi po wszystkich klasach przedstawiono w tab.~\ref{tab:wyniki_detekcji}, natomiast rozbicie wyników na kategorie toksykologiczne – w tab.~\ref{tab:wyniki_kategorie}. Asymetrię ryzyka konsumpcyjnego oceniono odrębnie w macierzy bezpieczeństwa (rozdz.~\ref{sec:safety_matrix}).

\begin{table}[h!]
 \centering
 \caption{Charakterystyka zbioru testowego oraz globalne metryki detekcji YOLO11m (wagi PyTorch, \texttt{split=test}). Źródło: opracowanie własne.}
 \label{tab:wyniki_detekcji}
 \begin{tabular}{|l|c|}
 \hline
 \textbf{Wielkość} & \textbf{Wartość} \\
 \hline \hline
 Liczba obrazów testowych & $32\,727$ \quad (tab.~\ref{tab:podzial_zbioru}) \\
 \hline
 Liczba anotowanych instancji & $73\,272$ \\
 \hline
 Izolacja po \texttt{obs\_id} & $0$ wspólnych z \texttt{train}/\texttt{val} \\
 \hline
 Liczba klas gatunkowych & $147$ \\
 \hline \hline
 $mAP_{50}$ & $0{,}8275$ \quad ($82{,}75\%$) \\
 \hline
 $mAP_{50\text{--}95}$ & $0{,}7823$ \quad ($78{,}23\%$) \\
 \hline
 Precyzja $P$ & $0{,}8090$ \quad ($80{,}90\%$) \\
 \hline
 Czułość $R$ & $0{,}7532$ \quad ($75{,}32\%$) \\
 \hline
 \end{tabular}
\end{table}

\begin{table}[htbp]
 \centering
 \caption{Wyniki detekcji na izolowanym zbiorze testowym w podziale na kategorie toksykologiczne. Źródło: opracowanie własne.}
 \label{tab:wyniki_kategorie}
 \begin{tabular}{|l|c|c|c|c|}
 \hline
 \textbf{Zakres ewaluacji} & $\boldsymbol{mAP_{50}}$ & $\boldsymbol{mAP_{50\text{--}95}}$ & $\boldsymbol{P}$ & $\boldsymbol{R}$ \\
 \hline \hline
 Cały zbiór testowy ($147$ klas) & $0{,}8275$ & $0{,}7823$ & $0{,}8090$ & $0{,}7532$ \\
 \hline
 Gatunki jadalne ($74$ klas) & $0{,}8156$ & $0{,}7740$ & $0{,}8032$ & $0{,}7497$ \\
 \hline
 Gatunki trujące i śmiertelnie trujące ($26$ klas) & $0{,}8133$ & $0{,}7667$ & $0{,}7878$ & $0{,}7373$ \\
 \hline
 \end{tabular}
\end{table}

Uzyskany poziom $mAP_{50}=82{,}75\%$ przy precyzji $80{,}90\%$ i czułości $75{,}32\%$ stanowi bazę operacyjną detektora na obserwacjach, których model nie widział w uczeniu. Rozbicie na podzbiory toksykologiczne nie wykazuje przewagi detektora na klasach trujących: $mAP_{50}$ obu grup jest zbliżone ($81{,}56\%$ dla jadalnych, $81{,}33\%$ dla trujących), a czułość na klasach toksycznych jest nawet nieco niższa ($R=73{,}73\%$ wobec $74{,}97\%$ dla jadalnych). $47$ klas niejadalnych nie wchodzi do tych dwóch wierszy tab.~\ref{tab:wyniki_kategorie}. Redukcja ryzyka konsumpcyjnego nie wynika z samej głowicy klasyfikacyjnej, lecz z pary progów $0{,}18$/$0{,}60$ oraz z macierzy ryzyka.

Zaobserwowano także systematyczną lukę pomiędzy $mAP_{50}$ a $mAP_{50\text{--}95}$, widoczną we wszystkich trzech zakresach ewaluacji. Niższy wynik przy surowych progach IoU ($0{,}50$--$0{,}95$) wynika z naturalnej zmienności morfologicznej kapeluszy oraz częściowego przesłonięcia owocników ściółką liściową, które utrudniają idealne dopasowanie geometryczne ramek. Różnica ta nie podważa użyteczności modelu w scenariuszu mobilnym, gdzie nadrzędne jest prawidłowe wykrycie i klasyfikacja gatunku, a nie milimetrowa precyzja obrysu.

\subsection{Porównanie PyTorch i TFLite FP16}\label{sec:tflite_parity}
Hipoteza $H_1$ mówi o modelu FP16 na telefonie, podczas gdy tablice $mAP$ powyżej pochodzą z wag \texttt{best.pt}. Na izolowanym teście wylosowano $2\,500$ obrazów (\texttt{seed}=42; $5\,568$ instancji). Dla \texttt{.pt} uruchomiono Ultralytics \texttt{val} na RTX~4070. Dla \texttt{best\_saved\_model\_fp16.tflite} nie użyto LiteRT \texttt{val}: ten dekoder nie parował tensora $[1,\;151,\;8400]$ z protokołem COCO. Zamiast tego Interpreter TFLite czyta wyjście w układzie aplikacji (\texttt{MushroomDetector.kt}: indeks ${}= \text{kanał}\times 8400+\text{punkt}$), letterbox $640$ z tłem $114$, NMS jak w \texttt{val} ($\mathrm{conf}=0{,}001$, $\mathrm{IoU}=0{,}70$, $\mathrm{max\_det}=300$), a $mAP$ liczy \texttt{ap\_per\_class} Ultralytics (\texttt{test\_tflite\_parity.py}).

Na tej samej próbie PyTorch dał $mAP_{50}=0{,}8467$, $mAP_{50\text{--}95}=0{,}8042$, $P=0{,}8139$, $R=0{,}7671$. TFLite FP16: $mAP_{50}=0{,}7914$, $mAP_{50\text{--}95}=0{,}7559$, $P=0{,}8169$, $R=0{,}7454$, po NMS $16\,917$ predykcji, średnio $177{,}2$\,ms/kadr na CPU (XNNPACK). Różnica $mAP_{50}$ wynosi $-0{,}0553$ (spadek względny $6{,}5\%$); $mAP_{50\text{--}95}$ spada o $0{,}0483$, czułość o $0{,}0217$, precyzja jest nieznacznie wyższa ($+0{,}0030$). Konwersja nie zeruje detekcji, ale \textbf{nie} jest ilościowym parytetem z \texttt{.pt}. Metryki pełnego testu w tab.~\ref{tab:wyniki_detekcji} pozostają pomiarem PyTorch; tab.~\ref{tab:tflite_parity} dotyczy wyłącznie tej próby. Czas $177{,}2$\,ms to Interpreter na CPU stacji, nie \texttt{detect()} na S23 Ultra.

\begin{table}[htbp]
 \centering
 \caption{Ta sama próba $2\,500$ obrazów izolowanego testu (\texttt{seed}=42). PyTorch: Ultralytics \texttt{val} (RTX~4070). TFLite: Interpreter CPU + dekoder Androida. Źródło: \texttt{tflite\_parity.json}.}
 \label{tab:tflite_parity}
 \begin{tabular}{|l|c|c|}
 \hline
 \textbf{Wielkość} & \textbf{PyTorch \texttt{best.pt}} & \textbf{TFLite FP16 (dekoder Android)} \\
 \hline \hline
 Obrazy / instancje & \multicolumn{2}{c|}{$2\,500$ / $5\,568$} \\
 \hline
 $mAP_{50}$ & $0{,}8467$ & $0{,}7914$ \\
 \hline
 $mAP_{50\text{--}95}$ & $0{,}8042$ & $0{,}7559$ \\
 \hline
 Precyzja $P$ & $0{,}8139$ & $0{,}8169$ \\
 \hline
 Czułość $R$ & $0{,}7671$ & $0{,}7454$ \\
 \hline
 Inferencja / obraz & GPU (RTX~4070) & $177{,}2$\,ms (CPU XNNPACK) \\
 \hline
 \end{tabular}
\end{table}

Wartości CFPR z~hipotezy $H_2$ obliczono wyłącznie na predykcjach PyTorch (\texttt{test\_h2\_baseline.py}). Analogiczna weryfikacja CFPR na artefakcie TFLite FP16 nie została przeprowadzona i~stanowi ograniczenie obecnej ewaluacji.

\section{Analiza macierzy ryzyka i asymetryczna polityka bezpieczeństwa}\label{sec:safety_matrix}
Szczególną kategorię błędów detekcji stanowią pomyłki pomiędzy gatunkami bliźniaczymi (ang. \textit{look-alikes}). Na zbiorze testowym oprócz $mAP$ liczono macierz ryzyka. Rejestr aplikacji (\texttt{MushroomRegistry}, $147$ klas) ma trzy etykiety biologiczne: jadalne (\texttt{EDIBLE}), niejadalne (\texttt{INEDIBLE}) oraz toksyczne (\texttt{POISONOUS} albo \texttt{DEADLY}). Warstwa wyświetlania (\texttt{displayStatus}) jest inna niż ta etykieta:
\begin{itemize}
 \item \textbf{asercja jadalności} (zielone „Jadalny”) -- tylko klasa \texttt{EDIBLE} i pewność $\ge 0{,}60$;
 \item \textbf{ostrzeżenie toksyczne} -- klasa toksyczna i pewność $\ge 0{,}18$;
 \item \textbf{brak asercji} (szary NEUTRAL) -- pozostałe przypadki, w tym \emph{wszystkie} predykcje niejadalne, niezależnie od pewności.
\end{itemize}
\texttt{CFPR} liczy się wyłącznie na instancjach toksycznych: $P(\text{asercja jadalności}\mid\text{gt toksyczny})$. Niejadalne nie są jadalne. Hipotezę $H_2$ rozstrzyga porównanie tej samej definicji przy parze $0{,}18$/$0{,}60$ oraz przy jednym progu $\tau=0{,}30$, na \emph{tym samym} izolowanym teście, w jednym przebiegu \texttt{predict} (\texttt{test\_h2\_baseline.py}, $32\,727$ obrazów, $73\,272$ instancji). Wiersze macierzy to etykiety rejestru: jadalne / niejadalne / toksyczne. Kolumny to \texttt{displayStatus}: asercja jadalności / brak asercji / ostrzeżenie. Starsza macierz ze zwiniętym wierszem ``nietoksyczne'' ($57\,432$ instancji, CFPR $2{,}50\%$) nie jest używana w weryfikacji $H_2$: składała jadalne z niejadalnymi i liczyła inną kolumnę ``SAFE''.

Na $15\,840$ instancjach toksycznych polityka asymetryczna dała $327$ asercje jadalności, a próg globalny $580$:
\begin{equation}
\mathrm{CFPR}_{0{,}18/0{,}60}
=
\frac{327}{15\,840}
=
0{,}0206
\quad (1{,}50\%),\qquad
\mathrm{CFPR}_{\tau=0{,}30}
=
\frac{580}{15\,840}
=
0{,}0366
\quad (3{,}11\%).
\end{equation}
Para progów obniża CFPR względem baseline ($H_2$ spełniona). $11\,385$ obiektów toksycznych ($98{,}50\%$) nie dostało zielonej asercji przy $0{,}18$/$0{,}60$. FAR na klasach jadalnych, $P(\text{ostrzeżenie}\mid\text{gt jadalny})$, wyniósł $1{,}26\%$ ($451/33\,301$) przy polityce asymetrycznej i $1{,}02\%$ ($388/33\,301$) przy $\tau=0{,}30$: niższy CFPR kupuje się niewielkim wzrostem ostrzeżeń na jadalnych. Tab.~\ref{tab:cfpr_porownanie}--\ref{tab:macierz_ryzyka_global}, rys.~\ref{fig:safety_matrix_h2}.

\begin{table}[htbp]
 \centering
 \caption{CFPR i FAR na izolowanym teście: para $0{,}18$/$0{,}60$ wobec jednego progu $\tau=0{,}30$. Źródło: opracowanie własne na podstawie \texttt{safety\_cfpr\_comparison.csv}.}
 \label{tab:cfpr_porownanie}
 \begin{tabular}{|l|c|c|c|c|}
 \hline
 \textbf{Polityka} & \textbf{CFP} & \textbf{CFPR} & \textbf{ostrzeżenia$\mid$jadalne} & \textbf{FAR jadalne} \\
 \hline \hline
 $0{,}18$ / $0{,}60$ & $327/15\,840$ & $1{,}50\%$ & $451/33\,301$ & $1{,}26\%$ \\
 \hline
 $\tau=0{,}30$ & $580/15\,840$ & $3{,}11\%$ & $388/33\,301$ & $1{,}02\%$ \\
 \hline
 \end{tabular}
\end{table}

\begin{table}[htbp]
 \centering
 \caption{Macierz ryzyka przy progach $0{,}18$/$0{,}60$ (\texttt{displayStatus}, bez bufora $4$ z $5$). Źródło: opracowanie własne na podstawie \texttt{safety\_matrix\_asymmetric.csv}.}
 \label{tab:macierz_ryzyka}
 \begin{tabular}{|l|c|c|c|c|}
 \hline
 \multirow{2}{*}{\textbf{Rejestr (gt)}} & \multicolumn{3}{c|}{\textbf{Predykcja}} & \multirow{2}{*}{\textbf{Razem}} \\
 \cline{2-4}
 & Jadalny & NEUTRAL & ostrzeżenie & \\
 \hline \hline
 jadalne & $22\,780$ & $12\,453$ & $451$ & $33\,301$ \\
 \hline
 niejadalne & $284$ & $25\,547$ & $199$ & $26\,030$ \\
 \hline
 toksyczne & $\mathbf{327}$ & $2\,388$ & $9\,020$ & $15\,840$ \\
 \hline
 \end{tabular}
\end{table}

\begin{table}[htbp]
 \centering
 \caption{Macierz ryzyka przy jednym progu $\tau=0{,}30$ (ten sam test i to samo dopasowanie IoU). Źródło: opracowanie własne na podstawie \texttt{safety\_matrix\_global.csv}.}
 \label{tab:macierz_ryzyka_global}
 \begin{tabular}{|l|c|c|c|c|}
 \hline
 \multirow{2}{*}{\textbf{Rejestr (gt)}} & \multicolumn{3}{c|}{\textbf{Predykcja}} & \multirow{2}{*}{\textbf{Razem}} \\
 \cline{2-4}
 & Jadalny & NEUTRAL & ostrzeżenie & \\
 \hline \hline
 jadalne & $27\,598$ & $7\,721$ & $388$ & $33\,301$ \\
 \hline
 niejadalne & $493$ & $25\,367$ & $170$ & $26\,030$ \\
 \hline
 toksyczne & $\mathbf{580}$ & $2\,724$ & $8\,475$ & $15\,840$ \\
 \hline
 \end{tabular}
\end{table}

\begin{figure}[htbp]
 \centering
 \begin{subfigure}[t]{0.48\textwidth}
  \centering
  \includegraphics[width=\linewidth]{wyniki/safety_matrix_asymmetric}
  \caption{Para $0{,}18$/$0{,}60$, $\mathrm{CFPR}=1{,}50\%$.}
  \label{fig:safety_matrix_asymm}
 \end{subfigure}
 \hfill
 \begin{subfigure}[t]{0.48\textwidth}
  \centering
  \includegraphics[width=\linewidth]{wyniki/safety_matrix_global}
  \caption{Próg $\tau=0{,}30$, $\mathrm{CFPR}=3{,}11\%$.}
  \label{fig:safety_matrix_glob}
 \end{subfigure}
 \caption{Macierze ryzyka na izolowanym teście. Wiersze: EDIBLE / INEDIBLE / TOXIC. Kolumny: JADALNY / NEUTRAL / OSTRZEZENIE. Źródło: opracowanie własne.}
 \label{fig:safety_matrix_h2}
\end{figure}

Uzyskany CFPR $1{,}50\%$ należy interpretować ostrożnie. $327$ z $15\,840$ obiektów toksycznych nadal dostało asercję jadalności: warstwa progów redukuje to ryzyko względem $\tau=0{,}30$ ($580$ przypadków), lecz go nie eliminuje i nie jest warunkiem wystarczającym do spożycia. Na przeciekającym podziale testowym ta sama polityka dawała $1{,}06\%$ przy innej definicji kubełków -- te liczby nie są porównywalne jeden do jednego z tab.~\ref{tab:cfpr_porownanie}. Bufor głosowania czasowego ($4$ z $5$ klatek) nie wchodzi w pomiar CFPR ani w treść $H_2$.

Dobór operacyjnych progów pewności aplikacji uzasadniono systematycznym przeciągnięciem progu detekcji (ang. \textit{confidence threshold sweep}) w zakresie od $0{,}05$ do $0{,}95$, zarejestrowanym skryptem \texttt{test3\_threshold\_sweep.py} na izolowanym zbiorze testowym, oraz niezależną ewaluacją macierzy bezpieczeństwa przy parach $(\mathrm{conf}_{\mathrm{toxic}},\,\mathrm{conf}_{\mathrm{edible}})=(0{,}18,\,0{,}60)$. Dla każdego węzła przeciągnięcia obliczano precyzję, czułość, miarę $F_1$ oraz $mAP_{50}$. Krzywa $F_1$ osiąga maksimum w okolicy niskich progów ($F_1^{\max}=0{,}7801$ przy $\mathrm{conf}=0{,}05$), jednak tak niski próg generowałby nadmierną liczbę fałszywych alarmów w warunkach terenowych.

We wcześniejszej iteracji polityki progów, poprzedzającej wprowadzenie podziału asymetrycznego, aplikacja korzystała z pojedynczego, globalnego progu operacyjnego $\mathrm{conf}=0{,}304$. Najbliższy zarejestrowany węzeł przeciągnięcia na izolowanym teście ($\mathrm{conf}=0{,}30$) daje: precyzja $P=80{,}41\%$, czułość $R=75{,}18\%$, miara $F_1=77{,}71\%$ oraz $mAP_{50}=71{,}01\%$. Pojedynczy próg globalny tego rzędu równoważy $P$ i $R$ w ujęciu zagregowanym, nie różnicuje jednak ryzyka konsumpcyjnego pomiędzy gatunkami jadalnymi a trującymi. Dlatego zastąpiono go parą $(\mathrm{conf}_{\mathrm{toxic}},\,\mathrm{conf}_{\mathrm{edible}})=(0{,}18,\,0{,}60)$. Jako docelowy punkt asercji jadalności przyjęto $\mathrm{conf}_{\mathrm{edible}}=0{,}60$, natomiast próg ostrzeżeń toksycznych ustalono na $\mathrm{conf}_{\mathrm{toxic}}=0{,}18$. Wybrane węzły przeciągnięcia zestawiono w tab.~\ref{tab:threshold_sweep}. Przebieg $F_1$ przedstawia rys.~\ref{fig:threshold_sweep_f1}.

\begin{table}[htbp]
 \centering
 \caption{Wybrane węzły przeciągnięcia progu pewności detekcji w zakresie $0{,}05$--$0{,}95$. Źródło: opracowanie własne na podstawie \texttt{threshold\_sweep.csv}.}
 \label{tab:threshold_sweep}
 \begin{tabular}{|c|c|c|c|c|p{4.4cm}|}
 \hline
 $\mathbf{conf}$ & $\boldsymbol{P}$ & $\boldsymbol{R}$ & $\boldsymbol{F_1}$ & $\boldsymbol{mAP_{50}}$ & \textbf{Uwagi} \\
 \hline \hline
 $0{,}05$ & $0{,}8090$ & $0{,}7532$ & $\mathbf{0{,}7801}$ & $0{,}7926$ & maksimum $F_1$ \\
 \hline
 $0{,}15$ & $0{,}8090$ & $0{,}7532$ & $0{,}7801$ & $0{,}7538$ & węzeł przy progu toksycznym \\
 \hline
 $0{,}20$ & $0{,}8090$ & $0{,}7532$ & $0{,}7801$ & $0{,}7390$ & dolna granica pasma neutralnego \\
 \hline
 $0{,}30$ & $0{,}8041$ & $0{,}7518$ & $0{,}7771$ & $0{,}7101$ & dawny próg globalny \\
 \hline
 $0{,}45$ & $0{,}8488$ & $0{,}6988$ & $0{,}7665$ & $0{,}6668$ & --- \\
 \hline
 $0{,}55$ & $0{,}8726$ & $0{,}6582$ & $0{,}7504$ & $0{,}6325$ & górna granica pasma neutralnego \\
 \hline
 $0{,}60$ & $0{,}8827$ & $0{,}6352$ & $0{,}7388$ & $0{,}6129$ & \textbf{próg asercji jadalności} \\
 \hline
 $0{,}75$ & $0{,}9125$ & $0{,}5448$ & $0{,}6822$ & $0{,}5326$ & --- \\
 \hline
 $0{,}90$ & $0{,}9354$ & $0{,}3604$ & $0{,}5204$ & $0{,}3570$ & załamanie czułości \\
 \hline
 \end{tabular}
\end{table}

Zestawienie ujawnia zasadniczą własność wykorzystaną przy doborze progu asercji jadalności: przejście od $\mathrm{conf}=0{,}30$ do $\mathrm{conf}=0{,}60$ podnosi precyzję z $80{,}41\%$ do $88{,}27\%$ kosztem spadku czułości z $75{,}18\%$ do $63{,}52\%$. W zastosowaniu ogólnym byłby to kompromis niekorzystny, ponieważ miara $F_1$ maleje z $0{,}7771$ do $0{,}7388$. W rozważanym zadaniu utrata czułości oznacza jednak wyłącznie brak zielonej etykiety dla części owocników rzeczywiście jadalnych (niedogodność, nie zagrożenie), podczas gdy przyrost precyzji bezpośrednio redukuje liczbę błędnych asercji jadalności. Symetryczna miara $F_1$ nie jest zatem właściwym kryterium doboru progu w warstwie decyzyjnej tego systemu.

\begin{figure}[htbp]
 \centering
 \includegraphics[width=0.85\textwidth]{wyniki/threshold_sweep_f1}
 \caption{Zależność miary $F_1$ od progu pewności detekcji w zakresie $0{,}05$--$0{,}95$; punkt asercji jadalności aplikacji $\mathrm{conf}_{\mathrm{edible}}=0{,}60$. Źródło: opracowanie własne.}
 \label{fig:threshold_sweep_f1}
\end{figure}

Skuteczność w przypadkach bliźniaczych nie sprowadza się wyłącznie do zdolności sieci do rozróżniania klas: zależy od sprzężenia detektora z polityką progów. Tab.~\ref{tab:cfpr_porownanie} pokazuje, że para $0{,}18$/$0{,}60$ obniża CFPR z $3{,}11\%$ do $1{,}50\%$ względem $\tau=0{,}30$ na tym samym teście. $H_2$ w tym ujęciu jest spełniona. Pozostałe $1{,}50\%$ przypomina, że system wspomaga identyfikację, a nie zapewnia bezpieczeństwa konsumpcyjnego.

\subsection{Pary gatunków bliźniaczych}\label{sec:lookalikes}
We wstępie wymieniono pary o wysokim koszcie pomyłki. Na izolowanym teście \texttt{test\_h2\_baseline.py} zliczał oba kierunki przy $\mathrm{IoU}\ge 0{,}50$ (tab.~\ref{tab:lookalike_pary}). Licznik nie jest CFPR: tu chodzi o konkretną parę taksonów, nie o asercję jadalności na wszystkich klasach toksycznych.

Najczęstsza pomyłka spośród tych par to \textit{Kuehneromyces mutabilis} $\to$ \textit{Galerina marginata} ($42/1\,018$, $4{,}13\%$) oraz odwrotnie ($13/808$, $1{,}61\%$). Muchomor sromotnikowy nie został pomylony z pieczarką polną ($0/780$ i $0/840$); z muchomorem cytrynowym: $7/780$ ($0{,}90\%$) i $4/305$ ($1{,}31\%$). Kurki nie sklasyfikowano jako lisówki ($0/1\,054$); dwa kadry kurki dostały etykietę \textit{Omphalotus olearius} ($2/1\,054$, $0{,}19\%$), w drugą stronę $0/482$.

\begin{table}[htbp]
 \centering
 \caption{Pomyłki pomiędzy parami bliźniaczymi na izolowanym teście. $n_{\mathrm{gt}}$ to instancje gatunku referencyjnego po dopasowaniu IoU. Źródło: opracowanie własne na podstawie \texttt{lookalike\_confusions.csv}.}
 \label{tab:lookalike_pary}
 \begin{tabular}{|p{4.2cm}|p{4.2cm}|r|r|r|}
 \hline
 \textbf{gt} & \textbf{pred} & $\boldsymbol{n_{\mathrm{gt}}}$ & \textbf{pomyłki} & \textbf{udział} \\
 \hline \hline
 \textit{A.\ phalloides} & \textit{A.\ campestris} & $780$ & $0$ & $0{,}00\%$ \\
 \hline
 \textit{A.\ campestris} & \textit{A.\ phalloides} & $840$ & $0$ & $0{,}00\%$ \\
 \hline
 \textit{A.\ phalloides} & \textit{A.\ citrina} & $780$ & $7$ & $0{,}90\%$ \\
 \hline
 \textit{A.\ citrina} & \textit{A.\ phalloides} & $305$ & $4$ & $1{,}31\%$ \\
 \hline
 \textit{A.\ virosa} & \textit{A.\ campestris} & $127$ & $1$ & $0{,}79\%$ \\
 \hline
 \textit{A.\ campestris} & \textit{A.\ virosa} & $840$ & $0$ & $0{,}00\%$ \\
 \hline
 \textit{O.\ olearius} & \textit{C.\ cibarius} & $482$ & $0$ & $0{,}00\%$ \\
 \hline
 \textit{C.\ cibarius} & \textit{O.\ olearius} & $1\,054$ & $2$ & $0{,}19\%$ \\
 \hline
 \textit{H.\ aurantiaca} & \textit{C.\ cibarius} & $464$ & $0$ & $0{,}00\%$ \\
 \hline
 \textit{C.\ cibarius} & \textit{H.\ aurantiaca} & $1\,054$ & $0$ & $0{,}00\%$ \\
 \hline
 \textit{G.\ marginata} & \textit{K.\ mutabilis} & $808$ & $13$ & $1{,}61\%$ \\
 \hline
 \textit{K.\ mutabilis} & \textit{G.\ marginata} & $1\,018$ & $42$ & $4{,}13\%$ \\
 \hline
 \end{tabular}
\end{table}

\section{Test odporności na zakłócenia leśne}
W warunkach terenowych jakość obrazu z kamery smartfonu jest systematycznie degradowana przez czynniki optyczne i środowiskowe: obniżone oświetlenie podszytu, rozmycie wynikające z ruchu ręki oraz szum elektroniczny matrycy przy podwyższonej czułości ISO. W celu ilościowej oceny odporności detektora na takie zaburzenia przeprowadzono test obciążeniowy (ang. \textit{stress test}) zaimplementowany w skrypcie \texttt{test2\_robustness\_noise.py}. Na izolowanym zbiorze testowym zasymulowano jednoczesne działanie trzech czynników: ściemnienia obrazu ze współczynnikiem jasności $0{,}6$ (redukcja o $40\%$), rozmycia ruchowego z jądrem konwolucyjnym o rozmiarze $5$ pikseli oraz addytywnego szumu Gaussa o odchyleniu standardowym $\sigma=8{,}0$, odpowiadającego typowemu szumowi matrycy przy podwyższonym ISO. Parametry symulacji oraz uzyskane wyniki zestawiono w tab.~\ref{tab:odpornosc}.

\begin{table}[htbp]
 \centering
 \caption{Parametry symulacji zakłóceń leśnych oraz degradacja skuteczności detekcji. Źródło: opracowanie własne na podstawie \texttt{robustness\_noise\_summary.json}.}
 \label{tab:odpornosc}
 \begin{tabular}{|p{5.0cm}|c|p{5.0cm}|}
 \hline
 \textbf{Czynnik zakłócający} & \textbf{Parametr} & \textbf{Odwzorowywane zjawisko} \\
 \hline \hline
 Ściemnienie obrazu & $0{,}6$ & zacienienie podszytu leśnego \\
 \hline
 Rozmycie ruchowe & jądro $5$\,px & drżenie ręki podczas akwizycji \\
 \hline
 Addytywny szum Gaussa & $\sigma=8{,}0$ & szum matrycy przy wysokim ISO \\
 \hline \hline
 \textbf{Wynik} & \textbf{Wartość} & \textbf{Komentarz} \\
 \hline \hline
 $mAP_{50}$ – obrazy czyste & $0{,}8275$ & wynik odniesienia \\
 \hline
 $mAP_{50}$. Obrazy zakłócone & $0{,}6145$ & wszystkie trzy czynniki łącznie \\
 \hline
 Degradacja bezwzględna & $21{,}30$\,pp & różnica wartości $mAP_{50}$ \\
 \hline
 Degradacja względna & $25{,}74\%$ & odniesiona do wyniku czystego \\
 \hline
 \end{tabular}
\end{table}

\begin{figure}[htbp]
 \centering
 \begin{tikzpicture}
 \def\sk{5.6} % cm na jednostke mAP
 % os pozioma
 \draw[semithick] (0,-0.55) -- (1.0*\sk,-0.55);
 \foreach \x/\e in {0/0{,}0, 0.2/0{,}2, 0.4/0{,}4, 0.6/0{,}6, 0.8/0{,}8, 1.0/1{,}0}{
 \draw[semithick] (\x*\sk,-0.55) -- (\x*\sk,-0.68) node[below, font=\scriptsize] {$\e$};
 }
 \node[font=\scriptsize, anchor=north] at (0.5*\sk,-1.1) {$mAP_{50}$};
 % slupki
 \node[font=\scriptsize, anchor=east] at (-0.15,0.85) {obrazy czyste};
 \fill[EEBlue!45, draw=black!50] (0,0.55) rectangle (0.8275*\sk,1.15);
 \node[font=\scriptsize, anchor=west] at (0.8275*\sk+0.12,0.85) {$0{,}8275$};

 \node[font=\scriptsize, anchor=east] at (-0.15,0.0) {obrazy zakłócone};
 \fill[EETangerine!60, draw=black!50] (0,-0.3) rectangle (0.6145*\sk,0.3);
 \node[font=\scriptsize, anchor=west] at (0.6145*\sk+0.12,0.0) {$0{,}6145$};

 % zaznaczenie degradacji
 \draw[latex-latex, semithick, red!70]
 (0.6145*\sk,1.55) -- node[above, font=\scriptsize, red!70] {$-21{,}30$\,pp} (0.8275*\sk,1.55);
 \draw[dashed, red!50] (0.6145*\sk,0.3) -- (0.6145*\sk,1.55);
 \draw[dashed, red!50] (0.8275*\sk,1.15) -- (0.8275*\sk,1.55);
 \end{tikzpicture}
 \caption{Degradacja skuteczności detekcji pod wpływem skumulowanych zakłóceń leśnych. Źródło: opracowanie własne na podstawie \texttt{robustness\_noise\_summary.json}.}
 \label{fig:odpornosc_wykres}
\end{figure}

Na obrazach niezakłóconych model YOLO11m osiągnął $mAP_{50}=0{,}8275$ ($82{,}75\%$). Po wprowadzeniu skumulowanych zakłóceń leśnych wartość ta spadła do $mAP_{50}=0{,}6145$ ($61{,}45\%$), co zilustrowano na rys.~\ref{fig:odpornosc_wykres}. Absolutna degradacja wyniosła zatem $21{,}30$ punktu procentowego, a degradacja względna – $25{,}74\%$ względem wyniku bazowego. Mimo istotnego spadku dokładności detektor zachował użyteczność operacyjną na poziomie $mAP_{50}>61\%$, co wskazuje, że reprezentacje cech wyuczone na zbiorze iNaturalist generalizują się do warunków o obniżonej jakości obrazu. Wynik ten uzasadnia jednocześnie konserwatywną politykę progów w aplikacji mobilnej: w warunkach leśnych część predykcji naturalnie traci pewność, a warstwa \texttt{displayStatus} zapobiega prezentowaniu niepewnych wskazań jadalności jako zielonego statusu „Jadalny”.

\section{Analiza porównawcza architektur (studium ablacyjne)}
Wybór wariantu YOLO11m zamiast lżejszego YOLO11n wymagał empirycznego uzasadnienia, ponieważ na urządzeniach mobilnych naturalną tendencją projektową jest minimalizacja liczby parametrów i rozmiaru artefaktu TFLite. Przeprowadzono zatem studium ablacyjne (ang. \textit{ablation study}), w którym obie architektury wytrenowano na tym samym zadaniu rozpoznawania $147$ gatunków, a następnie zmierzono na identycznym izolowanym teście po \texttt{obs\_id} (tab.~\ref{tab:podzial_zbioru}). Metryki w tab.~\ref{tab:ablation_yolo11} pochodzą z tego samego protokołu Ultralytics \texttt{val} (\texttt{split=test}, \texttt{imgsz=640}). Nano osiągnął $mAP_{50}=48{,}54\%$ oraz $mAP_{50\text{--}95}=42{,}82\%$; Medium -- $82{,}75\%$ / $78{,}23\%$. Różnica $34{,}21$ punktu procentowego na $mAP_{50}$ jest luką pojemności reprezentacyjnej, a nie artefaktem protokołu ewaluacji.

\begin{table}[htbp]
 \centering
 \caption{Porównanie ablacyjne YOLO11m i YOLO11n na izolowanym teście po \texttt{obs\_id} ($32\,727$ obrazów). Źródło: opracowanie własne.}
 \label{tab:ablation_yolo11}
 \begin{tabular}{|l|c|c|}
 \hline
 \textbf{Parametr} & \textbf{YOLO11m (wybrany)} & \textbf{YOLO11n (bazowy)} \\
 \hline \hline
 Liczba parametrów & $20{,}14$\,M & $2{,}66$\,M \\
 \hline
 Złożoność obliczeniowa & $68{,}3$\,GFLOPs & $6{,}7$\,GFLOPs \\
 \hline
 Rozmiar TFLite FP16 & $38{,}58$\,MB & $5{,}60$\,MB \\
 \hline
 $mAP_{50}$ & $82{,}75\%$ & $48{,}54\%$ \\
 \hline
 $mAP_{50\text{--}95}$ & $78{,}23\%$ & $42{,}82\%$ \\
 \hline
 Precyzja $P$ & $80{,}90\%$ & $51{,}49\%$ \\
 \hline
 Czułość $R$ & $75{,}32\%$ & $48{,}71\%$ \\
 \hline
 Opóźnienie inferencji (GPU) & $8{,}62$\,ms & $1{,}10$\,ms \\
 \hline
 \end{tabular}
\end{table}

\begin{figure}[htbp]
 \centering
 \begin{tikzpicture}
 \def\sk{0.058} % cm na punkt procentowy
 \def\gw{0.42} % wysokosc slupka
 % osie
 \draw[semithick] (0,-1.75) -- (100*\sk,-1.75);
 \foreach \x in {0,20,40,60,80,100}{
 \draw[semithick] (\x*\sk,-1.75) -- (\x*\sk,-1.88) node[below, font=\scriptsize] {\x};
 }
 \node[font=\scriptsize, anchor=north] at (50*\sk,-2.3) {wartość metryki [\%]};

 % mAP50
 \node[font=\scriptsize, anchor=east] at (-0.15,0.15) {$mAP_{50}$};
 \fill[EEBlue!45, draw=black!50] (0,0.15) rectangle (82.75*\sk,0.15+\gw);
 \node[font=\scriptsize, anchor=west] at (82.75*\sk+0.1,0.15+\gw/2) {$82{,}75$};
 \fill[EETangerine!60, draw=black!50] (0,0.15-\gw) rectangle (48.54*\sk,0.15);
 \node[font=\scriptsize, anchor=west] at (48.54*\sk+0.1,0.15-\gw/2) {$48{,}54$};

 % mAP50-95
 \node[font=\scriptsize, anchor=east] at (-0.15,-1.0) {$mAP_{50\text{--}95}$};
 \fill[EEBlue!45, draw=black!50] (0,-1.0) rectangle (78.23*\sk,-1.0+\gw);
 \node[font=\scriptsize, anchor=west] at (78.23*\sk+0.1,-1.0+\gw/2) {$78{,}23$};
 \fill[EETangerine!60, draw=black!50] (0,-1.0-\gw) rectangle (42.82*\sk,-1.0);
 \node[font=\scriptsize, anchor=west] at (42.82*\sk+0.1,-1.0-\gw/2) {$42{,}82$};

 % legenda
 \fill[EEBlue!45, draw=black!50] (0,1.35) rectangle (0.42,1.62);
 \node[font=\scriptsize, anchor=west] at (0.52,1.485) {YOLO11m (wybrany)};
 \fill[EETangerine!60, draw=black!50] (3.4,1.35) rectangle (3.82,1.62);
 \node[font=\scriptsize, anchor=west] at (3.92,1.485) {YOLO11n (bazowy)};
 \end{tikzpicture}
 \caption{Porównanie $mAP$ wariantów YOLO11m i YOLO11n na izolowanym teście po \texttt{obs\_id}. Źródło: opracowanie własne.}
 \label{fig:ablacja_wykres}
\end{figure}

Wariant Nano, mimo ponad siedmiokrotnie mniejszej liczby parametrów i prawie siedmiokrotnie mniejszego artefaktu FP16, osiągnął $48{,}54\%$ $mAP_{50}$ na tym samym izolowanym teście (rys.~\ref{fig:ablacja_wykres}). YOLO11m ma $82{,}75\%$. Precyzja i czułość Nano ($P=51{,}49\%$, $R=48{,}71\%$) pozostają wyraźnie poniżej Medium ($P=80{,}90\%$, $R=75{,}32\%$). Nano nie modeluje 147 gatunków bliźniaczych z akceptowalną rozdzielczością semantyczną (blaszki, pierścień, łuski, barwa trzonu). Opóźnienie GPU wariantu Medium ($8{,}62$\,ms na RTX~4070) nie jest pomiarem telefonu; budżet aplikacji to $250$\,ms na S23 Ultra (tab.~\ref{tab:latency_s23}). Zysk prędkości Nano na GPU nie rekompensuje utraty jakości. Dlatego YOLO11m przyjęto jako wariant produkcyjny.

\section{Profilowanie opóźnień potoku inferencji}
W celu rozdzielenia kosztów obliczeniowych poszczególnych etapów potoku przeprowadzono profilowanie na akceleratorze NVIDIA GeForce RTX~4070 przy wejściu $640{\times}640$ pikseli. Pomiar obejmował $100$ przebiegów i rozróżniał trzy fazy: wstępne przetwarzanie obrazu, inferencję sieci oraz postprocessing z tłumieniem niemaksymalnym (NMS). Średni czas preprocessingu wyniósł $1{,}64$\,ms (mediana $p_{50}=1{,}68$\,ms, percentyl $p_{95}=1{,}76$\,ms), średni czas inferencji $8{,}62$\,ms ($p_{50}=8{,}62$\,ms, $p_{95}=8{,}75$\,ms), a średni czas NMS $0{,}50$\,ms ($p_{50}=0{,}50$\,ms, $p_{95}=0{,}52$\,ms). Pełny rozkład statystyczny czasów poszczególnych faz zestawiono w tab.~\ref{tab:latency}.

\begin{table}[htbp]
 \centering
 \caption{Profilowanie opóźnień potoku inferencji na akceleratorze NVIDIA GeForce RTX~4070 ($100$ przebiegów, wejście $640\times640$). Wszystkie wartości w milisekundach. Źródło: opracowanie własne na podstawie \texttt{latency\_breakdown.json}.}
 \label{tab:latency}
 \begin{tabular}{|l|c|c|c|c|c|}
 \hline
 \textbf{Faza potoku} & \textbf{Średnia} & $\boldsymbol{\sigma}$ & $\boldsymbol{p_{50}}$ & $\boldsymbol{p_{95}}$ & \textbf{Maksimum} \\
 \hline \hline
 Wstępne przetwarzanie & $1{,}64$ & $0{,}08$ & $1{,}68$ & $1{,}76$ & $1{,}80$ \\
 \hline
 Inferencja sieci & $8{,}62$ & $0{,}08$ & $8{,}62$ & $8{,}75$ & $8{,}96$ \\
 \hline
 Postprocessing (NMS) & $0{,}50$ & $0{,}01$ & $0{,}50$ & $0{,}52$ & $0{,}54$ \\
 \hline \hline
 \textbf{Całkowity czas e2e} & $\mathbf{10{,}94}$ & $0{,}16$ & $10{,}98$ & $11{,}15$ & $11{,}45$ \\
 \hline \hline
 \multicolumn{6}{|l|}{Przepustowość równoważna: $1000/10{,}94 = 91{,}4$\,FPS} \\
 \hline
 \multicolumn{6}{|l|}{Wynik na karcie PC; pomiar na S23 Ultra: tab.~\ref{tab:latency_s23}} \\
 \hline
 \end{tabular}
\end{table}

\begin{figure}[htbp]
 \centering
 \begin{tikzpicture}
 \def\sk{0.70} % cm na milisekunde
 \def\gw{0.62}
 % skladowe czasu sredniego
 \fill[EEBlue!35, draw=black!50] (0,0) rectangle (1.64*\sk,\gw);
 \fill[EEUltramarine!45, draw=black!50] (1.64*\sk,0) rectangle (10.26*\sk,\gw);
 \fill[EETangerine!55, draw=black!50] (10.26*\sk,0) rectangle (10.76*\sk,\gw);
 \node[font=\scriptsize, anchor=east] at (-0.15,\gw/2) {czas średni};

 % etykiety skladowych
 \node[font=\scriptsize] at (0.82*\sk,\gw+0.32) {$1{,}64$};
 \node[font=\scriptsize] at (5.95*\sk,\gw+0.32) {$8{,}62$};
 \node[font=\scriptsize] at (10.51*\sk,\gw+0.72) {$0{,}50$};
 \draw[gray, thin] (10.51*\sk,\gw+0.58) -- (10.51*\sk,\gw+0.08);

 % os
 \draw[semithick] (0,-0.3) -- (12*\sk,-0.3);
 \foreach \x in {0,2,4,6,8,10,12}{
 \draw[semithick] (\x*\sk,-0.3) -- (\x*\sk,-0.43) node[below, font=\scriptsize] {\x};
 }
 \node[font=\scriptsize, anchor=north] at (6*\sk,-0.88) {czas [ms]};

 % legenda
 \fill[EEBlue!35, draw=black!50] (0,-2.15) rectangle (0.4,-1.88);
 \node[font=\scriptsize, anchor=west] at (0.5,-2.015) {wstępne przetwarzanie};
 \fill[EEUltramarine!45, draw=black!50] (4.0,-2.15) rectangle (4.4,-1.88);
 \node[font=\scriptsize, anchor=west] at (4.5,-2.015) {inferencja sieci};
 \fill[EETangerine!55, draw=black!50] (7.3,-2.15) rectangle (7.7,-1.88);
 \node[font=\scriptsize, anchor=west] at (7.8,-2.015) {NMS};
 \end{tikzpicture}
 \caption{Udział poszczególnych faz w średnim opóźnieniu potoku inferencji na akceleratorze NVIDIA GeForce RTX~4070. Źródło: opracowanie własne na podstawie \texttt{latency\_breakdown.json}.}
 \label{fig:latency_wykres}
\end{figure}

Całkowite opóźnienie ściany zegara (ang. \textit{end-to-end wall clock}) na stacji roboczej wyniosło średnio $10{,}94$\,ms ($p_{50}=10{,}98$\,ms, $p_{95}=11{,}15$\,ms), co odpowiada przepustowości rzędu $91{,}4$\,FPS. Dominującym składnikiem pozostaje inferencja sieci ($\approx 79\%$ średniego czasu e2e), natomiast preprocessing i NMS wnoszą łącznie ok.\ $2{,}1$\,ms w typowym przebiegu; proporcje te zilustrowano na rys.~\ref{fig:latency_wykres}. Wynik RTX~4070 jest górnym ograniczeniem wydajności potoku algorytmicznego. Hipotezę $H_1$ rozstrzyga jednak pomiar na telefonie, a nie na karcie graficznej PC.

Dlatego powtórzono profilowanie na urządzeniu produkcyjnym: Samsung Galaxy S23 Ultra (SM-S918U1, Snapdragon~8~Gen~2, Android~16). Mierzono wyłącznie wywołanie \texttt{MushroomDetector.detect()} zegarem monotonicznym \texttt{SystemClock.elapsedRealtimeNanos}. Wejście było identyczne z produkcją ($640\times640$ po letterboxingu CameraX). Pierwsze $8$ klatek odrzucono jako rozgrzewkę interpretera TFLite; zapisano $1\,899$ próbek przy wymuszonej kadencji $4$\,FPS (heartbeat $1{,}5$\,s wyłączono, aby SoC nie ostygł na statycznej scenie). Seria zimna to pierwsze $150$ klatek po rozgrzewce, w temperaturze pokojowej (bateria $31{,}2^{\circ}\mathrm{C}$). Seria gorąca to ostatnie $150$ klatek po ok.\ $8$--$9$\,min ciągłej inferencji (bateria $36{,}3^{\circ}\mathrm{C}$). Rozkład przedstawiono w tab.~\ref{tab:latency_s23}.

\begin{table}[htbp]
 \centering
 \caption{Czas wywołania \texttt{detect()} na Samsung Galaxy S23 Ultra. Wartości w milisekundach. Źródło: opracowanie własne na podstawie \texttt{inference\_profile.csv}.}
 \label{tab:latency_s23}
 \begin{tabular}{|l|c|c|c|}
 \hline
 \textbf{Wielkość} & \textbf{Zimny SoC} & \textbf{Po obciążeniu} & \textbf{Cała seria} \\
 \hline
  & ($n=150$) & ($n=150$) & ($n=1\,899$) \\
 \hline \hline
 Średnia & $156{,}5$ & $245{,}9$ & $216{,}4$ \\
 \hline
 $\sigma$ & $11{,}7$ & $36{,}9$ & $43{,}8$ \\
 \hline
 $p_{50}$ & $159{,}3$ & $262{,}5$ & $232{,}0$ \\
 \hline
 $p_{95}$ & $171{,}5$ & $278{,}3$ & $272{,}3$ \\
 \hline
 Minimum & $129{,}1$ & $145{,}3$ & $129{,}1$ \\
 \hline
 Maksimum & $177{,}2$ & $287{,}4$ & $406{,}4$ \\
 \hline
 Udział klatek $>250$\,ms & $0\%$ & $64\%$ & $27\%$ \\
 \hline
 Temperatura baterii & $31{,}2^{\circ}\mathrm{C}$ & $36{,}3^{\circ}\mathrm{C}$ & $31{,}2$--$36{,}3^{\circ}\mathrm{C}$ \\
 \hline
 \end{tabular}
\end{table}

Na zimnym SoC średnie $156{,}5$\,ms ($p_{95}=171{,}5$\,ms) mieści się w interwale throttlingu $250$\,ms: żadna z $150$ klatek nie przekroczyła tego progu. Po kilkuminutowym obciążeniu średnia rośnie do $245{,}9$\,ms, a $p_{95}$ do $278{,}3$\,ms; $64\%$ klatek gorącej serii przekracza $250$\,ms. Te liczby dotyczą wyłącznie \texttt{detect()}. Pobranie klatki CameraX i letterboxing $640\times640$ nie wchodzą w pomiar, więc z tab.~\ref{tab:latency_s23} nie wynika, że pełny cykl kamery zawsze zamyka się poniżej projektowych $300$\,ms. Przyrost o ok.\ $90$\,ms i wzrost temperatury baterii o $5{,}1$\,K to thermal throttling Snapdragona bez aktywnego chłodzenia (rozdz.~\ref{chap:przeglad_rozwiazan}). W produkcji heartbeat $1{,}5$\,s po zablokowaniu sceny ogranicza ten reżim: ciągłe $4$\,FPS przez ok.\ $8$--$9$\,min jest scenariuszem obciążeniowym, nie typową sesją zbieracza. Warstwa czasowa $H_1$ opiera się na zmierzonym \texttt{detect()} ($0{,}16$--$0{,}25$\,s średnio) i pracy offline; pełnego cyklu CameraX nie profilowano.

\section{Jakościowa analiza detekcji w warunkach rzeczywistych (dane zewnętrzne)}\label{sec:ood_field}
Metryki z poprzednich podrozdziałów pochodzą z izolowanego zbioru testowego iNaturalist. To ewaluacja wewnątrz rozkładu (ang. \textit{in-distribution}): to samo źródło, podobna optyka i ten sam protokół anotacji co w uczeniu. Aplikacja ma jednak działać w lesie, na strumieniu CameraX, przy oświetleniu podszytu i na matrycy konkretnego smartfonu. Taki reżim jest z definicji poza rozkładem uczącym (ang. \textit{out-of-distribution}, OOD). Dlatego uzupełniono ocenę ilościową jakościową inspekcją w terenie.

Testy przeprowadzono bezpośrednio na urządzeniu Samsung Galaxy S23 Ultra, opisanym w załączniku ze specyfikacją sprzętową. Autor skanował owocniki w naturalnym biotopie leśnym kamerą telefonu (lipiec--sierpień 2026), bez podawania modelowi zdjęć ze zbioru iNaturalist. Zebrano kilkadziesiąt zrzutów ekranu nakładki CameraX + Jetpack Compose; w załączniku~\ref{app:ood_screens} pozostawiono tylko interfejs, przykład detekcji wieloobiektowej, aktywację progu $0{,}18$ oraz dwa błędy modelu. Predykcje renderowała ta sama warstwa \texttt{MushroomRegistry.displayStatus}, która obowiązuje w produkcji. Ponieważ kadry terenowe nie mają ramek referencyjnych, nie liczono $mAP$. Oceniano zgodność gatunku, poprawność kategorii toksykologicznej oraz czytelność interfejsu.

Zaobserwowane scenariusze, uporządkowane od najprostszych do najtrudniejszych, zestawiono w tab.~\ref{tab:ood_scenariusze}.

\begin{table}[htbp]
 \centering
 \small
 \caption{Scenariusze jakościowej ewaluacji aplikacji na danych zewnętrznych (kamera smartfonu, biotop leśny). Źródło: opracowanie własne.}
 \label{tab:ood_scenariusze}
 \begin{tabular}{|p{2.55cm}|p{4.15cm}|p{3.55cm}|p{3.15cm}|}
 \hline
 \textbf{Scenariusz} & \textbf{Przykładowe gatunki} & \textbf{Status UI} & \textbf{Obserwacja} \\
 \hline \hline
 Pojedynczy owocnik jadalny & \textit{Pluteus cervinus}, \textit{Gyroporus castaneus}, \textit{Xerocomus chrysenteron} & zielona ramka, status Jadalny & poprawna asercja przy pewnym ujęciu \\
 \hline
 Wiele owocników śmiertelnych & \textit{Amanita phalloides} ($3+$ w kadrze) & czerwona ramka, Śmiertelnie trujący & jednoczesne wykrycie wszystkich widocznych owocników \\
 \hline
 Sceny mieszane i kępy & \textit{Amanita pantherina} / \textit{Amanita phalloides}, \textit{Hypholoma fasciculare} & pomarańczowa ramka, status Trujący & wieloklasowa detekcja w jednym kadrze \\
 \hline
 Gatunki nadrzewne & \textit{Tremella mesenterica}, \textit{Schizophyllum commune} & szara ramka, status Niejadalny & status neutralny, bez asercji jadalności \\
 \hline
 Tryb oszczędzania baterii & dowolny ustabilizowany kadr & baner trybu oszczędzania baterii & przejście na \textit{heartbeat} $1{,}5$\,s \\
 \hline
 \end{tabular}
\end{table}

W ujęciach pojedynczych gatunków jadalnych (\textit{Pluteus cervinus}, \textit{Gyroporus castaneus}, \textit{Xerocomus chrysenteron}, a także \textit{Russula cyanoxantha}) model lokalizował kapelusz zieloną ramką i przypisywał status „Jadalny” przy pewnym, zbliżonym kadrze. Nie obserwowano w tych warunkach fałszywej eskalacji do ostrzeżenia toksycznego. Przykład takiego ujęcia, zestawiony z kadrem o najwyższym priorytecie bezpieczeństwa, przedstawia rys.~\ref{fig:ood_jadalny_toksyczny}.

W scenach z wieloma owocnikami muchomora sromotnikowego (\textit{Amanita phalloides}) detektor oznaczał jednocześnie trzy lub więcej owocników w jednym kadrze. Każdy otrzymywał natychmiast czerwoną ramkę i etykietę „Śmiertelnie trujący”, zgodnie z progiem $\mathrm{conf}_{\mathrm{toxic}}=0{,}18$. W kadrze mieszanym obok muchomora sromotnikowego pojawiał się muchomor plamisty (\textit{Amanita pantherina}) -- system rozróżniał obie klasy i oznaczał ramką pomarańczową „Trujący” dla plamistego, nie degradując statusu sromotnikowego. Gęste kępy maślanki wiązkowej (\textit{Hypholoma fasciculare}) oraz kadry wielogatunkowe z \textit{Mycena galericulata} potwierdziły, że wieloobiektowa głowica YOLO11m nie sprowadza całej sceny do jednej etykiety: sąsiadujące kępy otrzymują osobne ramki.

Gatunki nadrzewne -- trzęsak pomarańczowożółty (\textit{Tremella mesenterica}) na korze oraz rozszczepka pospolita (\textit{Schizophyllum commune}) na leżaninie -- otrzymywały szarą ramkę „Niejadalny”. Warstwa prezentacji nie eskalowała ich ani do zielonej asercji jadalności, ani do fałszywego alarmu toksycznego. Po ustabilizowaniu kadru nakładka wyświetlała baner „Zablokowano -- tryb oszczędzania baterii”: \texttt{isSceneLocked()} podnosiła interwał inferencji z $250$\,ms do \textit{heartbeat} $1{,}5$\,s, zgodnie z implementacją z rozdz.~\ref{chap:implementacja}. Oba zjawiska -- detekcję nadrzewną i blokadę energetyczną -- widać na rys.~\ref{fig:ood_battery_nadrzewne}.

\begin{figure}[H]
 \centering
 \begin{subfigure}[t]{0.48\textwidth}
  \centering
  \includegraphics[width=\linewidth,height=11.2cm,keepaspectratio]{screenshots/ood_01_pluteus_cervinus.jpg}
  \caption{\textit{Pluteus cervinus}, jadalny, zielona ramka.}
  \label{fig:ood_pluteus}
 \end{subfigure}
 \hfill
 \begin{subfigure}[t]{0.48\textwidth}
  \centering
  \includegraphics[width=\linewidth,height=11.2cm,keepaspectratio]{screenshots/ood_05_amanita_mixed.jpg}
  \caption{\textit{Amanita phalloides} (czerwone) i \textit{Amanita pantherina} (pomarańczowa) w jednym kadrze.}
  \label{fig:ood_amanita_mix}
 \end{subfigure}
 \caption{Detekcja w warunkach terenowych: pojedynczy owocnik jadalny oraz równoczesne oznaczenie gatunków śmiertelnych i trujących. Źródło: opracowanie własne (zrzuty aplikacji, CameraX + Jetpack Compose).}
 \label{fig:ood_jadalny_toksyczny}
\end{figure}

\begin{figure}[H]
 \centering
 \begin{subfigure}[t]{0.48\textwidth}
  \centering
  \includegraphics[width=\linewidth,height=11.2cm,keepaspectratio]{screenshots/ood_09_tremella_battery.jpg}
  \caption{\textit{Tremella mesenterica}, niejadalny; baner trybu oszczędzania baterii.}
  \label{fig:ood_tremella}
 \end{subfigure}
 \hfill
 \begin{subfigure}[t]{0.48\textwidth}
  \centering
  \includegraphics[width=\linewidth,height=11.2cm,keepaspectratio]{screenshots/ood_07_hypholoma_fasciculare.jpg}
  \caption{Kępy \textit{Hypholoma fasciculare}, trujący, pomarańczowe ramki.}
  \label{fig:ood_hypholoma}
 \end{subfigure}
 \caption{Praca aplikacji poza zbiorem iNaturalist: gatunek nadrzewny z blokadą energetyczną oraz wieloobiektowa detekcja kęp trujących. Źródło: opracowanie własne.}
 \label{fig:ood_battery_nadrzewne}
\end{figure}

Ograniczenia tej ewaluacji są jasne i nie należy ich ukrywać. Próba jest okazjonalna, sezonowa (lipiec--sierpień) i nie pokrywa wszystkich $147$ klas. Częściowe przesłonięcie ściółką, źdźbłami trawy lub gałązkami obniża pewność i bywa przyczyną pominięcia drugiego, mniejszego owocnika w kadrze; gęste kępy generują nachodzące etykiety nakładki. Zmienne oświetlenie podszytu nie uniemożliwiało jednak poprawnej klasyfikacji kategorii toksykologicznej w zebranych ujęciach poprawnych. Fałszywy pozytyw na kamieniu oraz pominięcie mniejszego skupiska \textit{Tremella} udokumentowano w załączniku~\ref{app:ood_screens}. W zebranym materiale terenowym nie zaobserwowano błędu typu „śmiertelnie trujący prezentowany jako jadalny”: czerwone i pomarańczowe ramki pojawiały się już od progu ostrzeżenia $0{,}18$, a zielony status wymagał pewności $\ge 0{,}60$. To zgodność z logiką \texttt{displayStatus} na niewielkiej próbie OOD, a nie pomiar CFPR. $H_2$ weryfikuje tab.~\ref{tab:cfpr_porownanie} na izolowanym teście iNaturalist. Bufor konsensusu $4$ z $5$ klatek, opisany w rozdz.~\ref{chap:implementacja}, działał w tle jako mechanizm stabilizacji nakładki; nie był przedmiotem osobnego porównania z wariantem bez bufora.

\chapter{Podsumowanie i wnioski}\label{chap:podsumowanie}
Zaprojektowano i wdrożono autonomiczny system na platformę Android. Aplikacja detekuje i klasyfikuje $147$ gatunków grzybów offline, przy kadencji do $4$\,FPS na S23 Ultra. Potok obejmował akwizycję danych, auto-anotację GroundingDINO, trening YOLO11m, konwersję TFLite FP16 oraz aplikację z asymetryczną warstwą progów, która priorytetyzuje czułość dla klas toksycznych.

\section{Wnioski z realizacji projektu}
Cel główny – autonomiczny, offline system identyfikacji grzybów na Androidzie – został osiągnięty. Rozwiązanie łączy detekcję wieloobiektową z nakładką kamery i trwałym zapisem historii znalezisk, spełniając założenia z rozdz.~\ref{chap:projekt_systemu}.

Potok półautomatycznej anotacji oparty na GroundingDINO zastąpił masowe ręczne znakowanie ramek auto-anotacją zero-shot, przy zachowaniu formatu YOLO ($\texttt{class\_id cx cy w h}$) dla ponad stu gatunków \cite{liu2023grounding}. Ręczny audyt $1\,065$ auto-anotowanych obrazów potwierdził poprawność lokalizacji dla $94{,}93\%$ istniejących ramek ($1\,011$ z $1\,065$). Na stratyfikowanej próbie $250$ kadrów izolowanego testu nauczyciel pominął $34$ z $441$ widocznych owocników ($7{,}71\%$; tab.~\ref{tab:audyt_fn}). Bez tego mechanizmu zbudowanie zbioru uczącego o takiej skali byłoby poza zasięgiem pracy inżynierskiej.

Model YOLO11m osiągnął na izolowanym zbiorze testowym wartości $mAP_{50}=82{,}75\%$ oraz $mAP_{50\text{--}95}=78{,}23\%$ (PyTorch, \texttt{split=test}, zero wspólnych \texttt{obs\_id} z korpusem uczącym). Na tym samym teście wariant Nano uzyskał $48{,}54\%$ / $42{,}82\%$, co potwierdza wybór Medium. Na próbie $2\,500$ kadrów (\texttt{seed}=42) artefakt TFLite FP16 z dekoderem Androida dał $mAP_{50}=0{,}7914$ wobec $0{,}8467$ dla \texttt{.pt} (różnica $-0{,}0553$; podrozdz.~\ref{sec:tflite_parity}). Na stacji roboczej (RTX~4070) potok algorytmiczny zamyka się w $10{,}94$\,ms; na S23 Ultra średni czas \texttt{detect()} wynosi $156{,}5$\,ms na zimnym SoC i $245{,}9$\,ms po kilkuminutowym obciążeniu. Pełnego cyklu CameraX nie profilowano, więc nie twierdzi się, że cały przebieg klatki mieści się poniżej $300$\,ms. Wyniki te potwierdzają praktyczną wykonalność podejścia \textit{edge AI} w zadaniu identyfikacji mykologicznej, bez zależności od zasobów chmurowych \cite{ultralytics2024yolo11}.

Polityka \texttt{displayStatus} nie prezentuje zielonej etykiety „Jadalny” przy pewności poniżej $0{,}60$ ani dla predykcji niejadalnych. Na izolowanym teście CFPR pary $0{,}18$/$0{,}60$ wynosi $1{,}50\%$ ($327/15\,840$), wobec $3{,}11\%$ ($580/15\,840$) przy $\tau=0{,}30$ -- $H_2$ spełniona na tym pomiarze. Jakościowa inspekcja w terenie (podrozdz.~\ref{sec:ood_field}) pokazała, że zielone ramki jadalne, czerwone i pomarańczowe ostrzeżenia oraz szary status niejadalny działają na strumieniu CameraX poza zbiorem iNaturalist. Nie zastępuje to oceny mykologa: $1{,}50\%$ obiektów toksycznych nadal może otrzymać asercję jadalności.

\section{Perspektywy rozwoju i dalsze prace}
Kolejnym krokiem może być integracja algorytmów śledzenia obiektów wideo (np. ByteTrack lub SORT) z potokiem CameraX. Stabilizacja ramek w czasie zmniejszyłaby chwilowe wahania predykcji pomiędzy klatkami i poprawiła płynność nakładki interfejsu użytkownika.

Drugim istotnym obszarem pozostaje optymalizacja rozmiaru modelu. Warto zbadać uczenie z uwzględnieniem kwantyzacji (ang. \textit{Quantization-Aware Training}, QAT) oraz skalibrowaną precyzję mieszaną (FP16/INT8), tak aby zredukować artefakt poniżej $20$\,MB przy zachowaniu wrażliwości na subtelne cechy makromorfologiczne (blaszki, pierścień), kluczowe dla rozróżniania gatunków bliźniaczych.

Trzecią perspektywą jest fuzja osadzeń wizyjnych z metadanymi kontekstowymi: klasyfikacją biotopu na podstawie GPS, typem lasu, sezonowością oraz lokalnymi warunkami pogodowymi. Takie wzbogacenie sygnału wejściowego mogłoby poprawić rozróżnianie gatunków o zbliżonej morfologii, lecz odmiennych wymaganiach siedliskowych.

\section{Repozytorium projektu}
Kod źródłowy aplikacji mobilnej w języku Kotlin, skrypty potoku przetwarzania danych w języku Python oraz narzędzia eksportu modeli zostały udostępnione w publicznym repozytorium pod adresem: \url{https://github.com/kirilhorobets/mushroom-recognition}.

W przygotowaniu niniejszej pracy wykorzystano narzędzia oparte na metodach sztucznej inteligencji jako wsparcie redakcji oraz korekty językowej i stylistycznej tekstu; zaproponowane zmiany były każdorazowo weryfikowane i akceptowane przez autora.
Narzędzia oparte na metodach sztucznej inteligencji wykorzystano również pomocniczo przy opracowaniu i korekcie kodu źródłowego; wszystkie sugestie były każdorazowo weryfikowane i akceptowane przez autora.
